%%%%%%%%%%%%%%%%%%%%%%6-variant 2017/06/2 %%%%%%%%%%%%%%%%%%%%%%%%%%%%%%%%%%%%
\chapter{Matrix representations of GA and related groups}

Description of GA as a matrix algebra may be used as a bridge that
connects equations in matrix or tensor form with their
counterparts in GA. It turns out that representations of real
algebras \cl{p}{q} can be built up as tensor products of matrices
over three basic fields, \bbR, \bbC\ and \bbH.  
They are summarized in a form of irreducible $8$-periodicity table.

The GA contains a number of important groups such as versor
$\Gamma_{p,q}$, pinor \Pin{p}{q}, and spinor \Spin{p}{q} groups.
The elements of the mentioned groups transform one MV to other
belonging to the same class, i.~e. the group elements describe
symmetry properties of transformations between various sets of
MV's in a given GA. Such transformations are often related  with
the reflections and rotations in $n$-dimensional spaces having
different metrics. When GA groups are represented in matrix form
one frequently recognizes them as the well-known classical
symmetry groups, however, in the GA form their elements frequently
acquire a more clear geometric interpretation.

\section{Matrix representation of \cl{p}{q} and $8$-periodicity}
In matrix representation, the GA basis elements are written as
square matrices and geometric product of elements is replaced by matrix
product. There exists many sets of irreducible matrices that
represent the same GA. The representation matrices might be related by some
similarity transformation. The representations are either simple,
hence isomorphic to matrix algebras over the real, complex or
quaternionic numbers, or semisimple, i.~e. isomorphic to the
direct sum of two real or quaternionic matrices.

%\noindent[Rausch~\cite{Rausch05}, [Venselaar07], [Degirmenci
%078613.pdf] Ablamowicz~\cite{Ablamowicz12},
%Porteous~\cite{Porteous04}]
%
\begin{unit}\textit{Matrix direct sum and direct product}.
If $\hat{A}$ and $\hat{B}$ is, respectively, $m\times n$ and
$p\times q$ matrix then their direct sum is $(m+p)\times (n+q)$
matrix
\[\hat{A}\oplus\hat{B}=
\begin{bmatrix}
\hat{A}&\bzero\\
\bzero&\hat{B}
\end{bmatrix}\,.\]
The direct product $\hat{C}=\hat{A}\otimes\hat{B}$ of matrices
$\hat{A}$ and $\hat{B}$ is $(mn)\times (pq)$ matrix $\hat{C}$ with
elements defined by $c_{\alpha\beta}=a_{ij}b_{kl}$, where
$\alpha\equiv p(i-1)+k$ and $\beta\equiv q(j-1)+l$.
\end{unit}
%
\begin{example}
The direct product of $2\times 2$ matrix
$\hat{A}=\left[\begin{smallmatrix}a_{11}&a_{12}\\a_{21}&a_{22}\end{smallmatrix}\right]$
and $m\times n$ matrix $\hat{B}=\left[\begin{smallmatrix}b_{11}&b_{12}&\cdots&b_{1n}\\
b_{21}&a_{22}&\cdots&b_{2n}\\
\cdots&\cdots&\cdots&\cdots\\
b_{m1}&b_{m2}&\cdots&b_{mn}
\end{smallmatrix}\right]$ is given by the following $2m\times 2n$ matrix
\[\hat{A}\otimes\hat{B}=
\begin{bmatrix}
a_{11}\hat{B}&a_{12}\hat{B}\\
a_{21}\hat{B}&a_{22}\hat{B}
\end{bmatrix}, \quad\textrm{where} \quad
\quad a_{ij}\hat{B}=\begin{bmatrix}
  a_{ij}b_{11}&a_{ij}b_{12}&\cdots&a_{ij}b_{1n}\\
a_{ij}b_{21}&a_{ij}b_{22}&\cdots&a_{ij}b_{2n}\\
  \cdots&\cdots&\cdots&\cdots\\
a_{ij}b_{m1}&a_{ij}b_{m2}&\cdots&a_{ij}b_{mn}\\
\end{bmatrix}\]
\end{example}
%
\begin{unit}\textit{Isomorphisms of matrix algebras}.\label{isomatalg}
Real, complex and quaternionic numbers, $\bbR$, $\bbC$ and $\bbH$,
make up the fields with two binary operations, i.~e. addition and
multiplication. The addition operation is commutative, while the
multiplication is commutative only for real and complex numbers.
Real, complex and quaternionic $m\times m$ matrices will be
denoted, respectively, by $\bbR(m)$, $\bbC(m)$ and $\bbH(m)$ with
$\bbR(1)\equiv\bbR$, etc. Complex  numbers are isomorphic
to $2\times 2$ real matrices, $\bbC\cong\bbR(2)$, %\cite{Porteous95}p.~11
and quaternions are isomorphic to complex as well
as real matrices, $\bbH\cong\bbC(2)\cong\bbR(4)$.%\cite{Porteous95} %.~69


The below listed  direct product isomorphisms of matrices and
Clifford numbers play an essential role in the construction of the
$8$-periodicity table in GA~\cite{Rausch05,Porteous95}, %p.4,p.84
\begin{align}
\bbR(m)\otimes_{\bbR}\bbR&\cong\bbR(m)
&\bbC\otimes_{\bbR}\bbC&\cong\bbC\oplus\bbC\nonumber\\
\bbR(m)\otimes_{\bbR}\bbR(n)
&\cong\bbR(mn)&\bbC\otimes_{\bbR}\bbH&\cong\bbC(2)\nonumber\\
\bbR(m)\otimes_{\bbR}\bbC&\cong\bbC(m)
&\bbH\otimes_{\bbR}\bbH&\cong\bbR(4)\nonumber\\
\bbR(m)\otimes_{\bbR}\bbH&\cong\bbH(m)&
&\nonumber
\end{align}
Here $\otimes_{\bbR}$ indicates that direct product of algebras is
computed over the real field. This means that GA base elements, or
respective matrices are allowed to be multiplied by real numbers only. The
field index ${}_{\bbR}$ frequently  is omitted.
\end{unit}
%
\begin{example}
To get the algebraic isomorphism
$\bbC\otimes\bbC\cong\bbC\oplus\bbC\equiv {^2\bbC}$ it is enough
to remember that ${^2\bbC}$ is generated as a real algebra by
subalgebras
$\left\{\left[\begin{smallmatrix}z&0\\0&z\end{smallmatrix}\right]:z\in\bbC\right\}$
and
$\left\{\left[\begin{smallmatrix}z&0\\0&z^*\end{smallmatrix}\right]:z\in\bbC\right\}$.

The isomorphism $\bbC\otimes\bbH\cong\bbC(2)$ is easy to
understand if one remembers that algebra over \bbR\ can be
complexified by making its underlying vector space complex valued. Indeed,
recall that the quaternion can be represented by Pauli matrices as
$\bbH\cong
q_0\sig{0}+q_1(-\ii\sig{1})+q_2(-\ii\sig{2})+q_3(-\ii\sig{3})$,
where $q_i\in\bbR$ and $\sig{0}$ is the unit matrix, so $\bbC(2)$
can be regarded as a complexification of the real algebra \bbH.
\end{example}
%
\begin{unit}\label{simplereps}\textit{The simplest real GA's and their representation by irreducible (fundamental) matrices} are
\[\begin{array}{lll}
 \cl{0}{0}\cong\bbR\qquad&\cl{1}{0}\cong\bbR\oplus\bbR\qquad&\cl{0}{1}\cong\bbC\\
 \cl{2}{0}\cong\bbR(2)\qquad& \cl{1}{1}\cong\bbR(2)\qquad&
 \cl{0}{2}\cong\bbH\cong\bbC(2)
  \end{array}\]
where $\bbR(n)$ and $\bbC(n)$ are, respectively, real and complex
$n\times n$ matrices. The explicit form of the above matrices
determines representations for all other \cl{p}{q} algebras. The
representations of \cl{p}{q} and \cl{p^\prime}{q^\prime} algebras
has the same structure (see unit~\ref{8foldperiod}) when
$(p-q)\mod(8)=(p^\prime-q^\prime)\mod(8)$.
\end{unit}
%
\begin{unit}\textit{Explicit matrix representation of the simplest real geometric algebras\label{GAreps}}.
Below are listed the simplest matrix representations of real GA's
(item~\ref{simplereps}) from which explicit representations of
higher algebras \cl{p}{q} later are constructed.
\begin{itemize}
\item[\cl{0}{0}] describes real number  algebra which is
represented by $1\times 1$ real matrices, i.~e. real numbers
$\bbR$.

\item[\cl{1}{0}] is the parabolic number algebra. It has a single
basis vector $\e{1}$ the square of which is
$\e{1}^2=\e{1}\e{1}=1$. There are $2^1=2$ basis elements that are
represented by $\bbR\oplus\bbR$ two real matrices,
\[\cl{1}{0}\cong\{1,\e{1}\}\cong\left\{
\begin{bmatrix}1&0\\0&1\end{bmatrix},
\begin{bmatrix}1&0\\0&-1\end{bmatrix}
\right\}.\]

\item[\cl{0}{1}] is the complex number algebra. It has one basis
vector which may be represented either by imaginary unit
$\e{1}\cong\ii=\sqrt{-1}$, or by real $\bbR(2)$ matrix. The matrix
representations of  $2^1=2$ basis elements are:
\[\cl{0}{1}\cong\{1,\e{1}\}\cong\{1,\ii\}\cong\left\{
\begin{bmatrix}1&0\\0&1\end{bmatrix},
\begin{bmatrix}0&1\\-1&0\end{bmatrix}
\right\}.\]
 \item[\cl{2}{0}] algebra describes the Euclidean plane and is called the coquaternion
algebra. It has two basis vectors $\{\e{1},\e{2}\}$ and $2^2=4$
basis elements: $\{1,\e{1},\e{2},\e{12}\}$. In matrix form, the
orthonormal basis elements of this algebra typically are
represented as
\[\cl{2}{0}\cong\{1,\e{1},\e{2},\e{12}\}\cong\left\{
\begin{bmatrix}1&0\\0&1\end{bmatrix},
\begin{bmatrix}0&1\\1&0\end{bmatrix},
\begin{bmatrix}1&0\\0&-1\end{bmatrix},
\begin{bmatrix}0&-1\\1&0\end{bmatrix}
\right\}.\] The role of orthonormal base vectors here plays two
Pauli matrices $\e{1}\cong\hat{\sigma}_1$ and
$\e{2}\cong\hat{\sigma}_3$. More generally, the base vectors of
\cl{2}{0} can be parameterized by angle $\varphi$ as  rotation
matrices,
\begin{align*}
\{\e{1},\e{2}\}&\cong\left\{
 \begin{bmatrix}-\cos\varphi&\sin\varphi\\ \sin\varphi&\cos\varphi\end{bmatrix},
\begin{bmatrix}\sin\varphi&\cos\varphi\\ \cos\varphi& -\sin\varphi\end{bmatrix}
\right\}\\
&\cong\left\{
 \begin{bmatrix}-\sin\varphi&\ii \cos\varphi\\ -\ii \cos\varphi&\sin\varphi\end{bmatrix},
\begin{bmatrix}-\cos\varphi&-\ii \sin\varphi\\ \ii \sin\varphi& \cos\varphi\end{bmatrix}
 \right\}.
\end{align*}
%
\item[\cl{1}{1}] algebra describes the Minkowski plane. Two basis
vectors $\{\e{1},\e{2}\}$ with signatures $\e{1}^2=1$ and
$\e{2}^2=-1$, and the bivector $\e{12}$ in matrix form are be
represented by $\bbR(2)$ matrices
\[\cl{1}{1}\cong\{1,\e{1},\e{2},\e{12}\}\cong\left\{
\begin{bmatrix}1&0\\0&1\end{bmatrix},
\begin{bmatrix}0&1\\1&0\end{bmatrix},
\begin{bmatrix}0&1\\-1&0\end{bmatrix},
\begin{bmatrix}-1&0\\0&1\end{bmatrix}
\right\}.\] A more general parametrization of \cl{1}{1}
orthonormal base vectors may be taken as
\begin{align*}
\{\e{1},\e{2}\}&\cong\left\{
\begin{bmatrix}\sin\varphi&\cos\varphi\\ \cos\varphi& -\sin\varphi\end{bmatrix},
 \begin{bmatrix}-\ii\cos\varphi&\ii\sin\varphi\\ \ii \sin\varphi& \ii \cos\varphi\end{bmatrix}
\right\}\\
%&\cong\left\{
%\begin{bmatrix}-\cos\varphi&\sin\varphi\\ \sin\varphi& \cos\varphi\end{bmatrix},
%\begin{bmatrix}-\ii \sin\varphi&-\ii \cos\varphi\\ -\ii \cos\varphi&\ii \sin\varphi\end{bmatrix}
%\right\}\\
%&\cong\left\{
%\begin{bmatrix}-\ii\sinh\varphi&-\ii\cosh\varphi\\ \ii \cosh\varphi& \ii \sinh\varphi\end{bmatrix},
%\begin{bmatrix}\ii \cosh\varphi&\ii \sinh\varphi\\ -\ii \sinh\varphi& -\ii \cosh\varphi\end{bmatrix}
%\right\}\\
&\cong\left\{
\begin{bmatrix}\cosh\varphi&\sinh\varphi\\ -\sinh\varphi& -\cosh\varphi\end{bmatrix},
\begin{bmatrix}-\sinh\varphi&-\cosh\varphi\\ \cosh\varphi& \sinh\varphi\end{bmatrix}
\right\}.
\end{align*}

\item[\cl{0}{2}] is the quaternion algebra. Its basis vectors
$\{\e{1},\e{2}\}$ have negative signature. Orthonormal base
elements of the algebra can be represented by $4\times 4$ real
matrices
  \[
\begin{gathered}
\cl{0}{2}\cong\bbH\cong\{1,\e{1},\e{2},\e{12}\}\cong\\
\Big\{
\renewcommand*{\arraystretch}{0.9}
\arraycolsep=0.9\arraycolsep
\scriptstyle
\begin{bmatrix}1&0&0&0\\0&1&0&0\\0&0&1&0\\0&0&0&1\end{bmatrix},
\begin{bmatrix}0&-1&0&0\\1&0&0&0\\0&0&0&-1\\0&0&1&0\end{bmatrix},
\begin{bmatrix}0&0&-1&0\\0&0&0&1\\1&0&0&0\\0&-1&0&0
\end{bmatrix},\begin{bmatrix}0&0&0&-1\\0&0&-1&0\\0&1&0&0\\1&0&0&0\end{bmatrix}
\Big\}.
\end{gathered}
\]
The quaternion also can be represented by $2\times 2$ Pauli matrices
$\hat{\sigma}_{i}$ as well,
\[\bbH\cong
q_0\hat{\sigma}_{0}+q_1(-\ii\hat{\sigma}_{1})+q_2(-\ii\hat{\sigma}_{2})+q_3(-\ii\hat{\sigma}_{3})\]
where $\hat{\sigma}_{0}$ is the unit matrix and $q_i\in\bbR$. More
generally, the two base vectors may be parameterized  by complex
matrices
\begin{align*}
\{\e{1},\e{2}\}
 &\cong\left\{
\begin{bmatrix}-\ii\sin\varphi&-\cos\varphi\\ \cos\varphi& \ii\sin\varphi\end{bmatrix},
\begin{bmatrix}-\ii \cos\varphi&\sin\varphi\\ -\sin\varphi&\ii \cos\varphi\end{bmatrix}
\right\}\\
&\cong\left\{
\begin{bmatrix}-\sinh\varphi&-\ii\cosh\varphi\\ -\ii \cosh\varphi& \sinh\varphi\end{bmatrix},
\begin{bmatrix}-\ii \cosh\varphi&\sinh\varphi\\ \sinh\varphi& \ii \cosh\varphi\end{bmatrix}
\right\}.
\end{align*}
\end{itemize}
\end{unit}
%
\begin{unit}
\textit{Generation of higher order  algebras by a direct product
of the fundamental GA's}~\cite{Rausch05,Song2013}\label{isomorphismGeneration}, %[Venselaar07]
\begin{alignat*}{3}
  &\cl{0}{0}\otimes\cl{1}{1\,}&\cong\cl{1}{1\hphantom{+2}}&&\qquad\cl{p}{q}\otimes\cl{1}{1}&\cong\cl{p+1}{q+1}\\
&\cl{0}{q}\otimes\cl{2}{0}&\cong\cl{q+2}{0}&&\cl{p}{q}\otimes\cl{2}{0}&\cong\cl{q+2}{p}
\\
&\cl{p}{0}\otimes\cl{0}{2}&\cong\cl{0}{p+2}&&\cl{p}{q}\otimes\cl{0}{2}&\cong\cl{q}{p+2}
\\
& & && \cl{p}{q}\otimes\cl{0}{4}&\cong\cl{p}{q+4}\cong\cl{p+4}{q}
%Ablamowicx etal Lectures on Clifford Alg p.35
\end{alignat*}
Note that after multiplication by $\cl{2}{0}$ or $\cl{0}{2}$, the
indices $p$ and $q$ in the right-hand side are interchanged.
\end{unit}
%
\begin{unit}
\textit{Isomorphisms between different GA's and their
products}~\cite{Rausch05,Lounesto97,Porteous95},
\label{periodicityTableExtension}
\begin{alignat*}{3}
&\cl{p}{q+1}&\cong\cl{q}{p+1}\qquad\quad &&\cl{p}{q}&\cong\cl{q+1}{p-1},\quad p\ge1\\
&&&&\cl{p}{q}&\cong\cl{p-4}{q+4},\quad p\ge4
\end{alignat*}
\[\begin{split}
&\cl{p}{q+4}\cong\cl{p}{q}\otimes\cl{0}{4}\cong\cl{p}{q}\otimes\bbH(2)\\
&\cl{8}{0}\cong\cl{0}{8}\cong\bbR(16)\\
&\cl{p}{q+8}\cong\cl{p}{q}\otimes\cl{0}{8}\cong\cl{p}{q}\otimes\bbH(2)\otimes\bbH(2)\cong\cl{p}{q}\otimes\bbR(16)\\
&\cl{p+8}{q}\cong\cl{p}{q}\otimes\cl{0}{8}\cong\cl{p}{q}\otimes\bbH(2)\otimes\bbH(2)\cong\cl{p}{q}\otimes\bbR(16)\\
&\cl{p+k}{q+k}\cong\cl{p}{q}\otimes\underbrace{\cl{1}{1}\otimes\dots\otimes\cl{1}{1}}_{\textrm{k
factors}}
\cong\cl{p}{q}\otimes\bbR(2k)\\
\end{split}\]
%
\[
\begin{split}
&\mathit{Cl}^{+}_{p,q}\cong \mathit{Cl}^{+}_{q,p}\qquad\textrm{isomorphism between even subalgebras}\\
&\cl{p}{q}\cong \cl{p}{q-1}\oplus \cl{p}{q-1}\quad \textrm{if $p-q=1\mod(8)$ and $p+q$ is odd}\\
&\cl{p}{q}\cong \cl{q}{p-1}\oplus \cl{q}{p-1}\quad \textrm{if $p-q=1\mod(8)$ and $p+q$ is odd}
\end{split}
\]
These isomorphisms imply that the Clifford algebras have a period
of~8.
\end{unit}
%
\begin{example}
Using the direct products for fundamental GA's and their matrix
representations (see item~\ref{isomatalg}) the following
isomorphisms between higher GA's and their matrix representations
can be established.
\begin{equation*}\begin{split}
&\cl{3}{0}\cong\cl{0}{1}\otimes\cl{2}{0}\cong\bbC\otimes\bbR(2)\cong\bbC(2)\\
&\cl{1}{3}\cong\cl{0}{2}\otimes\cl{1}{1}\cong\bbH\otimes\bbR(2)\cong\bbH(2)\\
&\cl{3}{1}\cong\cl{2}{0}\otimes\cl{1}{1}\cong\bbR(2)\otimes\bbR(2)\cong\bbR(4)\\
&\cl{4}{1}\cong\cl{0}{1}\otimes\cl{2}{0}\otimes\cl{1}{1}=\bbC\otimes\bbR(2)\otimes\bbR(2)
\cong\bbC\otimes\bbR(4)\cong\bbC(4)\\
&\cl{1}{4}\cong\cl{1}{0}\otimes\cl{0}{4}\cong{^2\bbR}\otimes\bbH(2)\cong{^2\bbH(2)}\equiv\bbH(2)\oplus\bbH(2)\,.
\end{split}\end{equation*}
\end{example}
%
\begin{example} The irreducible matrix representation \cl{0}{7}$\cong
{^2\bbR(8)}$ can be obtained in the following way:
\[\begin{split}
\cl{0}{7}&\cong\cl{5}{0}\otimes\cl{0}{2}
\cong\cl{0}{3}\otimes\cl{2}{0}\otimes\cl{0}{2}
\cong\cl{1}{0}\otimes\cl{0}{2}\otimes\cl{2}{0}\otimes\cl{0}{2}\\
&\cong(\bbR\oplus\bbR)\otimes\bbH\otimes\bbR(2)\otimes\bbH
\cong(\bbH\oplus\bbH)\otimes\bbR(2)\otimes\bbH\\
&\cong(\bbR(4)\oplus\bbR(4))\otimes\bbR(2)
\cong(\bbR(8)\oplus\bbR(8))\equiv{^2\bbR(8)\,.}
\end{split}\]
Similarly for  \cl{0}{8} we derive \cl{0}{8}$\cong
{\bbR(16)}$,
\[\begin{split}
&\cl{0}{8}\cong\cl{6}{0}\otimes\cl{0}{2}\cong
\cl{0}{4}\otimes\cl{2}{0}\otimes\cl{0}{2}\cong
\cl{2}{0}\otimes\cl{0}{2}\otimes\cl{2}{0}\otimes\cl{0}{2}\cong\\
&\bbH\otimes\bbH\otimes\bbR(2)\otimes\bbR(2)\cong\bbR(4)\otimes\bbR(4)
\cong\bbR(16)\,.
\end{split}\]
\end{example}
%
\begin{example}
Using the above isomorphisms the irreducible matrix representation
for \cl{2}{8} algebra is found as
\[\cl{2}{8}\cong\cl{2}{0}\otimes\cl{0}{8}\cong\bbR(2)\otimes\bbR(16)\cong\bbR(32) \,.\]
Thus, \cl{2}{8} algebra elements can be represented by real
$32\times 32$ matrices and the algebra consists of  $2^{2+8}=1024$
such matrices.
\end{example}
%
\begin{unit}\textit{Construction of matrix representations for higher GA's}.\label{constructionOfRepresentation}
A number of different schemes to construct Clifford algebra
representations exist \cite{HILE199051,Poole1982,Ablamowicz2014}.

The matrix representation of arbitrary  GA can be obtained
directly from the isomorphisms of unit~\ref{isomorphismGeneration}.
Sequential application of these isomorphisms allows to get matrix base vectors for higher
dimensional GA's \cl{q+2}{p}, \cl{q}{p+2} or \cl{p+1}{q+1} once
the vectors of \cl{p}{q} and, respectively, bivectors of
\cl{2}{0}, \cl{0}{2} and \cl{1}{1} are known. Indeed, let us
assume that matrix representations of base vectors of \cl{p}{q}
are already known. The matrix representations of low dimensional
algebras, for example, of \cl{2}{0}, etc. can be taken from unit
\ref{GAreps}. The isomorphism
$\cl{p}{q}\otimes\cl{2}{0}\cong\cl{q+2}{p}$  can be constructively
proved simply by checking that $p+q+2$ orthonormal base vectors of
\cl{q+2}{p} are
\begin{equation*}
\{\e{1}\otimes\f{12},\dotsc,\e{p}\otimes\f{12},\e{p+1}\otimes\f{12},\dotsc,\e{p+q}\otimes\f{12},
1\otimes\f{1},1\otimes\f{2}\}\,,
\end{equation*}
where $\f{12}=\f{1}\f{2}$ is a bivector of $\cl{2}{0}$. Indeed let
us define the \cl{q+2}{p} algebra
\begin{equation*}
\e{i}^\prime\e{j}^\prime+\e{j}^\prime\e{i}^\prime=2\eta_{ij},\quad \textrm{where}\quad
\eta_{ii} =
\begin{cases}
+1 &\mbox{when}\quad 1\le i \le q+2, \\
-1 &\mbox{when}\quad q+2<i\le p+q+2,
\end{cases}
\end{equation*}
and $\eta_{ij}=0$, if $i\neq j$. Calculating the squares of
vectors $\e{i}^\prime=\e{p}\otimes\f{12}$, with $\e{p}^2=+1$ and
$1\le i\le q+2$, we see that the sign was reversed
\begin{equation*}
\begin{split}
(\e{p}\otimes\f{12})(\e{p}\otimes\f{12})+(\e{p}\otimes\f{12})(\e{p}\otimes\f{12})&=2(\e{p}^2\otimes\f{12}^2)\\
&=2(1\otimes(-1))=-2.
\end{split}
\end{equation*}
It means that signatures of $p+q$ vectors $\e{i}\otimes\f{12}$
become oposite, i.~e. $p$ vector turns into $q$ vector and vice
versa. The signature of the rest two base vectors $1\otimes\f{1}$
and $1\otimes\f{2}$ remains the same as that of  $\f{1}$ and
$\f{2}$, respectively. It follows, that the above constructed base
defines \cl{q+2}{p} algebra. The two other isomorphisms are
checked in exactly the same way. Substitution of matrix
representations for corresponding vectors and bivectors into the
expression of base of \cl{q+2}{p}, therefore, yields explicit
matrix representation of the algebra.
\end{unit}
%
\begin{example}\label{excl30}
Let us find matrix representation of \cl{3}{0} in explicit form.
Multiplication of lower order algebras from unit~\ref{GAreps}
gives
\[\begin{split}
&\cl{3}{0}\cong\cl{0}{1}\otimes\cl{2}{0}
\xRightarrow{\textrm{elements}} \{1,\e{1}\}\otimes\{1,\e{1},\e{2},\e{12}\}
 \xRightarrow{\textrm{tensor product basis in matrix form}}\\
&\{1,\ii\}\otimes \Big\{
\begin{bmatrix}1&0\\0&1\end{bmatrix},
\begin{bmatrix}0&1\\1&0\end{bmatrix},
\begin{bmatrix}1&0\\0&-1\end{bmatrix},
\begin{bmatrix}0&-1\\1&0\end{bmatrix}\Big\}= \\
&\Big\{ \begin{bmatrix}1&0\\0&1\end{bmatrix},
\begin{bmatrix}0&1\\1&0\end{bmatrix},
\begin{bmatrix}1&0\\0&-1\end{bmatrix},
\begin{bmatrix}0&-1\\1&0\end{bmatrix},
\begin{bmatrix}\ii&0\\0&\ii\end{bmatrix},
\begin{bmatrix}0&\ii\\\ii&0\end{bmatrix},
\begin{bmatrix}\ii&0\\0&-\ii\end{bmatrix},
\begin{bmatrix}0&-\ii\\\ii&0\end{bmatrix}
\Big\}.
\end{split}\]
In the above calculations we have used the representation
$\{1,\ii\}$ for \cl{0}{1} and  real matrix representation for
\cl{2}{0}. In order to know which of matrices play the role of
base vectors in \cl{3}{0}, one must select three matrices that
satisfy the GA criterion, $\e{i}\e{j}+\e{j}\e{i}=2\delta_{i,j}$.
This would not be an easy task in general. Fortunately, from 
the theory above we know that vectors should be of the form
$\{1\otimes\e{1},\:1\otimes\e{2},\:\e{1}\otimes\e{12}\}$, i.~e.
they are made of the scalar of the \cl{0}{1} algebra multiplied by
vectors of the \cl{2}{0} or, alternatively, the vectors of the
\cl{0}{1} algebra multiplied by bivector of the \cl{2}{0}.
Therefore, instead of search which are the vectors in the above
table, we can write the table with graded elements in a canonical
order immediately,
\[\begin{split}
&\cl{3}{0}\cong\cl{0}{1}\otimes\cl{2}{0}
\xRightarrow{\textrm{base vectors}} \{1\otimes\e{1},\:1\otimes\e{2},\:\e{1}\otimes\e{12}\}
 \xRightarrow{\textrm{canonical order}}\\
\Big\{
 &\begin{bmatrix}1&0\\0&1\end{bmatrix},
\begin{bmatrix}0&1\\1&0\end{bmatrix},
\begin{bmatrix}1&0\\0&-1\end{bmatrix},
\begin{bmatrix}0&-\ii\\\ii&0\end{bmatrix},
\begin{bmatrix}\ii&0\\0&-\ii\end{bmatrix},
\begin{bmatrix}0&-1\\1&0\end{bmatrix},
\begin{bmatrix}0&\ii\\\ii&0\end{bmatrix},
\begin{bmatrix}\ii&0\\0&\ii\end{bmatrix}
\Big\}\\
&=\{\hat{1},\:\hat{\sigma}_1,\;\hat{\sigma}_3,
\;\hat{\sigma}_2,\;\hat{\sigma}_1\hat{\sigma}_2,
\;\hat{\sigma}_1\hat{\sigma}_3,\;\hat{\sigma}_2\hat{\sigma}_3,
\;\hat{\sigma}_1\hat{\sigma}_2\hat{\sigma}_3\},
\end{split}\]
where $\hat{\sigma}_i$ are the respective Pauli matrices. From the
last line we see that the direct product gives all eight elements
of \cl{3}{0} in a form of products of Pauli matrices.
\end{example}
%
\begin{unit}\label{evenoddalgebras}\textit{Even and odd subalgebras}.
Let us denote a set of even multivectors of \cl{p}{q} as $\cl{p}{q}^+$.
\red{buvusio sakinio dalis "consists of all even products of basis vectors" nera suprantama, nes neasiku kas tai yra even product.} 
An important property is that the even elements make a subalgebra, i.~e. the set of 
even elements is closed under GP. $\cl{p}{q}^+$ also is invariant
with respect to vector inversion, $\e{i}\rightarrow -\e{i}$. The
even subalgebras are isomorphic to lower order GA
algebras,~\cite{Porteous04}%p137
% Trautman in Clifford algebras and their representations (Elsevier 2006, vol 1 p. 518-530
% on page of formula (16) in very end of page writes Cl_{k,l} isomorphics to Cl^+_{k+1,l}
% which is a typo
\[\cl{p}{q+1}^+\cong\cl{p}{q},\quad\cl{p+1}{q}^+\cong\cl{q}{p}. \]
Similarly, a collection of odd grade elements of \cl{p}{q} is
denoted as $\cl{p}{q}^-$. Odd elements, however, do not form
subalgebra because they are not closed under GP. Any GA can be
split into sum of even and odd parts:
$\cl{p}{q}=\cl{p}{q}^{+}+\cl{p}{q}^-$.
\end{unit}
%
\begin{example}\label{evenoddCL13}
Two nonisomorphic GA's, namely \cl{1}{3} and \cl{3}{1}, are widely
used in physics to describe the relativity theory. From formulas
in~\ref{evenoddalgebras}  it follows that they share the same
even subalgebra, namely, $\cl{3}{1}^{+}\cong\cl{3}{0}$ and
$\cl{1}{3}^{+}\cong\cl{3}{0}$, where $\cl{3}{0}\cong\bbC(2)$ is
the algebra of classical (Newtonian) physics.

Using the even-grade isomorphism for lower grades, in a
similar way one can draw the following GA subordination diagram\\
%\setlength{\unitlength}{4pt}
\begin{picture}(20,25)
 \put(0,20){\cl{1}{3}} \put(0,10){\cl{3}{0}} \put(0,0){\cl{0}{2}}
 \put(20,20){\cl{1}{2}} \put(20,10){\cl{2}{0}} \put(20,0){\cl{0}{1}}
 \put(40,20){\cl{1}{1}}
 \put(60,20){\cl{1}{0}}
 %
 \put(10,20){\vector(1,0){5}} \put(30,20){\vector(1,0){5}} \put(50,20){\vector(1,0){5}}
 \put(3,19){\vector(0,-1){5}} \put(23,19){\vector(0,-1){5}}
 \put(3,9){\vector(0,-1){5}} \put(23,9){\vector(0,-1){5}}
\end{picture}\\[5pt]
which shows that the relativistic algebra \cl{1}{3} contains the
complex number algebra \cl{0}{1}, quaternion algebra \cl{0}{2},
Euclidean and Minkowski planes, \cl{2}{0} and \cl{1}{1}
respectively, and 3D Euclidean and Minkowskian algebras \cl{3}{0}
and \cl{1}{2}. Since we have started from even \cl{1}{3} algebra,
the elements of lower dimensional algebras should be selected from the
latter. For example the basis vectors of \cl{3}{0} are
$\{\sig{1}\equiv\e{12},\sig{2}\equiv\e{31},\sig{3}\equiv\e{41}\}$
the squares of which give $\sig{i}^2=+1$. The pseudoscalar of
\cl{3}{0} is
$I_3=\sig{1}\sig{2}\sig{3}=\e{21}\e{31}\e{41}=\e{1}\e{2}\e{3}\e{4}=I_4$,
which coincides with \cl{1}{3} pseudoscalar. The bivectors of
\cl{3}{0} are
$\Isig{1}=(\sig{1}\sig{2}\sig{3})\sig{1}=\sig{23}=\e{43}$,
$\Isig{2}=\e{42}$, $\Isig{3}=\e{32}$ the squares of which give
$-1$. From this follows that multiplication table of even grade
elements $\{1,\e{21},\e{31},\e{41},\e{32},\e{42},\e{43},I_4\}$ of
\cl{1}{3} coincides with \cl{3}{0} algebra multiplication table
constructed from elements
$\{1,\sig{1},\sig{2},\sig{3},\sig{12},\sig{13},\sig{23},I_3\}$.
Similarly, it is easy to show that the quaternion algebra
\cl{0}{2} may be represented by elements
$\{1,\sig{12},\sig{13},\sig{23}\}=\{1,\e{32},\e{42},\e{43}\}$, Any
two sets of bivectors may serve as basis vectors for~\cl{0}{2}.
\end{example}
%
\begin{example}
Relations of even subalgebras of unit~\ref{evenoddalgebras} can help 
to derive inverse formula for MV of lower dimensional algebra once
formula of higher dimension algebra is known. For example, from table~\ref{inverseLeft} the inverse MV for $p+q=4$ can be presented in the following form
\[ \m{A}^{-1}=\frac{\m{A}_{\bar 2,\bar 3}(\m{A}\m{A}_{\bar 2,\bar 3})_{\bar 1,\bar 4}}
{\m{A}\m{A}_{\bar 2,\bar 3}(\m{A}\m{A}_{\bar 2,\bar 3})_{\bar
1,\bar 4}}\] where ${\bar i}$ denotes the negated $i$-th grade (see unit~\ref{gradenegation}). If $\m{A}$ consists only of
even grades then $\m{A}$ is invariant to odd grade negations,
so that the above formula for even MV's reduces to
\[ \m{A}^{-1}=\frac{\m{A}_{\bar 2}(\m{A}\m{A}_{\bar 2})_{\bar 4}}
{\m{A}\m{A}_{\bar 2}(\m{A}\m{A}_{\bar 2})_{\bar 4}}\] In
example~\ref{evenoddCL13} we have found that even grade elements
of \cl{1}{3}  go to vectors and bivectors of \cl{3}{0} and the
pseudoscalar $I_4$ goes to pseudoscalar $I_3$. Thus, in \cl{3}{0}
algebra notation the last formula  should be rewritten as
\[ \m{A}^{-1}=\frac{\m{A}_{\bar1,\bar 2}(\m{A}\m{A}_{\bar1,\bar 2})_{\bar 3}}
{\m{A}\m{A}_{\bar 1,\bar 2}(\m{A}\m{A}_{\bar 1,\bar 2})_{\bar
3}}\] which coincides with the result for
$p+q=3$ of table~\ref{evenoddalgebras}.
\end{example}
%
\begin{unit}\textit{$8$-fold periodicity table}\label{8foldperiod}
classifies  all \cl{p}{q} algebras by their irreducible matrix
representations ({\'E}.~Cartan and R.~Bott). The aim of the
classification is three-fold. Firstly, the knowledge of
fundamental \cl{1}{0}, \cl{0}{1}, \cl{2}{0}, \cl{1}{1} and
\cl{0}{2} algebras to which we may add the real number algebra
$\cl{0}{0}=\bbR$ determines all other \cl{p}{q} algebras (see
unit~\ref{constructionOfRepresentation}). Secondly, the
classification demonstrates that all GA's for which $(p-q)\mod{8}$ is the same number share similar properties. So, the table may be used to efficiently generate
matrix representations of algebras with $p,q\ge 8$.
Thirdly, the classification table is helpful if one wants to
transform equations given in matrix form, e.~g. equations
encountered physics, to GA form. The multiplication of $n\times n$
matrices $\bbR(n)$, $\bbC(n)$ or $\bbH(n)$ then is equivalent to
the geometric product in GA. The two tables below relate \cl{p}{q}
algebras with $n\times n$ matrices (the third form of the same table is given in unit~\ref{complexification}).
%
\begin{center}\label{TableOfPeriodicity1}
%\setlength{\extrarowheight}{1ex}
%\arraycolsep=0pt
\parbox[t]{0.9\textwidth}{
\textit{Period-$8$ table of irreducible
matrix representations of GA (called the fundamental) for
Cl$_{p,q}$, where $p$ and  $q$ indicate positive and negative signature.}}
%\setlength{\tabcolsep}{1pt}
%\resizebox{\textwidth}{!}{%
%\renewcommand*{\arraystretch}{1}
%\noalign{\vskip -3pt}\hphantom{m}

% Acaus preambule
%\begin{tabular}{x{0.1em}| >{$}l<{$} >{$}l<{$} >{$}l<{$} >{$}l<{$}
%>{$}l<{$} >{$}l<{$} >{$}l<{$} >{$}l<{$}}
%\diag{.1em}{.1em}{\mbox{$\!\!p\,$}}{\mbox{\hphantom{n}\smash{$q$}}}
%Dargio preambule
\begin{tabular}{>{$}l<{$} | >{$}l<{$} >{$}l<{$} >{$}l<{$} >{$}l<{$}
>{$}l<{$} >{$}l<{$} >{$}l<{$} >{$}l<{$}}
% tabular body
&0&1&2&3&4&5&6&7\\
\cline{2-9}
0 &\bbR&\bbC&\bbH&{}^2\bbH&\bbH(2)&\bbC(4)&\bbR(8)&{}^2\bbR(8)\\
1 &{}^2\bbR&\bbR(2)&\bbC(2)&\bbH(2)&{}^2\bbH(2)&\bbH(4)&\bbC(8)&\bbR(16)\\
2 &\bbR(2)&{}^2\bbR(2)&\bbR(4)&\bbC(4)&\bbH(4)&{}^2\bbH(4)&\bbH(8)&\bbC(16)\\
3 &\bbC(2)&\bbR(4)&^2\bbR(4)&\bbR(8)&\bbC(8)&\bbH(8)&^2\bbH(8)&\bbH(16)\\
4 &\bbH(2)&\bbC(4)&\bbR(8)&{}^2\bbR(8)&\bbR(16)&\bbC(16)&\bbH(16)&^2\bbH(16)\\
5 &{}^2\bbH(2)&\bbH(4)&\bbC(8)&\bbR(16)&^2\bbR(16)&\bbR(32)&\bbC(32)&\bbH(32)\\
6 &\bbH(4)&{}^2\bbH(4)&\bbH(8)&\bbC(16)&\bbR(32)&{}^2\bbR(32)&\bbR(64)&\bbC(64)\\
7&\bbC(8)&\bbH(8)&{}^2\bbH(8)&\bbH(16)&\bbC(32)&\bbR(64)&{}^2\bbR(64)&\bbR(128)
\end{tabular}
%
\end{center}
The left superscript denotes direct sum of matrices, e.~g.,
$^2\bbR(2)\equiv\bbR(2)\oplus\bbR(2)$.
%
\vspace{3mm}
\begin{center}
\parbox[t]{0.9\textwidth}{
\textit{8-periodicity classification listed according to value of
$(q-p)\mod 8$} [according to~\cite{Coquereaux09} but corrected].}
\begin{tabular}{l|cccc}
  $(q-p) \mod 8 $ & $ 0 $ & $ 1 $ & $ 2 $ & $ 3 $\\
\hline
$\cl{p}{q}$ matrix & $\bbR(2^{n\over 2})$ & $\bbC(2^{{n-1 \over 2}})$ & $\bbH(2^{n-2\over 2})$ & $\bbH(2^{{n-3 \over 2}}) \oplus \bbH(2^{{n-3 \over 2}})$\\
\omit \\[12pt]
$(q-p)\mod 8$  & $ 4 $ & $ 5 $ & $ 6 $ & $ 7 $ \\
\hline $\cl{p}{q}$ matrix & $\bbH(2^{n-2\over 2})$ & $
\bbC(2^{{n-1 \over 2}}) $ & $ \bbR(2^{{n\over 2}}) $ &
 $\bbR(2^{{n -1 \over 2}}) \oplus \bbR(2^{{n-1 \over 2}})$
\end{tabular}
\end{center}
Here as usually $n=p+q$. From the above classification it follows
that $\cl{n}{n}\cong\bbR(2^{n\over2})$ and representations of 
\cl{p}{q} and \cl{q}{p} differ, $\cl{p}{q}\not\cong\cl{q}{p}$.
\end{unit}
%
\begin{example}
It is easy to check that, for example, when $p=4$ and $q=5$ we
have $((q-p)\mod 8) = 1$ and, therefore, matrix representation is
$\bbC(2^{\tfrac{5+4-1}{2}})=\bbC(16)$, which matches the
corresponding entry in the period-8 table.

For larger values, say, $p=9$ and $q=12$ we get $((q-p)\mod 8)
=3$. The matrix representation for these $p$ and $q$ thus is
$\bbH(2^{(9+12-3)/2})\oplus\bbH(2^{(9+12-3)/2})=\bbH(2^{9})\oplus\bbH(2^{9})=\bbH(512)\oplus\bbH(512)=\bbR(16)\otimes(\bbH(32)\oplus\bbH(32))=
\bbR(16)\otimes\bbR(16)\otimes(\bbH(2)\oplus\bbH(2))$. This is
compatible with the explicit table entry at $p=1$ and $q=4$
multiplied by  additional matrix factor
$\bbR(16)\otimes\bbR(16)=\bbR(256)$. The result also can be
cross checked using GA isomorphism rules for $\cl{p+8}{q}$ and
$\cl{p}{q+8}$ of \ref{periodicityTableExtension} unit below, i.~e.
$\cl{9}{12}\cong
\bbR(16)\otimes\cl{1}{12}\cong\bbR(16)\otimes\bbR(16)\otimes\cl{1}{4}$.
\end{example}
\vspace{2mm} Representation of GA basis elements by matrices is not
unique. Frequently the different sets of representation matrices
may be related by some unitary transformation. Examples of
particular matrix representations of simple GA's may be found in
unit~\ref{GAreps}, where individual matrices are described.
%
\begin{example} Transformation of  matrix to MV form. Let the
Hermitian matrix is
\[\hat{M}=\begin{bmatrix}
\alpha_1&\beta_1+\ii\beta_2\\
\beta_1-\ii\beta_2&\alpha_2\\
\end{bmatrix},\quad\alpha_i,\beta_i\in\bbR  \,.\]
This  matrix belongs to $\bbC(2)$ and as the periodicity
table~\ref{GAreps} tells it may be rendered to either \cl{3}{0} or
\cl{1}{2} algebra. The basis representation matrices of \cl{3}{0}
consists of the unit matrix $\hat{1}$, three Pauli matrices
$\hat{\sigma}_i$, their products
$\hat{\sigma}_{ij}=\hat{\sigma}_i\hat{\sigma}_j$, and pseudoscalar
$\hat{I}=\hat{\sigma}_1\hat{\sigma}_2\hat{\sigma}_3$ (see
Example~\ref{excl30}). With $\textrm{Tr}(\hat{M})/2$,
$\textrm{Tr}(\hat{M}\hat{\sigma}_i^{\dagger})/2$,
$\textrm{Tr}(\hat{M}\hat{\sigma}_{ij}^{\dagger})/2$ and
$\textrm{Tr}(\hat{M}\hat{I}^{\dagger})/2$ one finds that real,
nonzero coefficients in the MV are at scalar and basis vectors
$\hat{\sigma}_i\leftrightarrow\e{i}$ only. Thus, the matrix
$\hat{M}$ in  \cl{3}{0} algebra may be represented as MV,
\[\hat{M}\rightarrow\m{M}=\frac12(\alpha_1+\alpha_2)+\beta_1\e{1}-\beta_2\e{2}+\frac12(\alpha_1-\alpha_2)\e{3},\]
which is the sum of  scalar and vector in coordinate system
represented by Pauli matrices.
\end{example}
%
\section{GA groups}
Some basic facts from group theory which are directly linked with
Clifford groups are singled out here.
\begin{unit}\textit{Definition of group}. \label{Definitionofgroup}A set of elements
(e.~g. multivectors) $A,B,C,\dots\in\mathsf{G}$ together with a
law of composition $\circ$, called multiplication (we will replace
it by GP), makes a group~$\mathsf{G}$ if the following axioms are
satisfied:
\begin{enumerate}
\item $A\circ B = C$ gives an element of the same group,
$C\in\mathsf{G}$. In general, group elements do  not commute,
i.~e., $A\circ B\ne B\circ A$. If they commute, the group is
called \textit{Abelian}. \item $(A\circ B)\circ C = A\circ (B\circ
C)$, i.~e. $\circ$ is the associative operation. \item $A \circ E
= E\circ  A = A$, i.~e., there exists a unit (multiplicative
identity) element~$E$, also denoted by~$1$, which belongs to
$\mathsf{G}$. \item $A^{-1}\circ A=A\circ A^{-1}=E$, i.~e., for
every element $A$ there always exists an multiplicative inverse $A^{-1}$.
\end{enumerate}
A subset of elements in $\mathsf{G}$ which is closed under
operation $\circ$ makes a subgroup in $\mathsf{G}$.
%
\begin{figure}[t]
\centering
\includegraphics[width=4.2cm]{figures//grComponents2}a)\quad
\includegraphics[width=2.8cm]{figures//grComponents1NotSimplyConnected}b)\quad
\includegraphics[width=2.8cm]{figures//grComponents1SimplyConnected}c)
\caption{Examples of parameter manifolds for various groups.
a)~The manifold $\cM$ which consists of two non-connected
components one of which necessarily contains the identity element
$1$ and the path for \textit{non-connected} group.  The identity
cannot be reached from the point $B$ located in the other manifold
component. Example of groups with two component manifold are:
$\mathsf{GL}(n,\bbR)$, $\mathsf{O}(n)$, $\mathsf{SO}(n,1)$. The
 $\mathsf{O}(n,1)$ is four-component manifold.
b)~\textit{Connected}, but not simply connected group manifold
$\cM$. Here the loop cannot be shrunk to the identity element.
Groups with connected but not simply connected manifold are
$\mathsf{GL}(n,\bbC)$, $\mathsf{SL}(n,\bbR)$, $\mathsf{SO}(n)$,
$\mathsf{U}(n)$ of any $n$, and $\mathsf{SO}(n,1)$ for $n\ge 2$.
c)~\textit{Simply connected} group manifold. Groups with simply
connected manifold are $\mathsf{SL}(n,\bbC)$, $\mathsf{SU}(n)$,
$\mathsf{SO}(1,1)$. Definitions of groups of type
$\mathsf{O}(n,1)$, $\mathsf{SO}(n,1)$ etc are given in
unit~\ref{classgroup} \label{fig:connectedManifold}}
\end{figure}

In \textit{continuous group} the elements of group are analytic
functions of parameter(s) in a such way that the identity element
is usually characterized by all parameters setting to zero. For example,
rotations in 3D space are represented by 3~parameters (spatial
angles), $R(\alpha,\beta,\gamma)$ such that $R(0,0,0)=E$. To
simplify notation the parameters of the group elements are
frequently omitted. The number of independent parameters $r$ is
called the \textit{dimension of group}.

An $r$-parameter \textit{Lie group} $\m{G}$ at the same time
represents an $r$-dimensional differentiable manifold in such a
way that the group composition $\m{G}\circ\m{G}\rightarrow\m{G}$
is a smooth map on the manifold. Therefore, the parameters of Lie
group make up a differentiable manifold of dimension~$r$. A Lie
group is called \textit{compact} if its manifold is a closed and
bounded submanifold of $\bbR^m$ for some finite $m$. It follows
that for a matrix group each of the matrix element must be
bounded.

A group is called (pathwise) \textit{connected} if an identity
element~$1$ can be reached from an arbitrary group
element~$A(\alpha,\beta\dots)$ by a continuous variation of the
parameters $\{\alpha,\beta\dots\}$, i.~e. coordinates of the
point $A(\alpha,\beta,\dots)$ on the manifold $\cM\ni A(\alpha,\beta,\dots)$ as shown in
Fig.~\ref{fig:connectedManifold}b and c. The group is
non-connected if such path for some points does not exist, for
example for a point $B$ in Fig.\ref{fig:connectedManifold}a. For a
group to be \textit{simply connected}, it is required that every
loop in group $\mathsf{G}$ manifold $\cM$ can be continuously shrunk
 to a point on $\cM$,~Fig.~\ref{fig:connectedManifold}c.

Elements of Lie groups typically are constructed by exponentiation
of generators of Lie algebra which are specific constant-matrices,
or differential operators of special kind (for example
$x\partial_y-y\partial_x$) multiplied by group parameters, see unit~\ref{expdiffmap}.
\end{unit}
%
\begin{unit}
\textit{Table of classical (matrix) Lie groups}.\label{classgroup}
The classical groups include general linear groups
$\mathsf{GL}(n,\bbF)$, special linear groups $\mathsf{SL}(n,\bbF)$
over $\bbF=\bbR,\bbC,\bbH$ and their continuous subgroups that
keep invariant some non-degenerate bilinear or sesquilinear forms.
The group name and classification is based on these invariant forms
(also called metric), which are unaltered under action of the
group elements. The bilinear form defined on a real $\bbR$,
complex $\bbC$ or quaternionic $\bbH$ vector space  $\cV^n$ is a
linear map $B(\bx,\by)$ from pair of vectors $\bx,\by\in\cV^n$
into a real, complex or quaternionic number, i.~e. $B:
\cV^n\otimes\cV^n \rightarrow \bbF$. The form is called bilinear
if $B(\alpha\bx,\beta\by)=\alpha B(\bx,\by)\beta$, for
$\bx,\by\in\cV^n$ and $\alpha,\beta\in\bbF$, and sesquilinear if
$B(\alpha\bx,\beta\by)=\bar\alpha B(\bx,\by)\beta$. Here
$\bar\alpha$ denotes field conjugation operation, i.~e. complex or
quaternionic conjugate. The linearity applies to both vector
arguments $\bx$ and $\by$. Every matrix group always is a Lie
group. However, there are Lie groups, which are not (finite
dimensional) matrix groups, the famous example being the
Heisenberg group.
% Andrew Baker, "Introduction to matrix groups and their applications, 2000, p.12
% www.maths.gla.ac.uk/~ajb


The table below summarizes main properties of classical groups
that are important in the context of GA groups. For more
information reader is referred to unit~\ref{groupproperties} and
books \cite{wybourne1974classical,knapp2002lie,Helgason78,Hall03}.
%Helgason \cite{Helgason78} p.449. \cite{Hall03} p.16.
\begin{center}
\begin{tabular}{>{$}l<{$} >{$}l<{$} >{$}l<{$} >{$}l<{$} >{}l<{} >{$}l<{$}}
\textrm{Group}&\textrm{Invariant}&\textrm{Condition}&\textrm{Order}&\textrm{Compact} \\
&\textrm{metric}&\textrm{for matrix}&\textrm{of group} & \\
\hline
% general linear groups, no invariant bilinear form
\mathsf{GL}(n,\begin{tabular}[t]{@{}c@{}}\bbR\\ \bbC \end{tabular}
)& - & \begin{tabular}[t]{@{}l@{}} $\det A \neq 0$\\ $\det A \neq
0$ \end{tabular} & \begin{tabular}[t]{@{}l@{}} $n^2$\\ $2n^2$
\end{tabular} & \begin{tabular}[t]{@{}c@{}}no\\ no \end{tabular} &
\begin{tabular}[t]{@{}c@{}}\\  \end{tabular}
\\
\mathsf{SL}(n,\begin{tabular}[t]{@{}c@{}}\bbR\\ \bbC
\end{tabular}) & - & \begin{tabular}[t]{@{}l@{}} $\det A =1$\\
$\det A =1$ \end{tabular} &
  \begin{tabular}[t]{@{}l@{}} $n^2-1$\\ $2(n^2-1)$ \end{tabular} &\begin{tabular}[t]{@{}c@{}}no\\ no \end{tabular}
\\
% orthogonal,  real matrices, symmetric bilinear form
  \mathsf{O}(n) & \bx^{\mathrm{T}}\by & A^{\mathrm{T}}A = \hat{1}_n & \tfrac12 n(n-1) &yes &
  \\
\mathsf{SO}(n) &  \bx^{\mathrm{T}}\by &
\mathllap{\Bigl\{}\begin{tabular}[c]{@{}l@{}} $A^{\mathrm{T}}A =
I_n$\\ $\det A = 1$ \end{tabular} &  \tfrac12 n(n-1) &yes
\\
\mathsf{O}(P,Q) & \bx^{\mathrm{T}}L\, \by & A^{\mathrm{T}} LA = L &
\tfrac12 n(n-1) &no
\\
\mathsf{SO}(P,Q) & \bx^{\mathrm{T}} L\,\by &
\mathllap{\Bigl\{}\begin{tabular}[c]{@{}l@{}} $ A^{\mathrm{T}} L A
= L$\\ $\det A = 1$ \end{tabular} & \tfrac12 n(n-1) & no
  \\
%unitary, complex matrices, symmetric bilinear form
\mathsf{U}(n) & \bx^{\dagger}\by & A^{\dagger}A = \hat{1}_n & n^2
&yes &
  \\
\mathsf{SU}(n) & \bx^{\dagger}\by &
\mathllap{\Bigl\{}\begin{tabular}[c]{@{}l@{}} $A^{\dagger}A =
I_n$\\ $\det A = 1$ \end{tabular} &  n^2-1 &yes
  \\
\mathsf{U}(P,Q) & \bx^{\dagger} L\,\by & A^{\dagger} L A = L & n^2
&no
\\
\mathsf{SU}(P,Q) & \bx^{\dagger} L\,\by &
\mathllap{\Bigl\{}\begin{tabular}[c]{@{}l@{}} $A^{\dagger} L A =
L$\\ $\det A = 1$ \end{tabular} & n^2-1 &no
\\
% sympletic, antisymmetric bilinear form
\mathsf{Sp}(2n,\begin{tabular}[t]{@{}c@{}}\bbR\\ \bbC
\end{tabular} ) &  \begin{tabular}[t]{@{}l@{}} $\bx^{\mathrm{T}}
J\, \by$\\ $\bx^{\dagger} J\,\by$ \end{tabular} &
\begin{tabular}[t]{@{}l@{}} $A^{\mathrm{T}} J A = J$\\
$A^{\dagger} J A = J$ \end{tabular} &
\begin{tabular}[t]{@{}l@{}} $n (2n+1)$\\ $2n (2n+1)$ \end{tabular}
&\begin{tabular}[t]{@{}c@{}}no\\ no \end{tabular}
%  \\
%\mathsf{SSp}(2n,\begin{tabular}[t]{@{}c@{}}\bbR\\ \\ \bbC
%\end{tabular} ) & \begin{tabular}[t]{@{}l@{}} $\bx^{\mathrm{T}} J
%\by$\\ \\ $\bx^{\dagger} J\by$ \end{tabular} &
%\begin{tabular}[t]{@{}l@{}}$\mathllap{\Bigl\{}\begin{array}{l}A^{\mathrm{T}}
%J A = J\\ \det A = 1 \end{array}$\\
%%$\mathllap{\Bigl\{}\begin{array}{l}A^{\dagger} J A = J\\ \det A =
%1 \end{array}$\end{tabular} & \begin{tabular}[t]{@{}l@{}} $n
%(2n+1)$\\ \\ $2n (2n+1)-1$ \end{tabular}
%&\begin{tabular}[t]{@{}c@{}}no\\ \\ no \end{tabular} &
%\begin{tabular}[t]{@{}c@{}}\\ \\  \end{tabular}&
%\begin{tabular}[t]{@{}c@{}}?\\ \\ yes \end{tabular}
%  \\
\end{tabular}
\end{center}
Here $A\equiv\hat{A}$ is a group element transformation
represented as a $n\times n$ matrix;  $\bx$ and $\by$ are
arbitrary column vectors of length $n$ that are transformed in $n$
dimensional vector space; $\bx^{\textrm{T}}$ represents a row
(transposed) vector and $\bx^{\dagger}$ is a complex conjugate row
vector. The letters $L$ and $J$ denote special matrices
$L=\Bigl(\begin{smallmatrix}-\hat{1}_P&0\\0&\hat{1}_Q\end{smallmatrix}\Bigr)$
and
$J=\Bigl(\begin{smallmatrix}0&\hat{1}_n\\-\hat{1}_n&0\end{smallmatrix}\Bigr)$,
where $\hat{1}_n$ is $n\times n$ unit matrix. Similarly,
$\hat{1}_P$ and $\hat{1}_Q$ are unit matrices of dimension $P$ and $Q$.
\end{unit}
%
\begin{unit}\textit{More properties of classical groups related with GA}.
\label{groupproperties}
%\blue{Mano kompiliatorius taip uzrasytu opciju su [..] nesupranta.
%Todel jas uzkomentavau ir parasiau taip, kaip rekomenduojama
%Goossens-Mittelbach-Samarin visiems zinomoje knygoje}
\renewcommand{\descriptionlabel}[1]%
      {\hspace{\labelsep}\textsf{#1}}
\begin{description}%[style=multiline,itemindent=0.em,labelsep=0.2em,itemsep=2pt,
%font=\normalfont,before={\renewcommand\makelabel[1]{##1}}]
\item[{\it Groups without invariant bilinear form}.]
\item[$\mathsf{GL}(n,\bbR)$,] $\mathsf{GL}(n,\bbC)$ and
$\mathsf{GL}(n,\bbH)$ denote the \textit{general linear group} of
all $n\times n$ invertible (i.~e. with $\det A \neq 0$) matrices
$A$ with real,  complex and quaternionic entries. The number of
entries, equivalently, the number of basis matrices is $n^2$. It
means that $\mathsf{GL}(n,\bbR)$ has $n^2$ parameters while
$\mathsf{GL}(n,\bbC)$ and $\mathsf{GL}(n,\bbH)$, respectively,
$2n^2$ and $4n^2$ parameters. The general linear groups do not
have invariant bilinear form.
%
\item[$\mathsf{SL}(n,\bbR)$,] $\mathsf{SL}(n,\bbC)$ and
$\mathsf{SL}(n,\bbH)$ are \textit{special linear} subgroups of
corresponding general linear groups with determinants restricted
to $1$. The determinant of quaternionic matrix is calculated as
determinant of square matrix regarded as a $2n\times 2n$ complex
matrix. %~\cite{Porteous95}~p.42;104.
The number of parameters are $n^2-1$, $2(n^2-1)$ and $4n^2-1$,
respectively. The special linear groups do not preserve invariance
of bilinear forms as well.\\[2pt]
%\red{Argi apribojimas kvaternionams neturetu
%sumazinti 3 laisves laipsnius? Kvaternionai turi 3 menamus
%vienetus. Tai jei kompleksiniams skaiciams sumazeja 1, tai
%kvaternionams turetu sumazeti 3?}\blue{Pasitikrinau. Porteous
%knygoje psl.104 kur taip pat parasyta $n^2-1$, $2n^2-2$, $4n^2-1$.
%Del $\bbR,\bbC,\bbH$ man noretusi rasyti: $1(n^2-1)$, $2(n^2-1)$,
%$4(n^2-1)$. Visdelto, radau knyga R. Gilmore "Lie groups,..."
%psl.48, kur parasyta $4n^2-1$ ir prideta "no metric". Taigi, pas
%mane  buvo gerai.}
%
\item[{\it Symmetric invariant bilinear forms, real matrices}.]
\item[$\mathsf{O}(n)$]$\equiv\mathsf{O}(n,\bbR)$ is an
\textit{orthogonal group} of real invertible $n\times n$ matrices
$A\in \mathsf{GL}(n,\bbR)$. Orthogonal group deals with Euclidean
(real, symmetric, bilinear) form,
\[
\bx^{\mathrm{T}}\by=[x_1,x_2,\ldots,
x_n]\begin{bmatrix}y_1\\y_2\\\vdots\\y_n\end{bmatrix}=\sum_{i=1}^n
x_i y_i.
\]
Requirement that action of $A$ leaves the metric invariant,
\[(\bx^\prime,\by^\prime)=(A\bx,A\by)=(A\bx)^{\mathrm{T}} (A\by)=\bx^{\mathrm{T}}(A^{\mathrm{T}}
A)\by=(\bx,\by),
\]
yields the condition $A^{\mathrm{T}} A = \hat{1}_n$.  Thinking of a matrix
as given by $n^2$ coordinate functions, the above condition is
equivalent to $n(n+1)/2$ constrains. The remaining
$n^2-n(n+1)/2=\tfrac{1}{2}n(n-1)$ free variables constitute free
parameter space of $\mathsf{O}(n)$. Since $\det
A=\pm 1$, this group is not simply connected.
%
\item[$\mathsf{SO}(n)$] is called a \textit{special othogonal
group}. It is a subgroup of $\mathsf{O}(n)$, and includes only
matrices with determinants restricted to $+1$. It has the same
number of parameters.
%
\item[$\mathsf{O}(P,Q)$] is an \textit{pseudo-orthogonal group} of
real $n\times n$ matrices $A\in \mathsf{GL}(n,\bbR)$, $n=P+Q$,
which deal with Lorentzian form,
\[
(\bx,\by)= \bx^{\mathrm{T}} L\by= \bx^{\mathrm{T}}
\begin{bmatrix}-\hat{1}_P&0\\0&\hat{1}_Q\end{bmatrix}\by=-\sum_{i=1}^P
x_i y_i +\sum_{i=P+1}^{P+Q=n} x_i y_i.
\]
The requirement that action of $A\in\mathsf{O}(P,Q)$ leaves the
metric invariant,
\[(\bx^\prime,\by^\prime)=(A\bx,A\by)=(A\bx)^{\mathrm{T}}L(A\by)=
\bx^{\mathrm{T}}(A^{\mathrm{T}}LA)\by=\bx^{\mathrm{T}}L\by,\]
yields the  condition $A^{\mathrm{T}}L A = L$. This gives
$\tfrac{1}{2}n(n-1)$ free variables which constitute free
parameter space of $\mathsf{O}(P,Q)$. Note, that definitions of
classical groups use opposite signature convention, i.~e., the
coordinates with the first $P$ indices have negative square.
%
\item[$\mathsf{SO}(P,Q)$] is a \textit{pseudo-special orthogonal}
group, the subgroup of $\mathsf{O}(P,Q)$ with restriction $\det A
=1$ and the same number of free parameters. \\[3pt]

\item[{\it Symmetric invariant bilinear forms, complex matrices}.]
\item[$\mathsf{U}(n)$] $\equiv\mathsf{U}(n,\bbC)$ is a
\textit{unitary group} of complex invertible $n\times n$ matrices
$A\in \mathsf{GL}(n,\bbC)$. This group deals with Hermitian
(complex, symmetric, bilinear) forms
\[
\bx^{\dagger}\by=[x_1^*,x_2^*,\ldots,
x_n^*]\begin{bmatrix}y_1\\y_2\\\vdots\\y_n\end{bmatrix}=\sum_{i=1}^n
x_i^* y_i,
\]
where $\dagger$ denotes transposition and complex conjugation.
The requirement that action of $\mathsf{U}(n)$ leaves Hermitian metric invariant
\[(\bx^\prime,\by^\prime)=(A\bx,A\by)=(A\bx)^{\dagger}(A\by)
=\bx^{\dagger}(A^{\dagger}A)\by=(\bx,\by)\] yields the condition
$A^{\dagger} A = \hat{1}_n$, which restricts number of (real)
parameters to~$n^2$.
%
\item[$\mathsf{SU}(n)$] denotes the \textit{special unitary
group}, the subgroup of $\mathsf{U}(n)$ with determinant
restricted to $1$. In physics, the groups $\mathsf{SU}(2)$,
$\mathsf{SU}(3)$, $\mathsf{SU}(4)$ are very popular. The Pauli
matrices  $2\times 2$, Gell-Mann matrices of dimension $3\times 3$
and Gell-Mann-type matrices of dimension $4\times 4$, the number
of which is $n^2-1$, i.~e.~$3$, $8$ and $15$, respectively,
typically are used as group generators (see unit~\ref{Liealgebra}
below).
%
\item[$\mathsf{U}(P,Q)$] is a \textit{pseudo-unitary group} of
complex $n\times n$ matrices $A\in \mathsf{GL}(n,\bbC)$. This
unitary group deals with Hermitian-Lorentzian metric (signature
convention again is opposite to GA),
\[
(\bx,\by)= \bx^{\dagger}L\by= -\sum_{i=1}^P x_i^* y_i
+\sum_{i=P+1}^{P+Q=n} x_i^* y_i.
\]
Requirement that action of $\mathsf{U}(P,Q)$ would leave the metric invariant
\[(\bx^\prime,\by^\prime)=(A\bx,A\by)=(A\bx)^{\dagger}L (A\by)=
\bx^{\dagger}(A^{\dagger}L A)\by=\bx^{\dagger} L\by\]
yields the condition $A^{\dagger}L A = L$. It restricts the number of free real
variables to $n^2$ which constitutes free parameter space of $\mathsf{U}(P,Q)$.
%
\item[$\mathsf{SU}(P,Q)$] is a \textit{pseudo-special unitary} group,
 the subgroup of $\mathsf{U}(P,Q)$ with $\det A =1$ and $n^2-1$ free real parameters.
 \\[3pt]
\item[{\it Skew-symmetric invariant bilinear forms}.]
\item[$\mathsf{Sp}(2n,\bbR)$,] $\mathsf{Sp}(2n,\bbC)$ and
$\mathsf{Sp}(n,\bbH)$ are \textit{symplectic groups} of $2n\times
2n$ matrices with real, complex and quaternionic entries. They
preserve the  sympletic metric (skew-symmetric bilinear form),
\[
(\bx,\by)= \bx^{\dagger}J \by= \bx^{\dagger}
\begin{bmatrix}0&\hat{1}_n\\-\hat{1}_n&0\end{bmatrix}\by=\sum_{i=1}^n
x_i^* y_{2n-i} -\sum_{i=1}^{n} x_{2n-i}^* y_i.
\]
The elements necessarily belong to even dimensional vector space, because in odd dimensions all matrices turns to zero.
The requirement that action of $\mathsf{Sp}(2n)$ leaves symplectic
metric invariant,
\[(\bx^\prime,\by^\prime)=(A\bx,A\by)=(A\bx)^{\dagger}J(A\by)=
\bx^{\dagger}(A^{\dagger}J A)\by=(\bx,\by),\] gives the condition
$A^{\dagger}J A = J$, which restricts a number of free parameters
of $\mathsf{Sp}(2n,\bbR)$ to $n(2n+1)$. For $\mathsf{Sp}(2n,\bbC)$
and $\mathsf{Sp}(2n,\bbH)$ a number of  free parameters is
$2n(2n+1)$ and $n(2n+1)$, respectively.
%see Gilmore p.48
\end{description}
The list of presented classical groups is not exhaustive. In
particular,  it lacks groups with invariant sesquilinear skew
Herminian bilinear forms, i.~e., an entire class of classical
groups over quaternionic field. Below we explicitly define
quaternionic groups that are important for GA group
theory~\cite{Trautman06},
\[\begin{split}
 &\mathsf{Sp}(2,\bbH)=\{\hat{A}\in\bbH(2)\,|\,\hat{A}^{\dagger}\hat{A}=\hat{1}\},\\
 &\mathsf{Sp}(1,1,\bbH)=\{\hat{A}\in\bbH(2)\,|\,\hat{A}^{\dagger}\hat{\sigma}_z\hat{A}=\hat{\sigma}_z\},\\
 &\mathsf{SL}(n,\bbH)=\{\hat{A}\in\bbH(n)\,|\,\det\hat{A}=1\}.
\end{split}\]
Here the dagger $\dagger$ means matrix transposition followed by
quaternionic conjugation, whereas $\hat{\sigma}_z$ denotes the
corresponding diagonal Pauli matrix.
%\cite{Trautman06} p.6
\end{unit}
%
\begin{unit}\textit{Local group isomorphisms}.\label{groupisomorph}
A Lie group can have several disconnected pieces as shown in
figure~\ref{fig:connectedManifold}a. The local isomorphism deals
only with the components in the vicinity of the identity point
$1$. In general, the local isomorphism between groups only implies that the
groups are, strictly speaking, homomorphic. The table below lists local group isomorphisms.
\[\begin{split}
n_{\mathrm{G}}=1\qquad&\mathsf{SO}(2,\bbR)\cong\mathsf{U}(1,\bbC)\\
n_{\mathrm{G}}=3\qquad&\mathsf{SL}(1,\bbH)\cong\mathsf{SU}(2,\bbC)\cong\mathsf{SO}(3,\bbR)\cong
\mathsf{U}(1,\bbH)\\
&\mathsf{Sp}(2,\bbR)\cong\mathsf{SL}(2,\bbR)\cong
\mathsf{SU}(1,1,\bbC)\cong\mathsf{SO}(2,1,\bbR)\\
n_{\mathrm{G}}=6\qquad&\mathsf{SO}(4,\bbR)\cong\mathsf{SU}(2,\bbC)\otimes\mathsf{SU}(2,\bbC)\\
&\mathsf{SL}(2,\bbC)\cong\mathsf{SO}(3,1,\bbR)\\
&\mathsf{SL}(2,\bbR)\otimes\mathsf{SL}(2,\bbR)\cong\mathsf{SO}(2,2,\bbR)\\
n_{\mathrm{G}}=10\qquad&\mathsf{Sp}(4,\bbR)\cong\mathsf{SU}(4,\bbC)\cong\mathsf{SO}(3,2,\bbR)\\
n_{\mathrm{G}}=15\qquad&\mathsf{SU}(4,\bbC)\cong\mathsf{SO}(6,\bbR)\\
&\mathsf{SU}(2,2,\bbC)\cong\mathsf{SO}(4,2,\bbR)\\
&\mathsf{SL}(4,\bbR)\cong\mathsf{SO}(3,3,\bbR)\\
&\mathsf{SL}(2,\bbH)\cong\mathsf{SO}(5,1,\bbR)\\
&\mathsf{SL}(3,1,\bbR)\cong\mathsf{SO}(6,\bbR)\\
\end{split}\]
%Gilmore~\cite{Gilmore74}~p.52, Steinkirch~\cite{Steinkirch11}~p.98
Here the $n_G$ denotes a number of group parameters (which is the
same as a number of Lie algebra generators). In the context of
\cl{p}{q} algebras and spinor groups it coinsides with the
number of bivectors, $n_G=\frac{(p+q)!}{2!(p+q-2)!}=\tfrac{1}{2}n(n-1)$; $n=p+q$.
%\red{yra sukeista vietomis p ir q skaiciai, t.y vietoje SO(3,1)
%ten SO(1,3)? } \blue{Spin grupems $Spin(p,q)=Spin(q,p)$. Ziur.
%toliau. Nurasiau is uzkomentuotu knygu}
\end{unit}
%
\begin{example}\label{SO2R}
A subset of elements $\{1,\e{12}\}$  of  \cl{2}{0} is related with
spin group $\m{Spin}_{2,0}\cong\mathsf{SO}(2,\bbR)$ (see
unit~\ref{TypesOfGroups}). It is a single parameter group
($n_G=1$, see~\ref{groupisomorph}) of $2\times 2$ matrices
$A(\beta)=\left[\begin{smallmatrix}\cos\beta&\sin\beta\\-\sin\beta&\cos\beta\end{smallmatrix}\right]$
that  describe rotation in plane by angle $\beta$ (the group parameter).
It is easy to check that group axioms are satisfied if group
multiplication~$\circ$ is identified with matrix multiplication:
\begin{enumerate}
 \item The product of two group elements (rotations)
gives $A(\alpha)A(\beta)=A(\alpha+\beta)\in\mathsf{SO}(2,\bbR)$ the
other rotation by angle $(\alpha+\beta)$.
 \item The associativity property is satisfied due to matrix multiplication properties.
 \item  The identity element is represented by unit $2\times 2$ matrix $\hat{1}$.
 \item Every rotation has inverse, $A^{-1}(\beta)=A(-\beta)$.
\end{enumerate}
The generator of Lie algebra may be identified with matrix
\[
\hat{e}_{12}=\frac{\partial A(\beta)}{\partial \beta} A^{-1}=
\begin{bmatrix}-\sin\beta&\cos\beta\\-\cos\beta&-\sin\beta\end{bmatrix}
\begin{bmatrix}\cos(-\beta)&\sin(-\beta)\\-\sin(-\beta)&\cos(-\beta)\end{bmatrix}=
\begin{bmatrix}0&1\\-1&0\end{bmatrix}.
\]
If we represent the base vectors of $\cl{2}{0}$ by matrices
$\hat{e}_{1}=\left[\begin{smallmatrix}1&0\\0&-1\end{smallmatrix}\right]$
and
$\hat{e}_{2}=\left[\begin{smallmatrix}0&1\\1&0\end{smallmatrix}\right]$,
then we immediately find that the expression for  generator
coincides with the bivector
$\hat{e}_{12}=\hat{e}_{1}\hat{e}_{2}=\left[\begin{smallmatrix}0&1\\-1&0\end{smallmatrix}\right]$ matrix.
We also can proceed in the opposite direction, i.~e., exponentiate
the bivector matrix $\hat{e}_{12}$. This gives the initial matrix
representation of the Lie group $\mathsf{SO}(2,\bbR)$ from the
algebra,
\[\hat{e}_{12}\rightarrow\exp{(\beta\hat{e}_{12})}=\hat{1}\cos\beta+\hat{e}_{12}\sin\beta=
\begin{bmatrix}\cos\beta&\sin\beta\\-\sin\beta&\cos\beta\end{bmatrix}\equiv
A(\beta),\] which can be rewritten in GA form as
$e^{\beta\e{12}}=\cos{\beta}+\e{12}\sin\beta$.
\end{example}
%
\begin{unit}\textit{Lie algebra}. \label{Liealgebra}
A Lie algebra $\mathfrak{g}$ is a vector space (i.~e. a linear
space with vector addition) endowed with a bilinear operator
$[\;,\;]\,:\mathfrak{g}\times\mathfrak{g}\rightarrow\mathfrak{g}$
called the Lie bracket (Lie product) which satisfies the following
axioms for $x,y,z\in\mathfrak{g}$ and $\alpha,\beta\in\bbR$:
\[\begin{split}
[\alpha x+\beta y,z]=\alpha[x,z]+\beta[y,z],&\quad\textrm{linearity}\\
[x,x]=0,& \quad\textrm{antisymmetry}\\
[x,[y,z]]+[z,[x,y]]+[y,[z,x]]=0,&\quad\textrm{Jacoby identity}
\end{split}\]
The axioms imply that $[x,y]=-[y,x]$. A set of linearly
independent vectors $x,y,z,\ldots \in\mathfrak{g}$ is called the
\textit{generators} of Lie algebra.

Since in the bivector commutator (see unit~\ref{commutatorProperties}) both the scalar and grade-$4$ parts cancel out, the commutator of bivectors always yields the other bivector. As a result the
space of bivectors $\cG_2^n$ under commutation is automatically
closed and forms the Lie algebra. Therefore, in GA the role of
algebra generators is played by $\cl{p}{q}$ bivectors, which make
$\tfrac{1}{2}n(n-1)$-dimensional vector space. Roughly speaking,
the Lie algebra generators (vectors) $\{x,y,z,\dots\}$ may be
replaced by bivectors $\{\cA,\cB,\cC,\dots\}$ of \cl{p}{q}. The
operator $[\;,\;]$ is the GA commutator
$[\cA,\cB]=\tfrac{1}{2}(\cA\cB-\cB\cA)$. Some of Lie algebras have
been realised as bivector algebras within GA, which in turn define
the so called \textit{spin group} (see unit~\ref{TypesOfGroups}).

Every algebra  always contains two {\it improper} subalgebras, ---
the algebra itself and null algebra, which contains only single
null element $\{0\}$. A Lie algebra may contain also proper Lie
algebras. A subalgebra is called \textit{Abelian} (sub)algebra if
the Lie product (commutator) of its any two elements vanishes,
$[\cA,\cB]=0$.  A subalgebra $\mathfrak{h}\subset \mathfrak{g}$ is
{\it invariant} if Lie product of any of its elements with any
element of the algebra $\mathfrak{g}$ remains in the subalgebra
$\mathfrak{h}$, i.~e. $[\cA,\cB]\in \mathfrak{h}$ for any $\cA\in
\mathfrak{g}$ and $\cB\in \mathfrak{h}$. An invariant subalgebra
in more general context is also called an {\it ideal}. The {\it
centre} of Lie algebra  is the set of elements whose Lie product
with all elements of the algebra vanishes, i.~e. the centre is the
Abelian ideal of the algebra.

A Lie algebra is {\it simple} if it  is non-Abelian and contains
no proper ideals. The algebra is {\it semi-simple} if it contains
no Abelian ideals (except $\{0\}$), i.~e., non-Abelian ideals are
allowed. A simple algebra is semi-simple as well. Semi-simple
algebra then is a direct sum of simple algebras. Semi-simple algebras
are of primary importance in physics.


The number  of generators in $\mathfrak{g}$, or, in other words
the dimension of vector space, or the number of independent
vectors required to form basis for the algebra,  is the same as
the number of group parameters in $\m{G}$. The elements of the
algebra $\mathfrak{g}$ and the group $\m{G}$ are related via
exponential and logarithmic-differential maps (see below),
respectively.
\end{unit}
%
\begin{example}\label{so2}
The matrix
$\mathfrak{a}=\bigl[\begin{smallmatrix}0&1\\-1&0\end{smallmatrix}\bigr]$ is
the generator of $\mathfrak{so}(2)$ algebra that gives rise to the
one-parameter group $\mathsf{SO}(2)$ that describes rotations in
plane. The group manifold $\mathcal{M}$ is a circle, also known as
1D sphere~$S^1$. The group parameter (the rotation angle $\beta$)
covers the entire perimeter of the circle, i.~e., $0\le\beta<
2\pi$. The general element of this group is related with
$\mathfrak{so}(2)$ generator~$\mathfrak{a}$ via exponential map,
$A(\beta)=\ee^{\beta \mathfrak{a}}=
\bigl[\begin{smallmatrix}\cos\beta&\sin\beta\\-\sin\beta&\cos\beta\end{smallmatrix}\bigr]$.
\end{example}
%
\begin{unit}\textit{Exponential and differential maps}.\label{expdiffmap}
The \textit{Lie Fundamental Theorem} says that for any finite
dimensional Lie algebra $\mathfrak{g}$, there is a unique simply
connected Lie group $\mathsf{G}$ whose Lie algebra is
$\mathfrak{g}$. The mutual correspondence is realized via
exponential and differential maps.
%~\cite{Hestenes84}.

If $\mathfrak{a}$ is one of the generators of the Lie algebra then the
element $A(\beta)$ of (simply connected) Lie group, which is
related to parameter $\beta\in\bbR$ can be obtained by the
exponential map $A(\beta)=\ee^{\beta \mathfrak{a}}$. This gives
parametrization of one parameter subgroups of the Lie group. If
instead, matrix representation of the group is given, then
respective algebra generator can easily be extracted with the
differential map $(\mathrm{d} A) A^{-1}$ (evaluated at some point
of manifold), where $A^{-1}$ is the inverse of $A$ and $\mathrm{d}
A$ is the differential of group element (matrix) with respect to
all group parameters. Typically, however, one is interesting only
in algebra element associated with single generator. Differential
and exponential maps do not account for global structure of group
manifold, i.~e. the differential map of not simply connected group
will yield the same algebra element as for a simply connected group.
The expression of differential map $\mathrm{d} (\log
A)=(\mathrm{d} A) A^{-1}$ follows from formal differentiation of
logarithm, i.~e. reverse of the above mentioned exponential map.
\end{unit}
%
\begin{example}\label{so2alg}
Let's check, that for given group element $A$, the differential
map $(\mathrm{d} A) A^{-1}$ always yields algebra (not group!)
element. Starting from the group element $A(\beta)=\ee^{\beta \mathfrak{a}}=
\bigl[\begin{smallmatrix}\cos\beta&\sin\beta\\-\sin\beta&\cos\beta\end{smallmatrix}\bigr]$
of the previous example we get $\frac{\partial
A(\beta)}{\partial\beta}A^{-1}(\beta)=\frac{\partial \ee^{\beta
\mathfrak{a}}}{\partial\beta} \ee^{-\beta \mathfrak{a}}=\mathfrak{a} \ee^{\beta \mathfrak{a}}
\ee^{-\beta \mathfrak{a}}=\mathfrak{a}$. Or repeating the same
calculations in matrix form
\[
\frac{\partial A(\beta)}{\partial\beta}A^{-1}(\beta)=
\begin{bmatrix}-\sin\beta&\cos\beta\\-\cos\beta&\sin\beta\end{bmatrix}
\begin{bmatrix}\cos(-\beta)&\sin(-\beta)\\-\sin(-\beta)&\cos(-\beta)\end{bmatrix}= \begin{bmatrix}0&1\\-1&0\end{bmatrix},
\]
which is again the sought algebra generator $\mathfrak{a}$ of
example~\ref{so2}.
\end{example}
%
\vspace{2mm}
 For the physicists it is important to remember that
the group has only single operation defined (i.~e. composition or
multiplication of group elements), while the Lie algebra has two
operations: multiplication (composition) and addition
(superposition) of elements.
%
\begin{example} 
$\mathsf{SO}(3)$ \textit{group} is the rotation group of physical
3D space. The respective $\mathfrak{so}(3)$ algebra has three
generators. In the  parameter space the group may be visualized as
a 3D solid ball. The point in the ball is characterized by a set of three
parameters $\{\theta,\varphi,R\}$, i.~e. two angles $0\le\theta<\pi$ and $0\le\varphi<2\pi$, 
and radius
of length $R\le\pi$. This ball, however, is a bit unusual. In
particular, the opposite (antipodal) points on surface of the ball
are identified. Topologically this means that two opposite points
on a diameter of the ball should be considered as a single point.
The simplest way to see this is to relate the ball (manifold in
the parameter space) with  rotations in the physical 3D space
using Rodrigues rotation formula. In \cl{3}{0} the Rodriques
formula is a transformation $\br\longrightarrow\br^{\prime}$ that
describes rotation  of arbitrary  vector $\br$ around unit vector
$\hat{\bn}=\cos\varphi \sin\theta\,\e{1}+\sin\varphi
\sin\theta\,\e{2}+ \cos\theta\,\e{3}$ by angle~$\phi$,
\[
  \br^{\prime}=\m{R}_{\hat{\bn}}(\br)=\cos\phi\,\br + (1-\cos\phi)(\br\d\hat{\bn})\hat{\bn}
  -\sin\phi(\hat{\br}\w\hat{\bn})I.
\]
Let, in addition, $\phi=R$. In the physical space the end of
vector~$\br^{\prime}$ will record a anticlockwise/clockwise circle
around $+\hat{\bn}$ and $-\hat{\bn}$, respectively,  as $\phi$ is
increased. In the ball,  the points  will move simultaneously from
the center to opposite directions along the diameter whose
direction is determined by angles $\theta$ and $\varphi$. At
$\phi=\pi$, the opposite  ends of the diameter will be reached.
Correspondingly, in the physical space the pairs of vectors
rotating in opposite directions will meet and merge into a single
vector at $\phi=\pi$. Indeed, according to the Rodrigues formula
the term with $\sin(\pi)=0$ vanishes and the remaining terms are
invariant to sign change $\hat{\bn}\rightarrow -\hat{\bn}$. This
means that two opposite points on the ball surface  should be
considered as the same point. The direct consequence of this
identification is that the $\mathsf{SO}(3)$ group  is not simply
connected. Indeed, two paths (loops), one of which closes itself
inside the ball without crossing the boundary of sphere, and the
other which goes away from one boundary point and reenters back from
the antipodal boundary point are not equivalent. The first path
can be contracted to a point (say, to the identity represented by
the center of ball), whereas the second cannot (compare
figure~\ref{fig:connectedManifold}b). This happens, because one
cannot move the opposite points of diameter on the sphere to the
center without cutting the path. The other way to say, the
opposite points cannot be contracted to a single point, because
both represent the same group element. The manifold of
$\mathsf{SO}(3)$, therefore, is connected, but not simply connected.

In order to link  group and algebra we will choose different
$\mathsf{SO}(3)$ parametrization. Namely, instead of single
rotation about one axis, we will use three orthogonal rotation
axes.  Rotation about each of the axis then can be written as one
parameter subgroups,\footnote{This will enable us to get all listed 
generator matrices from differential map expression evaluated at the same
single point. Most of widely used parameterizations, for example,
Wigner-$D$ or Wigner-$U$, use multiple rotations (by different
angles) about the same axes. In these cases in order to extract all
generators (in the form of the above given matrices) we would need to evaluate differential map at few
different points of manifold.}
\[
\begin{split}
& R_1(\alpha)=\begin{bmatrix}
1&0&0\\
0&\cos\alpha&-\sin\alpha\\
0&\sin\alpha&\cos\alpha
\end{bmatrix},\quad
R_2(\beta)=\begin{bmatrix}
\cos\beta&0&\sin\beta\\
0&1&0\\
-\sin\beta&0&\cos\beta
\end{bmatrix},\\
&R_3(\gamma)=\begin{bmatrix}
\cos\gamma&-\sin\gamma&0\\
\sin\gamma&\cos\gamma&0\\
0&0&1
\end{bmatrix}.
\end{split}
\]
Note, that despite symmetric form of matrices the angles ranges are
different: $0\le\alpha,\beta<\pi$ and $0\le\gamma< 2\pi$. This is
required in order to cover all possible 3D rotations $R(\alpha,
\beta,\gamma)=R_1(\alpha)R_2(\beta)R_3(\gamma)$ exactly once,
i.~e. to ensure one to one correspondence between the parameter space and 3D
rotations.

From these matrices the following generators are found
\[
\begin{split}
& L_1= \left(\frac{\partial R(\alpha,
0,0)}{\partial\alpha}R^{-1}(\alpha, 0,0)\right)\Big\rvert_{\alpha
=0} = \begin{bmatrix}
0&0&0\\
0&0&-1\\
0&1&0
\end{bmatrix},\\
& L_2=\left(\frac{\partial
R(0,\beta,0)}{\partial\beta}R^{-1}(0,\beta,0)\right)\Big\rvert_{\beta
=0} = \begin{bmatrix}
0&0&1\\
0&0&0\\
-1&0&0
\end{bmatrix},\\
& L_3=\left(\frac{\partial
R(0,\gamma,0)}{\partial\gamma}R^{-1}(0,\gamma,0)\right)\Big\rvert_{\gamma
=0}= \begin{bmatrix}
0&-1&0\\
1&0&0\\
0&0&0
\end{bmatrix}.
\end{split}
\]
Note, that differential map was calculated at a single point in
the parameter space, $(\alpha,\beta,\gamma)=(0,0,0)$, i.~e. at the
point which corresponds to the identity element on the manifold
represented by the center point of the parameter ball. It is rarely
noted, though, that the resulting generator matrices are centered
at zero matrix. The 3D~hyperplane of algebra generators (which is
a bivector plane and should not to be confused with the parameter
space itself), of course, is parallel to the tangent plane of
group manifold at the point of identity. It differs, however, from
the later by translation.

The chosen group parametrization was in advance adjusted to get
the well-known matrix form of generators, which satisfy standard
commutation relations
\[
[L_1,L_2]=L_3,\quad[L_2,L_3]=L_1,\quad[L_3,L_1]=L_2\, .
\]
These commutation relations exactly match the commutators of
angular momentum operators in Euclidean 3D space.
\end{example}

%
\begin{example}\textit{$\mathsf{SU}(2)$ group} \label{SU2group}locally is isomorphic
to $\mathsf{SO}(3)$ which describes 3D~rotations too. Globally
both groups differ, thought their algebras coinside. The
manifold of $\mathsf{SU}(2)$ is simply connected and can be
represented by three dimensional sphere $S^3$ in 4D Euclidean
space $\bbR^4$. $\mathsf{SU}(2)$ group covers $\mathsf{SO}(3)$
exactly twice, i.~e. for each rotation in $\mathsf{SO}(3)$ there
exist exactly two different rotations in $\mathsf{SU}(2)$.
Therefore $\mathsf{SU}(2)$ is called \textit{double cover} of
$\mathsf{SO}(3)$. By definition a covering space is called a
\textit{universal covering space} if it is simply connected.
Because $S^3$ is simply connected, $\mathsf{SU}(2)$ also is called
a universal cover of $\mathsf{SO}(3)$.


The role of $\mathsf{SO}(3)$ Rodrigues matrices in
$\mathsf{SU}(2)$ is played by so called Wigner~$U$ matrices, which
describe rotations in exactly the same way, i.~e., by specifying a
single rotation axis and the rotation angle about this axis. The
other widely used $\mathsf{SU}(2)$ parametrization scheme is known
as Wigner~$D$ matrices. They parameterize rotations using two
fixed rotation axes. In particularly these matrices are
constructed by first rotating arbitrary vector about one axis,
then about the second axis and lastly about the first axis again, but by different angle.
Both mentioned parameterizations are widely used in theory of
angular momentum.

In order to write $\mathfrak{su}$(2) generators in terms of Pauli matrices,
we again shall start from three matrices each depending on just one parameter.
Each parameter represents rotation about different rotation axis by corresponding
angle $\alpha$, $\beta$ or $\gamma$~\cite{Sattinger86}
\[
U_1=\begin{bmatrix}\cos\tfrac{\alpha}{2}&\ii\sin\tfrac{\alpha}{2}\\\ii\sin\tfrac{\alpha}{2}&\cos\tfrac{\alpha}{2}\end{bmatrix},\
U_2=\begin{bmatrix}\cos\tfrac{\beta}{2}&-\sin\tfrac{\beta}{2}\\\cos\tfrac{\beta}{2}&\sin\tfrac{\beta}{2}\end{bmatrix},\
U_3=\begin{bmatrix}\ee^{\ii\gamma/2}&0\\0&\ee^{-\ii\gamma/2}\end{bmatrix}\,.
\]
%The product of matrices makes general $\m{SU}(2)$ element once
%parameter ranges are chosen in the interval $0\le\alpha,\beta\le
%2\pi$ and $0\le\gamma < 4\pi$ \red{butina patikrinti kampu kitimo
%rezius}.
The product of matrices makes  up the general $\m{SU}(2)$ element
once parameter ranges are chosen in the interval $0\le\alpha<
2\pi$, $0\le\beta\le \pi$ and $0\le\gamma < 2\pi$ (note assymetric ranges for parameters $\alpha$ and $\beta$). 
% parameters are given in  Biedenharn-Louck, rusiskas vertimas, psl.35, 
 The generators of
$\mathfrak{su}$(2) algebra then can be obtained from the above
matrices by logarithmic-differentiation map calculated at single
point $(\alpha,\beta,\gamma)=(0,0,0)$ of parameter space,
\[
J_1=\begin{bmatrix}0&\frac{\ii}{2}\\\frac{\ii}{2}&0\end{bmatrix},\quad
J_2=\begin{bmatrix}0&-\frac{1}{2}\\\frac{1}{2}&0\end{bmatrix},\quad
J_3=\begin{bmatrix}\frac{\ii}{2}&0\\0&-\frac{\ii}{2}\end{bmatrix}\,.
\]
These generators can be expressed through Pauli matrices,
$J_1=\tfrac{\ii}{2}\hat{\sigma}_x$,
$J_2=-\tfrac{\ii}{2}\hat{\sigma}_y$,
$J_3=\tfrac{\ii}{2}\hat{\sigma}_z$, and satisfy the same
commutation relations as angular momentum operators
\[[J_1,J_2]=J_3,\quad[J_2,J_3]=J_1,\quad[J_3,J_1]=J_2\,. \]
From GA point of view it is preferable to express the generators
as products of Pauli matrices
\[
J_1=\tfrac{1}{2}\hat{\sigma}_y\hat{\sigma}_z,\quad
J_2=\tfrac{1}{2}\hat{\sigma}_x\hat{\sigma}_z,\quad
J_3=\tfrac{1}{2}\hat{\sigma}_x\hat{\sigma}_y \,.
\]
In this form the operators are represented by bivectors of
\cl{3}{0} algebra, which are closed under commutation relation
(see unit~\ref{commutatorProperties}). It follows that spin group
of \cl{3}{0} is (locally) isomorphic to $\mathsf{SU}(2)$ (and
$\mathsf{SO}(3)$). The fact that the generators of
$\mathfrak{su}(2)$ can be expressed through Pauli matrices as well
as their products (bivectors) is an exceptional case due to a
lucky coincidence of number of vectors and bivectors in \cl{3}{0}.
For other GA algebras the number of vectors and bivectors differ.
It follows that names (which specify the dimension) of Lie and spin algebras will
also be different.
\end{example}
%
\begin{unit}\textit{Structure constants, Killing form, adjoint and Casimir operators.
Relations between GA bivectors and Lie algebra generators.}
\label{StructureConstants}

The structure constants of infinitesimal generators of Lie algebra
play the key role in the Lie group theory, since they allow an
easy calculation of Killing form and Casimir operators which are
the invariants of the group. The Killing form, in turn, is
important in finding root system of the group, which is essential
for group classification. Since the Lie algebra is linear (with respect to generators)
and closed under Lie bracket, the commutators of infinitesimal
generators $\cX_i$ always can be expressed through the  same
generators:
\[[\cX_i,\cX_j]=\sum_k c_{ij}^k\cX_k,\]
where the numerical coefficients $c_{ij}^k$ are called the
\textit{structure constants}. They  determine the structure of
algebra completely.  The superscript $k$ in $c_{ij}^k$ denotes the index in
the expansion of the product. The (numerical) values of
structure constants dependent on selected basis, i.~e. their
values will change once the vector base of algebra is changed. The
antisymmetry of commutator and  Jacobi identity
(see~\ref{Liealgebra}) ensure, respectively, the following relations between
structure constants
\[
\begin{split}
  &c_{ij}^k=-c_{ji}^k,\quad \textrm{and}\quad 
  & \sum_{l=1}^d c_{il}^m c_{jk}^l+c_{jl}^m c_{ki}^l+c_{kl}^m c_{ij}^l
  =0\,.
\end{split}
\]
Structure constants are closely related to linear \textit{adjoint
operator} denoted as $\ad(\cA)$, on the algebra
$\mathfrak{g}\ni\cA$, which if applied to the element $\cX$ is
defined through the commutator as $\ad(\cA)\cX = [\cA,\cX]$ for
every algebra element $\cX\in\mathfrak{g}$. The set of all these
operators $\ad(\cA)$ makes up a Lie algebra of the same dimension
$d$ as the original algebra $\mathfrak{g}$. The consecutive action
of two adjoint operators  $\ad(\cA)$ and $\ad(\cB)$ is understood
as $\ad(\cA)\ad(\cB)\cX=\ad(\cA)[\cB,\cX]=[\cA,[\cB,\cX]]$.
The adjoint commutator for arbitrary $\cX$ is defined as
\[
\begin{split}
  \ad([\cA,\cB])\cX& = [ [\cA,\cB],\cX]= [\cA ,[\cB,\cX]]+ [\cB ,[\cX,\cA]]\\
           &=  \ad(\cA)  \ad(\cB)\cX -\ad(\cB)  \ad(\cA) \cX= [\ad(\cA),\ad(\cB)]
           \cX\,,
\end{split}
\]
where Jacobi identity was used in the third step. Because the above 
equation is valid for all $\cX$,  i.~e. $\ad([\cA,\cB])=
[\ad(\cA),\ad(\cB)]$ we obtained a natural mapping
$\cA\to\ad(\cA)$. This mapping, therefore, is a representation of
the algebra $\mathfrak{g}$, known as {\it adjoint representation},
which plays an important role in many group theoretical
constructs. In particular, the matrix elements of adjoint
representation of algebra generators coincide with the structure constants
of the algebra, i.~e., 
\[(\ad(\cX_k))_{i\,j}=c_{kj}^i.\] 
This can be seen after
comparing results of the following calculations:
$\ad(\cX_i)\cX_j=\sum_k (\ad(\cX_i))_{kj} \cX_k$ and
$\ad(\cX_i)\cX_j= [\cX_i,\cX_j]= \sum_k c_{ij}^k\cX_k$. Here the
index sequence in the first expression ensures that the result
will be a matrix with row and column indices $i$ and $j$,
respectively.

The adjoint representation matrices constitute the algebra
representation with the same structure constants, i.~e.
\[
  ([\ad(\cX_k),\ad(\cX_l)])_{ij}=\sum_k c_{kl}^m\, \ad(\cX_m)_{ij}\,.
\]
The adjoint representation allows to define invariant bilinear
form called the {\it Killing form} (we will see that in GA terms it is a metric of
bivector subspace $\cG_2^m$), which plays an important role in
many areas of group theory. In particularly, the Killing form is
defined as a trace of simple product of any two elements of
algebra $\mathfrak{g}$,
\[
  g_{\cA\cB}=\Tr(\ad(\cA)\ad(\cB ))\,.
\]
Killing metric was denoted by the same letter as GA metric tensor
$g_{ik}=(\e{i}\d\e{k})$ of the GA vector subspace $\cG_1^m$.
The range of indices of Killing metric, however goes over all GA
bivectors (not GA vectors). So, simply speaking, the Killing form
is a metric of bivector space (which at the same time also is a
metric of generator space, i.~e. vector space of Lie algebra). In
order to avoid confusion, in GA we will denote the Killing form by
different letter, i.~e., as
 \[K_{\{ij\}\{kl\}}=(\e{ij}\d\e{kl}),\]
or $K_{mn}$ when the bivectors are enumerated in a canonical order.
In fact this formula can be taken as a definition of \textit{Killing matrix in
GA}.

Note, that in case of Lie algebra only the Lie bracket operation
between two elements is defined. However, though in the Lie
algebra a simple product of elements is not defined, we can define such a
product in any representation of the algebra in terms of operators
on vector space or in terms of matrices. So, the simple product of
two elements can be defined either as an consecutive action of two
operators (like in the trace formula above) or as a matrix product
of two matrix representations of Lie algebra elements. Because
under trace operation terms can be permuted cyclically, the
Killing form is symmetric, $g_{\cA\cB}=g_{\cB\cA}$. From cyclic
permutation of three elements the property
$g_{[\cA,\cB]\cC}=g_{\cA[\cB,\cC]}$ also can be easily derived.

The Killing form of base vectors can be completely rewritten in
terms of structure constants only,
\[
  g_{ij}=\Tr(\ad(\cX_i)\ad(\cX_j))= \delta\sum_{k,l=1}^d c_{il}^k
  c_{jk}^l\,,
\]
where $\delta$ is the normalization (Dynkin) constant. When Killing
form of algebra is calculated, we can easily determine wheter the given Lie algebra is semi-simple or not. 
There exists a well known {\it Cartan criterion}, which says that Lie
algebra is semi-simple if and only if its Killing metric is
nondegenerate (the matrix determinant is not zero).

For a semi-simple algebra the Killing metric can be inverted
$(g_{ij})^{-1}=g^{ij}$. Then, in turn, fully antisymmetric
structure constants can be defined as
\[
c_{ijk} =\sum_l g_{il}  c_{jk}^l\,.
\]
Total antisymmetry of $c_{ijk}$ in its all three indices can be
verified by expressing $g_{ij}$ in terms of structure constants
and then applying Jacobi identity to the two last structure
constants. For semi-simple algebra the above formula has
inverse $c_{jk}^i =\sum_l g^{il}  c_{ljk}$.

With the help of inverse Killing metric we can define invariant group
operators. The simplest one is the {\it Casimir
operator} $C_2$ (the index shows a number of algebra generators),
\[
  C_2 =\sum_{i,j=1}^d g^{ij}  \cX_i \cX_j,\qquad\cX_i\in\mathfrak{g}.
\]
The Casimir operator commutes with all generators $\cX_i$ and,
thus with all algebra elements. Higher order invariant operators
are defined as
\[
  C_n =\sum c_{i_1j_1}^{j_2} c_{i_2 j_2}^{j_3} \cdots c_{i_n j_n}^{j_1} g^{i_1l_1} g^{i_2l_2} \cdots g^{i_nl_n} \cX_{l_1} \cX_{l_2} \cX_{l_n}\,.
\]
This formula reduces to Casimir operator for $n=d=2$. As in the
case of Casimir operator, the invariant operators $C_n$ commute
with all elements of the semi-simple Lie algebra as well.

A real Lie algebra is defined to be {\it compact} (for geometric
meaning see unit~\ref{Definitionofgroup}) if its Killing form is
negative definite. A compact algebra is necessarily semi-simple.
Complex Lie algebras are not compact~\cite{Kirson}.
% Michael W. Kirson "Introductory Algebra for Physicists"
%@BOOKLET{Kirson,
%   author="M.~W. Kirson",
%   title="Introductary Algebra for Physicists",
%   publisher="",
%   address="",
%   year="",
%   note=""
%   }

%https://webhome.weizmann.ac.il/home/fnkirson/Alg17/course.html
%\red{Dabar supazindinus  su tradicine Lie teorija reikia surasyti
%visus aprasytus atitikmenis per GA savokas}
\end{unit}
%
\begin{unit}\textit{Bivector subspace as Lie algebra vector space}.
As it was noted in unit~\ref{Liealgebra}, the bivector subspace
$\cG_2^n$ is  closed under commutator product, i.~e.,
$[\cA,\cB]\in\cG_2^n$ for all $\cA,\cB\in\cG_2^n$. Therefore the
bivectors realize some finite-dimensional real or complex Lie
algebra.  

Given any algebra $\cl{p}{q}$ there exists a vector 
basis $\{\bff_1,\dots,\bff_n\}$ in $\cV^n$ where $n=p+q$, such that
any general bivector $\cB\in\cG_2^n$ can be written in the following form 
\[
\cB=\beta_{12}\bff_1\w\bff_2+\beta_{34}\bff_3\w\bff_4+\dots+\beta_{n-1,n}\bff_{n-1}\w\bff_{n},
\]
There are $\bigl(\begin{smallmatrix}n\\2\end{smallmatrix}\bigr)=n(n-1)/2$ basis bivectors $\bff_i\w\bff_j$. %\cite{Lundholm09}~p.68
The structure constants for bivector algebra of $\cl{p}{q}$ 
then can easily be extracted from the commutators of bivectors
\[
\tfrac{1}{2}[\e{ij},\e{kl}]=(\e{i}\d\e{k})\e{lj}+(\e{i}\d\e{l})\e{jk}+
(\e{j}\d\e{k})\e{il}+(\e{j}\d\e{l})\e{ki}\,.
\]
where $(\e{i}\d\e{k})=g_{ik}$ are GA vector 
$\cV^n$ space metric tensor components (see~\ref{reciprocalFrame}). In particular, in 
orthonormal bivector base the structure constants are
\begin{alignat*}{4}
&c_{\{ij\}\{ik\}}^{\{jk\}}=&-c_{\{ij\}\{ik\}}^{\{kj\}}=&-\e{i}^2,&\qquad
    c_{\{ij\}\{ki\}}^{\{jk\}}=&-c_{\{ij\}\{ki\}}^{\{kj\}}=&\e{i}^2,\\
&c_{\{ji\}\{ik\}}^{\{jk\}}=&-c_{\{ji\}\{ik\}}^{\{kj\}}=&\e{i}^2,&
    c_{\{ji\}\{ki\}}^{\{jk\}}=&-c_{\{ji\}\{ki\}}^{\{kj\}}=&-\e{i}^2.
   \end{alignat*}
Here the indices in the curly brackets $\{ij\}$ enumerate base
bivectors $\e{ij}$. Thus the structure constants can be expressed
through \cl{p}{q} algebra basis vectors.

As was already noted in the previous unit, the Killing form gives
the metric of the Lie algebra, or equivalently the metric of
$\cG_2^n$ subspace. In GA the Killing form  can be
expressed through the bivector inner product~\cite{Lundholm09},% p.70
\[K_{ij}\equiv K(\cA_i,\cB_j)=\cA_i\d\cB_j=\langle\cA_i\cB_j\rangle_0,\quad\cA_i,\cB_j\in\cG_2^n. \]
All possible inner products of bivectors compose the Killing matrix.
If bivectors are expressed through GA basis vectors $\e{i}$ then
the Killing matrix is diagonal, since in this case $K_{ij}=0$ if
$i\ne j$ and $K_{ii}=\pm 1$.

\red{Once Killing metrix is calculated, the Casimir operator $C_2$ in GA can be written as 
\[
  C_2= \sum_{i,j} K_{ij}^{-1} \cA_i \cA_j\,,
\]
where $K_{ij}^{-1}$ is the inverse of Killing metric and bivectors $\cA_i$ constitute base of $\cG_2^n$ space. In orthormal base $K_{ij}^{-1}=K_{ij}$.
The second order Casimir operator $C_2$ in GA has an interesting interpretation: it simply counts number of bivectors in GA. Note that similar construction in vector subspace $\sum_{i,j} g_{ij}^{-1} a_i a_j$ also yields GA invariant, i.~e. the number of vectors. The higher Casimir operators, when exists, can be interpreted in a similar way, i.~e. as a number of trivectors, four-vectors, etc... }
\end{unit}
% galima veliau pratest C_3. Turetu buti susijusios su invariantu trivektoriu erdveje (reiktu dar atidziai patikrint). Faktiskai geometrine grupiu teorija tam tikra prasme!
%
\begin{example}%~\cite{Doran93}.
In this example we will explicitly check that $\mathfrak{su}(2)$
Lie algebra is closely related to $\cl{3}{0}$ bivector algebra.
 It then will follow that the Lie
group $\mathsf{SU}(2,\bbC)$ and spin group $\Spin{3}{0}$ (for
definition see next Section) are isomorphic,
$\Spin{3}{0}\cong\mathsf{SU}(2,\bbC)$.

The elements of Lie algebra satisfy the axioms
in~\ref{Liealgebra}. In GA, the Lie algebra can be generated by
commutators of matrices that represent GA bivectors. Therefore, we
first have to check that the number of generators in both algebras
is the same. Indeed, looking at item~\ref{classgroup} we find
that the number of generators in the $\mathsf{SU}(2,\bbC)$ group 
is $n_{\mathrm{G}}=n^2-1=3$, which is equal to the number of
bivectors  $\tfrac{1}{2}n(n-1)=3$ of \cl{3}{0}.

Now let's construct commutator tables of both algebras. The
standard matrices that represent $\mathfrak{su}(2)$ algebra may be
expressed via Pauli matrices as $-\ii\hat{\sigma}_1$,
$-\ii\hat{\sigma}_2$, and $\ii\hat{\sigma}_3$. The
table~\ref{bivectorAlgebraToLie} contains explicit matrix
representations for other Lie algebras. Calculating all
commutators of Pauli matrices and commutators of bivectors of
$\cl{3}{0}$ yield the following commutation tables
\[\begin{array}{c|ccc}
\tfrac{1}{2}[\d,\d] &\e{12}&\e{23}&\e{31}\\
\hline
 \e{12}&0 &-\e{31} &\e{23} \\
 \e{23}&\e{31} &0 &-\e{12} \\
 \e{31}&-\e{23}&\e{12} &0
\end{array}\qquad
\begin{array}{c|ccc}
\tfrac{1}{2}[\d,\d] &-\ii\hat{\sigma}_1&-\ii\hat{\sigma}_2&\ii\hat{\sigma}_3\\
\hline
 -\ii\hat{\sigma}_1&0 &-\ii\hat{\sigma}_3 &-\ii\hat{\sigma}_2 \\
 -\ii\hat{\sigma}_2&\ii\hat{\sigma}_3 &0 &\ii\hat{\sigma}_1 \\
\phantom{-}\ii\hat{\sigma}_3&\ii\hat{\sigma}_2&-\ii\hat{\sigma}_1 &0
\end{array}
\]
From the tables one can read the structure constants \red{(coefficients in front of bivectors)}. It
is easy to see that both algebras are isomorphic. The explicit
isomorphism is given by
\[
\e{12}\Leftrightarrow
\begin{bmatrix}0&-\ii\\-\ii&0\end{bmatrix}=-\ii\hat{\sigma}_1,\quad
\e{23}\Leftrightarrow\begin{bmatrix}0&-1\\1&0\end{bmatrix}=-\ii\hat{\sigma}_2,\quad
\e{31}\Leftrightarrow\begin{bmatrix}\ii&0\\0&-\ii\end{bmatrix}=\ii\hat{\sigma}_3\,.
\]
To obtain the Lie group in GA form or matrix form it is enough to
take an exponential map of obtained matrices multiplied by scalar
parameter (see the item~\ref{Liealgebra}),
\[\begin{split}
&\ee^{\alpha\e{12}/2}=\cos\tfrac{\alpha}{2}+\e{12}\sin\tfrac{\alpha}{2}\Leftrightarrow\ee^{-\alpha\ii\hat{\sigma}_1/2}=
\begin{bmatrix}\cos\tfrac{\alpha}{2}&-\ii\sin\tfrac{\alpha}{2}\\-\ii\sin\tfrac{\alpha}{2}&\cos\tfrac{\alpha}{2}\end{bmatrix},\\
&\ee^{\beta\e{23}/2}=\cos\tfrac{\beta}{2}+\e{23}\sin\tfrac{\beta}{2}\Leftrightarrow\ee^{-\beta\ii\hat{\sigma}_2/2}=
\begin{bmatrix}\cos\tfrac{\beta}{2}&-\sin\tfrac{\beta}{2}\\\sin\tfrac{\beta}{2}&\cos\tfrac{\beta}{2}\end{bmatrix},\\
&\ee^{\gamma\e{31}/2}=\cos\tfrac{\gamma}{2}+\e{31}\sin\tfrac{\gamma}{2}\Leftrightarrow\ee^{\gamma\ii\hat{\sigma}_3/2}=
\begin{bmatrix}\ee^{\ii\gamma/2}&0\\0&\ee^{-\ii\gamma/2}\end{bmatrix}.
\end{split}\]
\red{The sign of the parameter $\alpha$ here is oposite to that of Example~\ref{SU2group} and is caused by different sign of the first generator, i.~e. $\ii\hat{\sigma}_1 =-J_1=-\ii\hat{\sigma}_x$.}
The determinant of matrices
is equal to $1$ and the matrices satisfy
$\hat{U}_i\hat{U}_i^{-1}=\hat{U}_i\hat{U}_i^{\dagger}=\hat{1}$. So
they belong to  $\mathsf{SU}(2)$  and describe rotations around
three orthogonal axes (or in terms of GA, rotations in
perpendicular bivector planes). In conclusion, we have shown that
starting from GA vector space $\cV^n$ the following chain
$\cV^n\rightarrow\cG_2^n\rightarrow\mathfrak{su}(n)\rightarrow\mathsf{SU}(n)$,
has generated Lie group representation.
\end{example}
%
\begin{example}\label{SL2R} The special linear group $\m{SL}(2,\bbR)$ is
connected, non-compact Lie group with three real parameters.
Generators of $\mathfrak{sl}(2,\bbR)$ algebra are traceless
matrices usually written as
\[H=\begin{bmatrix}1&0\\0&-1\end{bmatrix},\quad
X=\begin{bmatrix}0&1\\0&0\end{bmatrix},\quad
Y=\begin{bmatrix}0&0\\1&0\end{bmatrix}, \] which
satisfy the commutation relations $[H,X]=2X$, $[H,Y]=-2Y$, and
$[X,Y]=H$. Assuming that generators are enumerated in the listed order it follows that nonzero structure constants of $\mathfrak{sl}(2,\bbR)$ algebra are: $c_{12}^2=2$, $c_{13}^3=-2$, $c_{23}^1=1$.
The constants satisfy $\sum_{l=1}^3 c_{il}^m c_{jk}^l+c_{jl}^m
c_{ki}^l+c_{kl}^m c_{ij}^l=0$, where terms with $i=1,2,3$ nullify
seperately.

To connect $\mathfrak{sl}(2,\bbR)$ algebra with bivectors of
\cl{2}{1} we at first make a linear combinations which makes the
above basis orthonormal with respect to Killing form (a symmetric
bilinear form on the Lie algebra) defined as $\Tr (\ad(a) \ad(b))$,
where $\ad(a)$ denotes the adjoint representation of the Lie
algebra (see unit~\ref{StructureConstants}),
\[x=\frac{1}{\sqrt{2}}\begin{bmatrix}0&1\\1&0\end{bmatrix},\quad
y=\frac{1}{\sqrt{2}}\begin{bmatrix}0&-1\\1&0\end{bmatrix},\quad
z=\frac{1}{\sqrt{2}}\begin{bmatrix}1&0\\0&-1\end{bmatrix}.
\] 
New generator base yields different structure constants: $c_{12}^3=1$,
$c_{23}^1=1$, $c_{31}^2=-1$. With these constants the Killing
matrix is found to be $g_{ij}=\delta\sum_{k,l=1}^3 c_{il}^k
c_{jk}^l=\delta\,\textrm{diag}(2,-2,2)$, or after normalization by Dynkin
constant the metric becomes $g_{ij}=\textrm{diag}(1,-1,1)$.

 It is easy to check that commutation table constructed from new
$\mathfrak{sl}(2,\bbR)$ generators coincides with a similar table
for \cl{2}{1} bivectors if one takes
\[x\rightarrow -\e{23},\quad y\rightarrow -\e{13},\quad z\rightarrow\e{12}.\]
These bivectors give the same Killing matrix, $g_{ij}=$
$g_{ij}=\textrm{diag}(\e{31}^2,\e{12}^2,\e{23}^2)=\textrm{diag}(1,-1,1)$.
\end{example}
%
\begin{example}\label{SL2C} In the
relativistic space-time, 4D rotations (boosts  and space
rotations)  are described by $\m{SL}(2,\bbC)$ group. This group is
 double cover of $4\times 4$ group of real matrices usually given
in textbooks under the name of Lorentz transformation $\mathsf{SO}(1,3,\bbR)$. The relationship between relativistic
group $\mathsf{SO}(1,3,\bbR)$ and its double covering $\m{SL}(2,\bbC)$ is
very similar to relationship between quantum groups
$\mathsf{SO}(3)$ and $\mathsf{SU}(2)$. Because $\m{SL}(2,\bbC)$ is
simply connected the double cover is also universal. Six
generators of respective $\mathfrak{sl}(2,\bbC)$ Lie algebra are
\[
\begin{split}
&k_1=\begin{bmatrix}0&1\\1&0\end{bmatrix},\quad
k_2=\begin{bmatrix}0&-\ii\\\ii&0\end{bmatrix},\quad
k_3=\begin{bmatrix}-1&0\\0&1\end{bmatrix},\\
&j_1=\begin{bmatrix}0&\ii\\\ii&0\end{bmatrix},\quad
j_2=\begin{bmatrix}0&1\\-1&0\end{bmatrix},\quad
j_3=\begin{bmatrix}-\ii&0\\0&\ii\end{bmatrix}.
\end{split}
\]
They form the so called $\bbR$-basis. The name implies that
the above matrices (base vectors) can be multiplied by real
parameters (coefficients) only. Alternatively, it is possible to
take only three matrices $j_1, j_2$ and $j_3$ as a base vectors,
which form so called $\bbC$-basis of the algebra, and then
multiply by complex parameters to describe the same algebra and
group. We will not pursue the later line further, because our primary
interest is $\mathfrak{sl}(2,\bbC)$ relation to real Clifford
algebras.

The squares of $k$ and $j$ generators differ by sign,
$k_1^2=k_2^2=k_3^2=1$ and $j_1^2=j_2^2=j_3^2=-1$, which are
characteristic of time-like (boosts) and space-like rotations,
respectively. Indeed, the exponentiation of matrices multiplied by
six real scalar parameters yields matrices, elements of which
contains hyperbolic and trigonometric functions that realize boost
and rotation transformations. It can be shown that commutation
table of generators $k_i$ and $j_i$ exactly matches the commutator
table of bivectors of \cl{1}{3} if we take $k_1\rightarrow\e{12}$,
$k_2\rightarrow\e{13}$, $k_3\rightarrow\e{14}$,
$j_1\rightarrow\e{34}$, $j_2\rightarrow\e{42}$,
$j_3\rightarrow\e{23}$. This means that \cl{1}{3} algebra
bivectors are isomorphic to $\mathfrak{sl}(2,\bbC)$ algebra
$\bbR$-basis vectors. We also note, that the generator set $j_1,
j_2$ and $j_3$ over the field of real coefficients make
$\mathfrak{su}(2)$ subalgebra.

In physics \cl{1}{3} algebra is constructed from 4-dimensional
Dirac matrices, so that bivectors of this algebra are represented
by larger, $4\times 4$ complex matrices. It can be shown that
these matrices are reducible in the Lie algebra context, i.~e., by
selecting a suitable unitary transformation they can be
transformed into a block-diagonal form where the upper and lower
blocks are isomorphic to $\mathfrak{sl}(2,\bbC)$ matrices.
\end{example}
%
\begin{example}
Relativistic transformations are described by bivectors $\e{i}\e{j}$ 
of \cl{1}{3} and \cl{3}{1} algebras. 
It can be easily verified that bivector commutation tables of both algebras are isomorphics and, consequently, both algebras share the same structure constants.
It then follows that exponentiation of bivectors $\ee^{\alpha_{ij}\e{i}\e{j}}$
of the both algebras represents the same group 
$\mathsf{Sp}(2,\bbC)$ (which is also isomorphics to
$\mathsf{SL}(2,\bbC)$). This group is called the Lorentz group, and is
characterized by six real parameters (angles $\alpha_{ij}$).

From $8$-periodicity table~\ref{8foldperiod} it also follows that
the basis vectors of \cl{1}{3} can be represented by quaternionic
$\bbH(2)$ rather than by complex matrices, i.~e. quaternionic representation of $\mathsf{SL}(1,3)$ group can be constructed.
\end{example}
%
\begin{unit}\textit{Equivalence relations between GA bivectors and Lie
  algebra matrix generators for low dimensional GA's}\label{bivectorAlgebraToLie}
\[\hspace*{-1.7cm}
\begin{array}{cc|l}
\textrm{GA}&\textrm{Lie algebra}&\textrm{Equivalence relations}\\
\hline\hline
%%%%%%%%%%%%%%%%%%%%%%%%%%%%%%%%%%%%%%%%%%%%%%%%%
% Ivan Todorov "Clifford algebras and spinors" arXiv: 1106.3197v2 [math-ph]
% provides large values. His table gives too large algebras: unreliable source
\cl{3}{0}&\mathfrak{su}(2,\bbC)&\e{12}\Leftrightarrow -\ii\sig{z},\,\e{13}\Leftrightarrow -\ii\sig{y},\,\e{23}\Leftrightarrow -\ii\sig{x}\\[3pt]
\cl{2}{1}&\mathfrak{sl}(2,\bbR)&\e{12}\Leftrightarrow \ii\sig{y},\,\e{13}\Leftrightarrow \sig{z},\,\e{23}\Leftrightarrow \sig{x}\\[3pt]
\cl{1}{2}&\mathfrak{sl}(2,\bbR)&\e{12}\Leftrightarrow -\sig{z},\,\e{13}\Leftrightarrow -\sig{x},\,\e{23}\Leftrightarrow -\ii\sig{y}\\[3pt]
\cl{0}{3}&\mathfrak{su}(2,\bbC)&\e{12}\Leftrightarrow \ii\sig{z},\,\e{13}\Leftrightarrow \ii\sig{y},\,\e{23}\Leftrightarrow \ii\sig{x}\\[3pt]
\hline
%%%%%%%%%%%%%%%%%%%%%%%%%%%%%%%%%%%%%%%%%%%%%%%%%
\cl{4}{0}& {^2\mathfrak{su}}(2,\bbC)&\e{ij}\Leftrightarrow -u_k,\textrm{\quad for the same $i,j,k$ values as for \cl{0}{4}}\\[3pt]
\cl{3}{1}&\mathfrak{sl}(2,\bbC)&\e{12}\Leftrightarrow\ii\sig{y},\,\e{13}\Leftrightarrow\ii\sig{x},\,\e{14}\Leftrightarrow\sig{z},\\
&  &\e{23}\Leftrightarrow-\ii\sig{z},\,\e{24}\Leftrightarrow\sig{x},\,\e{34}\Leftrightarrow-\sig{y}\\[3pt]
\cl{2}{2}& {^2\mathfrak{sp}}(2,\bbR)&\e{12}\Leftrightarrow p_1,\,\e{13}\Leftrightarrow p_2,\,\e{14}\Leftrightarrow p_3,\\
&  &\e{23}\Leftrightarrow p_4,\,\e{24}\Leftrightarrow p_5,\,\e{34}\Leftrightarrow p_6\\[3pt]
%
\cl{1}{3}&\mathfrak{sl}(2,\bbC)&\e{12}\Leftrightarrow-\sig{z},\,\e{13}\Leftrightarrow-\sig{x},\,\e{14}\Leftrightarrow\sig{y},\\
&  &\e{23}\Leftrightarrow-\ii\sig{y},\,\e{24}\Leftrightarrow-\ii\sig{x},\,\e{34}\Leftrightarrow\ii\sig{z}\\[3pt]
%
\cl{0}{4}& {^2\mathfrak{su}}(2,\bbC)&\e{12}\Leftrightarrow u_1,\,\e{13}\Leftrightarrow u_2,\,\e{14}\Leftrightarrow u_3,\\
& &\e{23}\Leftrightarrow u_4,\,\e{24}\Leftrightarrow u_5,\,\e{34}\Leftrightarrow u_6\\[3pt]
%%%%%%%%%%%%%%%%%%%%%%%%%%%%%%%%%%%%%%%%%%%%%%%%%
\hline
\cl{3}{2}& {\mathfrak{sp}}(4,\bbR)&\e{12}\Leftrightarrow P_1,\,\e{13}\Leftrightarrow P_2,\,\e{14}\Leftrightarrow P_3,\,\e{15}\Leftrightarrow P_4,\,\e{23}\Leftrightarrow P_5,\\
        & &\e{24}\Leftrightarrow P_6,\,\e{25}\Leftrightarrow P_7,\,\e{34}\Leftrightarrow P_8,\,\e{35}\Leftrightarrow P_9,\,\e{45}\Leftrightarrow P_{10}\\[3pt]

\cl{2}{3}& {\mathfrak{sp}}(4,\bbR)&\e{12}\Leftrightarrow -P_{10},\,\e{13}\Leftrightarrow -P_3,\,\e{14}\Leftrightarrow -P_6,\,\e{15}\Leftrightarrow -P_8,\,\e{23}\Leftrightarrow -P_4,\\
        & &\e{24}\Leftrightarrow -P_7,\,\e{25}\Leftrightarrow -P_9,\,\e{34}\Leftrightarrow -P_1,\,\e{35}\Leftrightarrow -P_2,\,\e{45}\Leftrightarrow -P_{5}\\[3pt]
%%%%%%%%%%%%%%%%%%%%%%%%%%%%%%%%%%%%%%%%%%%%%%%%%%%
\hline
\cl{6}{0}& {\mathfrak{su}}(4,\bbC)&\e{ij}\Leftrightarrow -U_k,\textrm{\quad for the same $i,j,k$ values as for \cl{0}{6}}\\[3pt]
%
\cl{4}{2}& {\mathfrak{su}}(2,2,\bbC)&\e{12}\Leftrightarrow S_1,\e{13}\Leftrightarrow S_2,\e{14}\Leftrightarrow S_3,\e{15}\Leftrightarrow S_4,\e{16}\Leftrightarrow S_5,\\
& &\e{23}\Leftrightarrow S_6,\e{24}\Leftrightarrow S_7,\e{25}\Leftrightarrow S_8,\e{26}\Leftrightarrow S_9,\e{34}\Leftrightarrow S_{10},\\
& &\e{35}\Leftrightarrow S_{11},\e{36}\Leftrightarrow S_{12},\e{45}\Leftrightarrow S_{13},\e{46}\Leftrightarrow S_{14},\e{56}\Leftrightarrow S_{15}\\[3pt]
%
\cl{3}{3}& {\mathfrak{sl}}(4,\bbR)&\e{12}\Leftrightarrow L_1,\e{13}\Leftrightarrow L_2,\e{14}\Leftrightarrow L_3,\e{15}\Leftrightarrow L_4,\e{16}\Leftrightarrow L_5,\\
& &\e{23}\Leftrightarrow L_6,\e{24}\Leftrightarrow L_7,\e{25}\Leftrightarrow L_8,\e{26}\Leftrightarrow L_9,\e{34}\Leftrightarrow L_{10},\\
& &\e{35}\Leftrightarrow L_{11},\e{36}\Leftrightarrow L_{12},\e{45}\Leftrightarrow L_{13},\e{46}\Leftrightarrow L_{14},\e{56}\Leftrightarrow L_{15}\\[3pt]
%
\cl{2}{4}& {\mathfrak{su}}(2,2,\bbC)&\e{12}\Leftrightarrow -S_{15},\e{13}\Leftrightarrow -S_4,\e{14}\Leftrightarrow -S_8,\e{15}\Leftrightarrow -S_{11},\e{16}\Leftrightarrow -S_{13},\\
& &\e{23}\Leftrightarrow -S_5,\e{24}\Leftrightarrow -S_9,\e{25}\Leftrightarrow -S_{12},\e{26}\Leftrightarrow -S_{14},\e{34}\Leftrightarrow -S_{1},\\
& &\e{35}\Leftrightarrow -S_{2},\e{36}\Leftrightarrow -S_{3},\e{45}\Leftrightarrow -S_{6},\e{46}\Leftrightarrow -S_{7},\e{56}\Leftrightarrow -S_{10}\\[3pt]
%
\cl{0}{6}& {\mathfrak{su}}(4,\bbC)&\e{12}\Leftrightarrow U_1,\e{13}\Leftrightarrow U_2,\e{14}\Leftrightarrow U_3,\e{15}\Leftrightarrow U_4,\e{16}\Leftrightarrow U_5,\\
& &\e{23}\Leftrightarrow U_6,\e{24}\Leftrightarrow U_7,\e{25}\Leftrightarrow U_8,\e{26}\Leftrightarrow U_9,\e{34}\Leftrightarrow U_{10},\\
& &\e{35}\Leftrightarrow U_{11},\e{36}\Leftrightarrow U_{12},\e{45}\Leftrightarrow U_{13},\e{46}\Leftrightarrow U_{14},\e{56}\Leftrightarrow U_{15}\\[3pt]
\end{array}
\]
Lie algebra matrix generators  on the right-hand side of in the
equivalence relations:
%&H=\begin{bmatrix}1&0\\0&-1\end{bmatrix},\quad
%X=\begin{bmatrix}0&1\\0&0\end{bmatrix},\quad
%Y=\begin{bmatrix}0&0\\1&0\end{bmatrix}.\\
%&z=H,\quad y=Y-X,\quad x=Y+X.\\\blue{Cia sumazinau srifta}
\[\small
\!\begin{aligned}
\mkern-18mu \sig{x}&=\begin{bmatrix}0&1\\1&0\end{bmatrix},\,&
\mkern-18mu \sig{y}&=\begin{bmatrix}0&-\ii\\\ii&0\end{bmatrix},\,&
\mkern-18mu \sig{z}&=\begin{bmatrix}1&0\\0&-1\end{bmatrix},
\end{aligned}
\]
%
\[\small
\setlength{\arraycolsep}{1pt}
\!\begin{aligned}
\mkern-18mu p_1&=\begin{bmatrix}-\ii\sig{y}&0\\0&\ii \sig{y}\end{bmatrix},&
\mkern-18mu p_2&=\begin{bmatrix}-\sig{z}&0\\0&\sig{z}\end{bmatrix},&
\mkern-18mu p_3&=\begin{bmatrix}-\sig{x}&0\\0&-\sig{x}\end{bmatrix},&
\mkern-18mu p_4&=\begin{bmatrix}\sig{x}&0\\0&-\sig{x}\end{bmatrix},\\
\mkern-18mu p_5&=\begin{bmatrix}-\sig{z}&0\\0&\sig{z}\end{bmatrix},&
\mkern-18mu p_6&=\begin{bmatrix}-\ii\sig{y}&0\\0&-\ii\sig{y}\end{bmatrix}.
\\[4pt]
\mkern-18mu u_1&=\begin{bmatrix}\ii\sig{z}&0\\0&\ii\sig{z}\end{bmatrix},&
\mkern-18mu u_2&=\begin{bmatrix}\ii\sig{y}&0\\0&\ii\sig{y}\end{bmatrix},&
\mkern-18mu u_3&=\begin{bmatrix}\ii\sig{x}&0\\0&-\ii\sig{x}\end{bmatrix},&
\mkern-18mu u_4&=\begin{bmatrix}\ii\sig{x}&0\\0&\ii\sig{x}\end{bmatrix},\\
\mkern-18mu u_5&=\begin{bmatrix}-\ii\sig{y}&0\\0&\ii\sig{y}\end{bmatrix},&
\mkern-18mu u_6&=\begin{bmatrix}\ii\sig{z}&0\\0&-\ii\sig{z}\end{bmatrix}.
\end{aligned}
\]
%
\[\small
\setlength{\arraycolsep}{1pt}
\!\begin{aligned}
\mkern-18mu P_1&=\begin{bmatrix}-\ii\sig{y}&0\\0&-\ii\sig{y}\end{bmatrix},&
\mkern-18mu P_2&=\begin{bmatrix}0&\sig{z}\\-\sig{z}&0\end{bmatrix},&
\mkern-18mu P_3&=\begin{bmatrix}0&\sig{z}\\\sig{z}&0\end{bmatrix},&
\mkern-18mu P_4&=\begin{bmatrix}-\sig{z}&0\\0&\sig{z}\end{bmatrix}\\
\mkern-18mu P_5&=\begin{bmatrix}0&-\sig{x}\\ \sig{x}&0\end{bmatrix},&
\mkern-18mu P_6&=\begin{bmatrix}0&-\sig{x}\\-\sig{x}&0\end{bmatrix},&
\mkern-18mu P_7&=\begin{bmatrix}\sig{x}&0\\0&-\sig{x}\end{bmatrix},&
\mkern-18mu P_8&=\begin{bmatrix}-1_2&0\\0&1_2\end{bmatrix},\\
\mkern-18mu P_9&=\begin{bmatrix}0&-1_2\\-1_2&0\end{bmatrix},&
\mkern-18mu P_{10}&=\begin{bmatrix}0&-1_2\\1_2&0\end{bmatrix}.
\end{aligned}
\]
%
\[\small
\setlength{\arraycolsep}{1pt}
\!\begin{aligned}
\mkern-18mu U_1&=\begin{bmatrix}\ii\sig{z}&0\\0&\ii\sig{z}\end{bmatrix},&
\mkern-18mu U_2&=\begin{bmatrix}\ii\sig{y}&0\\0&\ii\sig{y}\end{bmatrix},&
\mkern-18mu U_3&=\begin{bmatrix}\ii\sig{x}&0\\0&-\ii\sig{x}\end{bmatrix},&
\mkern-18mu U_4&=\begin{bmatrix}0&\ii\sig{x}\\\ii\sig{x}&0\end{bmatrix},\\
\mkern-18mu U_5&=\begin{bmatrix}0&-\sig{x}\\\sig{x}&0\end{bmatrix},&
\mkern-18mu U_6&=\begin{bmatrix}\ii\sig{x}&0\\0&\ii\sig{x}\end{bmatrix},&
\mkern-18mu U_7&=\begin{bmatrix}-\ii\sig{y}&0\\0&\ii\sig{y}\end{bmatrix},&
\mkern-18mu U_8&=\begin{bmatrix}0&-\ii\sig{y}\\-\ii\sig{y}&0\end{bmatrix},\\
\mkern-18mu U_9&=\begin{bmatrix}0&\sig{y}\\-\sig{y}&0\end{bmatrix},&
\mkern-18mu U_{10}&=\begin{bmatrix}\ii\sig{z}&0\\0&-\ii\sig{z}\end{bmatrix},&
\mkern-18mu U_{11}&=\begin{bmatrix}0&\ii\sig{z}\\\ii\sig{z}&0\end{bmatrix},&
\mkern-18mu U_{12}&=\begin{bmatrix}0&-\sig{z}\\\sig{z}&0\end{bmatrix},\\
\mkern-18mu U_{13}&=\begin{bmatrix}0&-1_2\\ 1_2&0\end{bmatrix},&
\mkern-18mu U_{14}&=\begin{bmatrix}0&-\ii 1_2\\-\ii 1_2&0\end{bmatrix},&
\mkern-18mu U_{15}&=\begin{bmatrix}\ii 1_2&0\\0&-\ii 1_2\end{bmatrix}.
\end{aligned}
\]
%
\[\small
\setlength{\arraycolsep}{1pt}
\!\begin{aligned}
\mkern-18mu L_1&=\begin{bmatrix}0&1_2\\- 1_2&0\end{bmatrix},&
\mkern-18mu L_2&=\begin{bmatrix}0&\ii\sig{y}\\\ii\sig{y}&0\end{bmatrix},&
\mkern-18mu L_3&=\begin{bmatrix}0&\sig{z}\\\sig{z}&0\end{bmatrix},&
\mkern-18mu L_4&=\begin{bmatrix}1_2&0\\0&-1_2\end{bmatrix},\\
\mkern-18mu L_5&=\begin{bmatrix}0&-\sig{x}\\-\sig{x}&0\end{bmatrix},&
\mkern-18mu L_6&=\begin{bmatrix}-\ii\sig{y}&0\\0&\ii\sig{y}\end{bmatrix},&
\mkern-18mu L_7&=\begin{bmatrix}-\sig{z}&0\\0&\sig{z}\end{bmatrix},&
\mkern-18mu L_8&=\begin{bmatrix}0&1_2\\1_2&0\end{bmatrix},\\
\mkern-18mu L_9&=\begin{bmatrix}\sig{x}&0\\0&-\sig{x}\end{bmatrix},&
\mkern-18mu L_{10}&=\begin{bmatrix}\sig{x}&0\\0&\sig{x}\end{bmatrix},&
\mkern-18mu L_{11}&=\begin{bmatrix}0&\ii\sig{y}\\-\ii\sig{y}&0\end{bmatrix},&
\mkern-18mu L_{12}&=\begin{bmatrix}\sig{z}&0\\0&\sig{z}\end{bmatrix},\\
\mkern-18mu
L_{13}&=\begin{bmatrix}0&\sig{z}\\-\sig{z}&0\end{bmatrix},&
\mkern-18mu
L_{14}&=\begin{bmatrix}\ii\sig{y}&0\\0&\ii\sig{y}\end{bmatrix},&
\mkern-18mu
L_{15}&=\begin{bmatrix}0&\sig{x}\\-\sig{x}&0\end{bmatrix}.
\end{aligned}
\]
%
\[\small
\setlength{\arraycolsep}{1pt}
\!\begin{aligned}
\mkern-18mu S_1&=\begin{bmatrix}-\ii\sig{z}&0\\0&\ii\sig{z}\end{bmatrix},&
\mkern-18mu S_2&=\begin{bmatrix}-\ii\sig{y}&0\\0&-\ii\sig{y}\end{bmatrix},&
\mkern-18mu S_3&=\begin{bmatrix}\ii\sig{x}&0\\0&\ii\sig{x}\end{bmatrix},&
\mkern-18mu S_4&=\begin{bmatrix}0&-\sig{z}\\-\sig{z}&0\end{bmatrix},\\
\mkern-18mu S_5&=\begin{bmatrix}0&-\ii\sig{z}\\\ii\sig{z}&0\end{bmatrix},&
\mkern-18mu S_6&=\begin{bmatrix}-\ii\sig{x}&0\\0&\ii\sig{x}\end{bmatrix},&
\mkern-18mu S_7&=\begin{bmatrix}-\ii\sig{y}&0\\0&\ii\sig{y}\end{bmatrix},&
\mkern-18mu S_8&=\begin{bmatrix}0&-\ii 1_2\\\ii 1_2&0\end{bmatrix},\\
\mkern-18mu S_9&=\begin{bmatrix}0&1_2\\1_2&0\end{bmatrix},&
\mkern-18mu S_{10}&=\begin{bmatrix}\ii\sig{z}&0\\0&\ii\sig{z}\end{bmatrix},&
\mkern-18mu S_{11}&=\begin{bmatrix}0&\sig{x}\\\sig{x}&0\end{bmatrix},&
\mkern-18mu S_{12}&=\begin{bmatrix}0&\ii \sig{x}\\-\ii \sig{x}&0\end{bmatrix},\\
\mkern-18mu
S_{13}&=\begin{bmatrix}0&\sig{y}\\\sig{y}&0\end{bmatrix},&
\mkern-18mu
S_{14}&=\begin{bmatrix}0&\ii\sig{y}\\-\ii\sig{y}&0\end{bmatrix},&
\mkern-18mu S_{15}&=\begin{bmatrix}\ii 1_2&0\\0&-\ii
1_2\end{bmatrix}.
\end{aligned}
\]
\end{unit}
%
\begin{unit}\textit{Properties of matrix Lie generators in equivalence relations in}~\ref{bivectorAlgebraToLie}
%
\begin{itemize}

%
\item[$\bullet$]  The Lie algebra is represented by a group of
infinitesimal generators $K_i$, the commutators of which satisfy
$[K_i,K_j]=\sum_k c_{ij}^k K_k$. In~~\ref{bivectorAlgebraToLie}
the elementary GA bivectors and their matrix representation,
$\e{ij}\Leftrightarrow K_i$, provide the simplest possible set
consisting of a single nonzero structure constant,
namely,  $c_{ij}^k=\pm 1$. \\
%
\item[$\bullet$]  The generators in the equivalence
relations~\ref{bivectorAlgebraToLie} represent the irreducible
matrix representation of Lie algebra. All Lie  generators are
expressed through Pauli, unit and zero $2\times 2$ matrices, which
reflect the properties of  elementary bivectors $\e{ij}$ and
therefore are related with the 8-fold periodicity table in
unit~\ref{8foldperiod}.\\
%
\item[$\bullet$] The squares of matrices
in~\ref{bivectorAlgebraToLie} yields plus or minus unit
matrix. The sign  coincides with $\e{ij}^2$ and determines  the (diagonal) 
Killing form  of the respective Lie algebra.\\
%
\item[$\bullet$] Using the exponential map $\ee^{\varphi\cK_i}$,
where $\varphi\in\bbR$, one can generate one-parameter basis of Lie group, from which any Lie group element in turn can be obtained by multiplying one-parameter matrices. Alternatively, since general Lie algebra element may be written as linear combination of bivectors $\cB=\sum_{i,j}\alpha_{ij}\e{ij}\in\cG_n^2$, $\alpha_{ij}\in\bbR$ the exponential map of it generate the general element of the Lie group in a form of single exponent 
$\pm\ee^{\cB}\in\m{Spin}^{+}$ as well.\\
%
\item[$\bullet$] The construction of equivalence
table~\ref{bivectorAlgebraToLie} is facilitated by following
observations: 1) The number of bivectors of $\cl{p+1}{q}$ is the same as total number of vectors and bivectors of one dimension smaller $\cl{q}{p}$ algebra 2) The set of vectors+bivectors is closed under commutator operation. It then follows that we can use matrix representations of smaller GA algebra which are of the required dimensions. The observations also imply the following theorem.
\newline \textbf{Theorem}. The $\cl{p+1}{q}$  and $\cl{p}{q+1}$  bivector commutator algebras are isomorphic to vector + bivector commutator algebras of $\cl{q}{p}$ and $\cl{p}{q}$, respectively.

In turn, bivectors of $\cl{q}{p}$ and $\cl{p}{q}$ then make (maximal) commutator subalgebras of $\cl{p+1}{q}$ and $\cl{p}{q+1}$, respectively. The statement is the generalization of the well known space-time split into $\cl{3}{0}$ bivectors (rotations, which make subalgebra) and vectors (boosts, which do not form subalgebra) of the relativistic $\cl{1}{3}$ algebra (see example~\ref{splitDemo}).

Vector+bivector algebras are not the only possible subalgebras of
bivector commutator algebra.  For $p+q\ge5$ there exists
subalgebras which include other grade elements as well. For
example, $\cl{1}{4}$ bivectors have subalgebra, which include
elements of all grades of $\cl{1}{3}$ algebra. It means that
we can make Lie algebra from specific sets of generators which not
all are GA bivectors. For  $\cl{0}{n}$ and $\cl{n}{0}$ algebras,
however, all subalgebras emerge form vector+bivector split.\\
\end{itemize}
\end{unit}
\begin{example}\red{Neaisku, kodel ismetete si grazu pavyzdi, atstaciau}\label{splitDemo}
Let's represent $\cl{1}{3}$ bivector commutator algebra by $\cl{1}{2}$ and $\cl{3}{0}$ vector+bivector commutator algebras.

Bivectors of $\cl{1}{3}$ are $\{\e{12},\e{13},\e{14},\e{23},\e{24},\e{34}\}$. If we put $\cl{1}{2}$ vectors and bivectors in the following order $\{\e{1},\e{13},\e{12},\e{3},\e{2},\e{23}\}$, then it is easy to check that $\cl{1}{3}\Rightarrow \cl{1}{2}$ substitution  $\{\e{12},\e{13},\e{14},\e{23},\e{24},\e{34}\}\Rightarrow \{-\e{1},-\e{13},-\e{12},-\e{3},-\e{2},-\e{23}\}$
makes both commutator tables identical.

In the case of $\cl{3}{0}$ the properly ordered set of vectors+bivectors is $\{\e{1},\e{2},\e{3},\e{12},\e{13},\e{23}\}$. The substitution 
$\{\e{12},\e{13},\e{14},\e{23},\e{24},\e{34}\}\Rightarrow \{-\e{1}, -\e{2}, -\e{3}, -\e{12}, -\e{13}, -\e{23}\}$ then gives well known result of time-space split of 4D relativistic algebra into 3D~vectors (boosts in 4D) and bivectors (space rotations). This explains why 4D~space rotations (3D bivectors) make subalgebra, whereas boosts (3D vectors) don't. Analogous splits can be easily constructed for any other Clifford algebra.
\end{example}
%
\begin{example}
Using the commutation tables we shall explicitly check that Lie
algebra $\mathfrak{cl}(2,\bbC)$ and GA algebra \cl{3}{1} (also
\cl{1}{3} algebra)  are isomorphic. From table in the
unit~\ref{classgroup} we find that $\mathfrak{sl}(2,\bbC)$ algebra
has $2(n^2-1)=2(2^2-1)=6$ infinitesimal generators
$\mathcal{X}_i$. The table~\ref{bivectorAlgebraToLie} shows that
standard basis matrices that represent the $\mathfrak{sl}(2,\bbC)$
generators can be expressed by a set of three Pauli matrices and
Pauli matrices multiplied by imaginary unit:
\[\begin{split}
&\mathcal{X}_1=\hat{\sigma}_x,\ \mathcal{X}_2 =\hat{\sigma}_y,\
\mathcal{X}_3=-\hat{\sigma}_z, \mathcal{X}_4=\ii\hat{\sigma}_x,\
\mathcal{X}_5=\ii\hat{\sigma}_y,\ \mathcal{X}_6=-\ii\hat{\sigma}_z
\end{split}\]
The table below shows the commutators
$[\mathcal{X}_i,\mathcal{X}_j]/2$
\[\begin{array}{c|cccccc}
\tfrac{1}{2}[\mathcal{X}_i,\mathcal{X}_j]&\mathcal{X}_1&\mathcal{X}_2&\mathcal{X}_3&\mathcal{X}_4&\mathcal{X}_5&\mathcal{X}_6\\
\hline
\mathcal{X}_1&0&\mathcal{X}_6&-\mathcal{X}_5&0&-\mathcal{X}_3&\mathcal{X}_2\\
\mathcal{X}_2&-\mathcal{X}_6&0&\mathcal{X}_4&\mathcal{X}_3&0&-\mathcal{X}_1\\
\mathcal{X}_3&\mathcal{X}_5&-\mathcal{X}_4&0&-\mathcal{X}_2&\mathcal{X}_1&0\\
\mathcal{X}_4&0&-\mathcal{X}_3&\mathcal{X}_2&0&-\mathcal{X}_6&\mathcal{X}_5\\
\mathcal{X}_5&\mathcal{X}_3&0&-\mathcal{X}_1&\mathcal{X}_6&0&-\mathcal{X}_4\\
\mathcal{X}_6&-\mathcal{X}_2&\mathcal{X}_1&0&-\mathcal{X}_5&\mathcal{X}_4&0\\
\end{array}\]
from which structure constants of $\mathfrak{sl}(2,\bbC)$ are seen
to be zero and $\pm 1$.

In the relativity, the basis vectors of  algebra usually are
represented by larger, $4\times 4$ complex Dirac matrices
$\hat{\gamma}_i$ (see...).  \cl{1}{3} algebra contains the same
number of bivectors: $\tfrac{1}{2}n(n-1)=\tfrac{1}{2}4(4-1)=6$, so
the bivector group $\cG_2^4$ and $\Spin{1}{3}$ group (see the next
Section) consist of a set of matrix products denoted as
 $\sig{1}=\hat{\gamma}_1\hat{\gamma}_0$,
 $\sig{2}=\hat{\gamma}_2\hat{\gamma}_0$,
 $\sig{3}=\hat{\gamma}_3\hat{\gamma}_0$,
 $\Isig{1}=\hat{\gamma}_3\hat{\gamma}_2$,
 $\Isig{2}=\hat{\gamma}_1\hat{\gamma}_3$,
 $\Isig{3}=\hat{\gamma}_2\hat{\gamma}_1$.
Then the commutators of the above bivectors multiplied by
$\tfrac{1}{2}$  yield table with exactly the same structure constants $\pm 1$ or $0$,
\[\begin{array}{c|cccccc}
[\,.\, ,\, .\,]&\sig{1}&\sig{2}&\sig{3}&\Isig{1}&\Isig{2}&\Isig{3}\\
\hline
\sig{1}&0&\Isig{3}&-\Isig{2}&0&-\sig{3}&\sig{2}\\
\sig{2}&-\Isig{3}&0&\Isig{1}&\sig{3}&0&-\sig{1}\\
\sig{3}&\Isig{2}&-\Isig{1}&0&-\sig{2}&\sig{1}&0\\
\Isig{1}&0&-\sig{3}&\sig{2}&0&-\Isig{3}&\Isig{2}\\
\Isig{2}&\sig{3}&0&-\sig{1}&\Isig{3}&0&-\Isig{1}\\
\Isig{3}&-\sig{2}&\sig{1}&0&-\Isig{2}&\Isig{1}&0
\end{array}\]
Thus both $2\times 2$ and $4\times 4$ sets of matrices can be used
to obtain the Lie group elements by the exponential map from
infinitesimal matrix generators.
\end{example}
%

IKI CIA PERSKAICIAU IR PATAISIAU. TOLIAU NESPEJAU



%
\section{Lipschitz-Clifford, pin and spin groups}

For groups in GA the binary operation between elements is played by GP (for group operation details see~\ref{groupActions}).
The GA group action frequently can be interpeted in terms of reflections and rotations of geometric objects. Roughly, four types of groups can
be distinguished in~GA:
$\textit{Cl}_{p,q}^*\supset\Gamma_{p,q}\supset\m{Pin}_{p,q}\supset\m{Spin}_{p,q}$,
which are, respectively, group of invertible MV's,
Lipschitz-Clifford, pin and spin groups. The table below provides a
more detailed classification~\cite{Lundholm09,Che55}.
%\bibitem{Che55} C. Chevalley, {\it The construction and study of certain
%important algebras}, Publications of Mathematical Society of Japan No 1
%(Herald Printing, Tokyo, 1955).
\begin{unit} \textit{Groups in Clifford algebra}.
\begin{center}\label{TypesOfGroups}
%\setlength{\extrarowheight}{1ex}
%\arraycolsep=0pt
%\parbox[t]{0.9\textwidth}{\textit{}}
%\setlength{\tabcolsep}{1pt}
%\resizebox{\textwidth}{!}{%
%\renewcommand*{\arraystretch}{1}
%\noalign{\vskip -3pt}\hphantom{m}

% Acaus preambule
%\begin{tabular}{x{0.1em}| >{$}l<{$} >{$}l<{$} >{$}l<{$} >{$}l<{$}
%>{$}l<{$} >{$}l<{$} >{$}l<{$} >{$}l<{$}}
%\diag{.1em}{.1em}{\mbox{$\!\!p\,$}}{\mbox{\hphantom{n}\smash{$q$}}}
%Dargio preambule
\begin{tabular}{>{$}l<{$}  >{$}c<{$} >{}l<{} }
% tabular body
\textrm{Group}&\textrm{Definition}&\textrm{Name}\\
\hline
%\cline{2-3}
\mathit{Cl}^{*}_{p,q} &\{\m{M}\in\cl{p}{q}\,|\,\exists\,
\m{M}^{-1}:&group of all invertible MV's $\m{M}$\\
&\qquad \m{M}\m{M}^{-1}=\m{M}^{-1}\m{M}=1\}&\\[2pt]
\widetilde{\Gamma}_{p,q} &\{\m{M}\in\mathit{Cl}^{*}_{p,q}: \gradeinverse{\m{M}}v\m{M}^{-1}\subseteq \cV^{p+q}\}
&Lipschitz-Clifford group\\[2pt]
\Gamma_{p,q} &\{v_1 v_2\cdots v_{k}\in\mathit{Cl}^{*}_{p,q}:
v_i\in \cV^*\}
&versor group (in fact, $\Gamma_{p,q}=\widetilde{\Gamma}_{p,q}$)\\[2pt]
\Gamma_{p,q}^{+} &\{\Gamma_{p,q}\cap \mathit{Cl}^{+}_{p,q}\}
&special Lipschitz-Clifford group\\[2pt]
\Pin{p}{q} & \{\m{R}\in \Gamma_{p,q}:  \m{R}\reverse{\m{R}}=\reverse{\m{R}} \m{R}=\pm 1\}
&pin group, i.~e. the group of\\
& & unit (normalised) versors\\[2pt]
\Spin{p}{q} & \{\m{R}\in(\Pin{p}{q}\cap \mathit{Cl}^{+}_{p,q}): \m{R}\reverse{\m{R}}=\pm 1\}
&spin group, i.~e. the group of  \\
& &even unit versors, $\Spin{p}{q}\cong\Spin{q}{p}$\\[2pt]
\Spin{p}{q}^{+} & \{\m{R}\in\Spin{p}{q}: \m{R}\reverse{\m{R}}=+1\}
&rotor group \\
\end{tabular}
%
\end{center}
Here $\cV^*$ is a space of invertible (nonisotropic) vectors,
i.~e., $\cV^*:= \{v\in\cV^n: v^2\neq 0\}$.  $v\in\cV^{n}$ denotes
a vector in $n=p+q$ dimension vector space $\cV^{n}$. The elements
of versor group are finite GP's of invertible vectors, which are
called versors.  Groups may also be generated by using just unit vectors, and
in the case of $\m{Spin}$ and $\m{Spin}^+$, only an even number of
such unit vector factors is needed. Elements made of product of even number of vectors are rotations. The \Pin{p}{q} group in addition includes reflections, i.~e. it's elements contain product of noneven number of vectors as well. The elements $\m{R}$ of \Pin{p}{q} and \Spin{p}{q} groups are called pinors and spinors, respectively. The elements of $\m{Spin}^+$ are also called rotors and are intimately connected to orthogonal groups. The highly non-trivial fact is that the Lipschitz-Clifford group presented in the table coincides with the
versor group (see unit~\ref{inverseMVfactorization}).
\end{unit}
%
\begin{unit} \textit{Canonical actions in GA groups}.\label{groupActions}
  The following group actions (operations) are consistent with geometric product
of GA.
\begin{center}\label{groupActions}
\begin{tabular}{>{$}l<{$}  >{$}c<{$} >{}l<{} }
% tabular body
\textrm{Operation}&\textrm{Definition of map}&\textrm{Name}\\
\hline
L_x &\{\cl{p}{q}\to \mathrm{End} \cl{p}{q}\}&left action\\
  &x\mapsto L_x: y\mapsto x y&\\[2pt]
R_x &\{\cl{p}{q}\to \mathrm{End} \cl{p}{q}\}&right action\\
&x\mapsto R_x: y\mapsto y  x &\\[2pt]
\Ad_x & \{\mathit{Cl}^{*}_{p,q}\to \mathrm{End} \cl{p}{q}\}
&adjoint action\\
&x\mapsto \Ad_x: y\mapsto x y x^{-1} &\\[2pt]
\widetilde{\Ad}_x & \{\mathit{Cl}^{*}_{p,q}\to \mathrm{End}
\cl{p}{q}\}
&twisted adjoin action \\
&x\mapsto \widetilde{\Ad}_x: y\mapsto \gradeinverse{x} y x^{-1} &($\gradeinverse{x}$ denotes grade inverse)\\
\end{tabular}
%
\end{center}
$x$ and $y$ are group elements. The $\mathrm{End}$ here denotes the space of
endomorphisms i.~e. space of structure preserving maps (see unit xxx). Example of a vector space endomorphism is a linear map  $f: \cV \to \cV$ of vector space to itself. Example of group $G$ endomorphism $f$ is a group homomorphism $f: G \to G$.

Left and right actions $L$ and $R$ are the {\it algebra
homomorphisms}, i.~e., they preserve both the summation and
multiplication structure, while $\Ad_x$ and $\widetilde{\Ad}_x$
are only {\it group homomorphisms} (i.~e. preserve multiplicative structure only).

One can verify, that $\Ad_x$ always is an outermorphism (preserves the outer
product, see Section~\ref{outermorphism}), while twisted adjoint action $\widetilde{\Ad}_x$ is not.
Both of these adjoint actions, however, agree on subgroup of even
invertable elements, i.~e. on $\mathit{Cl}^{*}_{p,q}\cap
\mathit{Cl}^{+}_{p,q}$.

The twisted adjoint action $\widetilde{\Ad}_x$ takes into account
the graded structure and, therefore, plays more important role in
geometric algebra than simple adjoint action $\Ad_x$.
\end{unit}
%Literature used: Porteous~\cite{Porteous95}, Hall~\cite{Hall03},
%Meinrenken~\cite{Meinreken05} labai abstrakti knyga,
%Gallier~\cite{Gallier08}, Garling~\cite{Garling11} about Lie and
%spin groups,
%\cite{Lundholm2009}.

%
\begin{figure}[b,t]
\centering
\includegraphics[width=4cm]{figures//triangleL}
\caption{Unilateral triangle in $\e{12}$ plane with two unit
vectors $\hat{a}$ and $\hat{b}$ that generate the triangle
group\label{fig:triangleL}}
\end{figure}

\begin{example}\label{unilateraltriangle} 
  Let's describe all possible transformations of planar unilateral triangle (see Figure~\ref{fig:triangleL}) into itself by means of geometric algebra group actions. As an initial vectors we will take two unit vectors  $\hat{a}$ and  $\hat{b}$ drawn from the center of the triangle. The vector $\hat{a}$ points to the trianle corner and the $\hat{b}$ to the center point of triangle side, i.~e. the $\hat{b}$ is perpendicular to that side. The
angle between the vectors is $\pi/3$. Three  rotation by
$\varphi=\{2\pi/3,4\pi/3,2\pi\}$ bring the triangle to itself. The
respective rotation transformation is
$\m{R}_{\varphi}(v)=Rv\widetilde{R}$, where $R$ is the half-angle
rotor
$R=\hat{a}\hat{b}=\ee^{\varphi\e{12}/2}=\cos\tfrac{\varphi}{2}+\e{12}\sin\tfrac{\varphi}{2}$,
and $v$ is an arbitrary in-plane vector. So, the clockwise
rotations by angles $\varphi=\{2\pi/3,4\pi/3,2\pi\}$ can be
described by three rotors which at the same time are  versors (i.~e.~products of two, four and six vectors, respectively),
\[\mathrm{R}=\{R=\hat{a}\hat{b},\,R^2=(\hat{a}\hat{b})^2,\,R^3=(\hat{a}\hat{b})^3=-1\}. \]
The triangle also possesses  three symmetry reflectors determined
by three lines perpendicular to sides of the triangle that are the
versors as well,
\[\mathrm{r}=\{\hat{b},\,\hat{b}R,\,\hat{b}R^2\}\]
If the identity transformation is added then we may write a set of
versors
$S=\{1,\mathrm{R},\mathrm{r}\}=\{1,R,R^2,R^3,\hat{b},\hat{b}R,\hat{b}R^2\}$
that is directly connected with triangle group. The vectors
$\hat{a}$ and $\hat{b}$ may be treated as generators since they
allow to construct all transformations that bring the triangle to
itself. Transformation that  preserve the structure of the group
(in our case the transformations of triangle to itself) is called
the \textit{endomorphism}. The rotation and  reflection
transformations of the arbitrary in-plane vector $v$ are
described, respectively,  by $R_iv\widetilde{R}_i$ and
$-\hat{r}_iv\hat{r}_i$, where $\hat{r}_i$ is the unit vector. Both
formulas can be united into single, $(-1)^pS_ivS_i^{-1}$, where
$S_i^{-1}=\widetilde{S}_i$ and $p$ is the parity  of number of
vectors (even or odd) in the versor $S_i$. With the help of
twisted adjoint action defined in unit \ref{groupActions} the
triangle symmetry transformations  can be written
\[\mathrm{Ad}_{S_i}{v}=\gradeinverse{S}_ivS_i^{-1}, \]
where grade inversion now takes into account even-odd parity
automatically. We see that the twisted adjoint transformation
rather then rotors/ reflectors make up a GA group of the regular
triangle.
\end{example}
%
\begin{example}  \blue{Naujas Example} Let us take a set of the versors
$S=\{1,R,R^2,R^3=-1,b,bR,bR^2\}$ from
Example~\ref{unilateraltriangle}. If we multiply from right by $R$
and remember the properties $\hat{a}^2=\hat{b}^2=1$,
$R=\hat{a}\hat{b}$ and $(\hat{a}\hat{b})^3=1$, we find that
$SR=\{R,R^2,R^3,-R,\hat{b}R,\hat{b}R^2,-\hat{b}\}$. The elements
now are in a different order, however, some of them acquire the
negative sign.  The same property is observed after multiplication
from left,
$RS=\{R,R^2,R^3,-R,-\hat{b}R^2,\hat{b},R\hat{b}R^2=\hat{b}R\}$.
Thus, left/right multiplication not only reshuffles the versors
but may change their signs. If instead the adjoint elements
$\gradeinverse{Ad}_{S_iR}$  and $\gradeinverse{Ad}_{RS_i}$   are
used then the minus sign will vanishes and we will recover the
full group of a regular triangle, albeit with elements in a
different positions. The same remains true if we multiply $S$ by
all remaining rotors and reflectors.

In a similar way employing two unit vectors that make a different
angle $\varphi$, the remaining planar point groups can be
constructed for other regular polygons:  $\varphi=\pi/2$ (two
points), $\pi/4$ (square) and $\pi/6$ (regular hexagon). The
pentagon does belong to planar point groups. The space groups in
3D can be constructed with three non-planar vectors-generators
$\hat{a}$, $\hat{b}$ and $\hat{c}$. More examples of
crystallographic groups including translation from the point of
view of GA can be found in~\cite{Hestenes07,Hitzer13}.
%@ARTICLE{Hestenes07,
%   author=" D.~Hestenes and J.~Holt",
%   title="The crystallographic space groups in geometric algebra",
%   year="2007",
%   journal="J.~Math. Phys.",
%   volume="48",
%   number="2",
%   pages="023514"
%}
%@ARTICLE{Hitzer13,
%   author=" E.~Hitzer and D.~Ichikawa",
%   title="Representation of crystallographic subperiodic groups in {C}lifford's geometric  algebra",
%   year="2013",
%   journal="Adv. Appl. {C}lifford Algebras",
%   volume="23",
%   number="",
%   pages="887-906"
%}
\end{example}

\blue{Cartano-Dieudonne teorema perkeliau i transformacijas, i
skyriu  'Orthogonal transformations'. Man atrodo ten jai geresne
vieta.}
%
\begin{unit}\textit{Properties of pin and spin groups.}
\begin{itemize}
\item[$\bullet$] The following isomorphism between
Lipschitz-Clifford groups and their classical counterparts exists
\blue{(examples are given in the second column)}:
\begin{align*}
  &\widetilde{\Ad}: \Pin{p}{q}\big/\{\pm 1\}\cong\mathsf{O}(q,p)&& \mathsf{O}(3)\cong \mathsf{SO}(3)\otimes \Z_2\quad \textrm{(globally)}\\
    &\widetilde{\Ad}: \Spin{p}{q}\big/\{\pm 1\}\cong\mathsf{SO}(q,p)&& \Spin{3}{0}\cong \mathsf{SU}(2)\\
    &\widetilde{\Ad}: \mathsf{Spin}^{+}_{p,q}\big/\{\pm 1\}\cong\mathsf{SO}^+(q,p)&& \mathsf{SU}(2)\cong  \mathsf{SO}(3)\otimes\Z_2\quad  \textrm{(fiber bundle)}
\end{align*}
\blue{Ka cia reiskia simbolis} $\widetilde{\Ad}:$ \blue {Ka jis
pasako? Su dvitaskiu man zinomas funkcijos apibrezimas}
$f:x\mapsto y$. \blue{Pagal mane adjoint simbolis nereikalingas.
Yra toks daiktas kaip quotient mapping} $q:\m{G}\mapsto\m{G}/\{\pm
1\}$, kur $\m{G}/\{\pm 1\}$ yra vadinama quotient group. Todel
$\m{SO}(3)\cong\m{SU}(2)/\{\pm 1\}$ t.y. eliminavus polius,
gauname izomorfizma. Terminas fiber bundle niekur zinyne nebus
sutinkamas, todel geriau sio termino nerasyti arba pavyzdi
uzrasyti kazkaip kitaip, gal, $\m{SO}(3)\cong\m{SU}(2)/\{\pm 1\}$

Here $\mathsf{O}(P,Q)$, $\mathsf{SO}(P,Q)$ and
$\mathsf{SO}^+(P,Q)$ are the orthogonal, special orthogonal and
special orthogonal groups with the determinant $+1$,
respectivelly. The $\{\pm 1\}$ denotes the two element kernel
$\Z_2=\{1,-1\}$ of the mapping. Thus, the Clifford groups are
double coverings of their classical orthogonal counterparts. These
coverings are universal. This has natural explanation in terms of
vector reflections (realised as twisted adjoin action) $w\to-v w
v^{-1}$, which   yields the same result if $v\to -v$.

From this we have that the number of bivectors of $\Spin{p}{q}$,
which is given by binomial
$\left(\begin{smallmatrix}n\\2\end{smallmatrix}\right)=
\tfrac{n!}{2!(n-2)!}=\tfrac{1}{2}n(n-1)$, where $n=p+q$, matches
the number of infinitesimal generators of $\mathsf{SO}(q,p)$
(see~\ref{classgroup}) and is equal to $n_{\mathrm{G}}=0$, $1$,
$3$, $6$, $10$ and $15$ for $p+q=1,2,3,4,5$ and $6$, respectively.

\item[$\bullet$]  The spin groups with opposite signature are
isomorphic
  \[
    \Spin{p}{q}\cong \Spin{q}{p}
  \]
 The respective pin groups, in general, are not isomorphic, $\Pin{p}{q}\ncong \Pin{q}{p}$, except when $p-q=0\ \mod (4)$
 (\red{neradau kur patikrint}).
%
\item[$\bullet$] \red{Manau abu zemiau uzrasyti izomorfizmai gali
buti NETEISINGI. Vietoje} $\otimes$ ir $\times$, cia turi buti
"semidirect products". Taip yra todel, kad kernelis nera
komutatyvus, Taigi, reikia savoku "normal subroup", "factor group"
ir pan. Atkreipkit demesi, kad Lundas toj vietoj paraso tiesiog
taska ir nieko nemini. Pas Varlamov tais klausimais daug daugiau
daugiau parasyta, bet siu iskaidymu ten nera, o panasiuose yra
semidirect products.
 For a mixed signature $p,q\ge 1$ and any pair
of orthogonal vectors $\e{+}$ and $\e{-}$ such that $\e{+}^2=1$,
$\e{-}^2=-1$ one has that~\cite{Lundholm09} p.65
\[\begin{split}
&\Pin{p}{q}=\mathsf{Spin}^{+}_{p,q}\otimes\{1,\e{+},\e{-},\e{+}\e{-}\}\\
&\Spin{p}{q}=\mathsf{Spin}^{+}_{p,q}\times\{1,\e{+}\e{-}\}.\\
\end{split}\]
\item[$\bullet$] For Euclidean $(p=n,q=0)$ and anti-Euclidean
$(p=0,q=n)$ signatures one has that
\[\Pin{p}{q}=\mathsf{Spin}^{+}_{p,q}\times\{1,\e{}\}, \]
where any $\e{}\in\mathcal{V}$ is such that $\e{}^2=\pm 1$. From
these properties follows  that it is sufficient to study only the
properties of the rotor group $\mathsf{Spin}^{+}_{p,q}$ to
understand the $\mathsf{Spin}$ and $\mathsf{Pin}$ groups.
\blue{nieko daugiau negaliu pasakyti}
\end{itemize}
\end{unit}
%
\begin{unit}\textit{Table of spin groups}.\label{tableSG}
The connected components (denoted by zero superscript) of real
spin groups $\Spin{p}{q}^0\cong\Spin{q}{p}^0$ that are associated
with orthogonal spaces of dimension $n\le 6$ and which are
isomorphic to classical groups are given
below~\cite{Trautman06,Porteous95}.
  \red{nuliu virsuje pazymejau connected component, nes + jau reiskia lygini,
   o zvaigzdute turincius inverse elementus pogrupius. Labai gali buti, kad  + ir 0
   reiskia ta pati (t.y rotoriu grupe, kuri yra identity componenteje).
    Reiktu pasiaiskinti dar. Bet tada reiktu tai paminet ir siose keliose
    vietose pakeist, nes dabar truputi klaidina.}
\begin{center}\label{tableSG}
%\setlength{\extrarowheight}{1ex}
%\arraycolsep=0pt
\parbox[t]{0.9\textwidth}{
\textit{Classical groups which are isomorphic to connected components of real spin groups $\Spin{p}{q}^0$}}
%\setlength{\tabcolsep}{1pt}
%\resizebox{\textwidth}{!}{%
%\renewcommand*{\arraystretch}{1}
%\noalign{\vskip -3pt}\hphantom{m}

% Acaus preambule
%\begin{tabular}{x{0.1em}| >{$}l<{$} >{$}l<{$} >{$}l<{$} >{$}l<{$}
%>{$}l<{$} >{$}l<{$} >{$}l<{$} >{$}l<{$}}
%\diag{.1em}{.1em}{\mbox{$\!\!p\,$}}{\mbox{\hphantom{n}\smash{$q$}}}
%Dargio preambule
\begin{tabular}{>{$}l<{$} | >{$}l<{$} >{$}l<{$} >{$}l<{$} >{$}l<{$}
>{$}l<{$} >{$}l<{$} >{$}l<{$} >{$}l<{$}}
% tabular body
&0&1&2&3&4&5&6\\
\cline{2-9}
0 &\{\pm 1\}&\mathsf{O}(1,\bbR)&\mathsf{U}(1,\bbC)&\mathsf{SU}(2,\bbC)&^2\mathsf{SU}(2,\bbC)&\mathsf{Sp}(2,\bbH)&\mathsf{SU}(4,\bbC)&\\
1 &\mathsf{O}(1,\bbR)&\mathsf{GL}(1,\bbR)&\mathsf{SL}(2,\bbR)&\mathsf{SL}(2,\bbC)&\mathsf{Sp}(1,1,\bbH)&\mathsf{SL}(2,\bbH)&&\\
2 &\mathsf{U}(1,\bbC)&\mathsf{SL}(2,\bbR)&^2\mathsf{Sp}(2,\bbR)&\mathsf{Sp}(4,\bbR)&\mathsf{SU}(2,2,\bbC) &&&\\
3 &\mathsf{SU}(2,\bbC)&\mathsf{SL}(2,\bbC)&\mathsf{Sp}(4,\bbR)&\mathsf{SL}(4,\bbR)&&&&\\
4 &^2\mathsf{SU}(2,\bbC)&\mathsf{Sp}(1,1,\bbH)&\mathsf{SU}(2,2,\bbC)&&&&&\\
5 &\mathsf{Sp}(2,\bbH)&\mathsf{SL}(2,\bbH)&&&&&&\\
6 &\mathsf{SU}(4,\bbC)&&&&&&&\\
\end{tabular}
%
\end{center}
Here $^2\mathsf{G}\equiv\mathsf{G}\oplus\mathsf{G}$.
The matrices that represent the spin groups are irreducible. The
first column and row of the table frequently are denoted as
$\mathsf{Spin}(n)$. These spin groups are compact and their connected component coinsides with the full spin group $\Spin{p}{0}^0\cong\Spin{p}{0}\cong\Spin{0}{p}$.

The larger spinor groups, $n>6$, are
related with the exceptional groups such as $E_6$, $E_7$, $E_8$,
$F_4$, $G_2$ rather then with classical ones~\cite{Meinrenken09}.
%
\end{unit}




Naudota literatura: Porteous book \cite{Porteous95}, Hall book
\cite{Hall03}, Lounesto book \cite{Lounesto97} p.224, Gilmore book
\cite{Gilmore74} p.52, Lundholm~\cite{Lundholm09}  p.71-82,
Garling~\cite{Garling11} p.151





\section{Complexification of GA}
There are two ways to complexify the GA: by complexification of
matrix representation and by vector space doubling.
\begin{unit}\textit{Complexification of matrix representation}.
\label{complexification}
In the real GA algebra the matrix representation given in the
table in~\ref{8foldperiod} is irreducible, i.~e.,  the basis
vectors cannot be represented by smaller matrices. In applications
this restriction may appear too demanding. For example, the table
in~\ref{8foldperiod} shows that \cl{1}{3} algebra  should be
represented by $2\times 2$ quaternionic matrices. However, the
relativistic Dirac equation with the same signature usually is
represented by larger matrices, the so called  $4\times 4$ Dirac
matrices.  The complex structure can be brought into GA by
complexifying the basis representation matrices. The result is a
complex quadratic space and complex GA denoted as
\textit{Cl}$_{n}^{\bbC}$ with $n=p+q$. In the complex
representation, all real GA's with the same $n=p+q$ will be
isomorphic:
\[\mathit{Cl}_{n}^{\bbC}\cong\cl{p}{q}\otimes\bbC,\qquad p+q=n.\]
The table below shows matrix representation  of the complexified
GA for different $p+q=n$.  Gallier~\cite{Gallier08}~p.37, Lounesto
book p.217

\vspace{2mm} \textit{Matrix representations of the complexified
GA's}
\[\begin{array}{c|ccccccccc}
n=p+q&0&1&2&3&4&5&6&7&8\\
\textit{Cl}_n^{\bbC}&\bbC&{2\bbC}&\bbC(2)&{2\bbC(2)}
&\bbC(4)&{2\bbC(4)}&\bbC(8)&{2\bbC(8)}&\bbC(16)
\end{array}\]
where ${2\bbC}$ is an abbreviation for
$\bbC\oplus\bbC$. From the table we have that
\[\textit{Cl}_{0,2n}\otimes\bbC\cong\textit{Cl}_{2n}^{\bbC}\cong\bbC(2^n),\qquad
\textit{Cl}_{0,2n+1}\otimes\bbC\cong\textit{Cl}_{2n+1}^{\bbC}\cong\bbC(2^n)\oplus\bbC(2^n)
\]
Therefore, at odd $n$ the complexified  representation matrices
have the block structure $\bbC(n)\oplus\bbC(n)$. This
may used to transform a pair of similar coupled equations to GA.
The complexification and periodicity of $2$ may be better
apprehended if the fundamental matrix
representation~\ref{8foldperiod} is plotted in coordinates $p+q=n$
vs. $p-q$.
%
$$
\setlength{\arraycolsep}{0pt}%
\begin{array}{c|ccccccccccccccc}
\multicolumn{16}{c}{\textit{Period-$8$ table of irreducible
matrix representations in the form $p+q=n$ vs. $p-q$}}\\[3pt]
 p-q&-7&-6&-5&-4&-3&-2&-1&0&1&2&3&4&5&6&7\\
 p+q& & & & & & & & & & & & & & &\\ \hline
 0&  & & & & & & &\bbR(1)&& & & & & &\\
 1& & & & & & &\bbC(1)&&^2\bbR(1)& & & & & &\\
 2& & & & & &\bbH(1) & &\bbR(2) &&\bbR(2)& & & & &\\
 3& & & & &^2\bbH(1) &&\bbC(2) &&^2\bbR(2)& &\bbC(2)& & & &\\
 4& & & &\bbH(2) &&\bbH(2) &&\bbR(4) &&\bbR(4)&&\bbH(2) & & &\\
 5& &&\bbC(4)&&^2\bbH(2) &&\bbC(4) &&^2\bbR(4)& &\bbC(4)&&^2\bbH(2) & &\\
 6& &\bbR(8)&&\bbH(4) &&\bbH(4) &&\bbR(8)&&\bbR(8)&&\bbH(4) &&\bbH(4) &\\
 7&^2\bbR(8)&&\bbC(8)& &^2\bbH(4)&&\bbC(8)&&^2\bbR(8)& &\bbC(8)&&^2\bbH(4)&&\bbC(8)\\
\end{array}
$$
This table shows that the relativistic theory in GA in the
complexified form ($p+q=4$ line) may be treated either by
\cl{1}{3} or \cl{3}{1} algebra.
 \end{unit}
%
\begin{example}
In $\bbH(2)$ the irreducible representation of basis vectors
$\{\e{1},\e{2},\e{3},\e{4}\}$ of \cl{1}{3} can be expressed as
\[\begin{split}
&\hat{e}_1=\left[\begin{array}{cc}1&0\\0&-1\end{array}\right],\quad
&\hat{e}_2=\left[\begin{array}{cc}0&\ii\\
\ii&0\end{array}\right],\quad
&\hat{e}_3=\left[\begin{array}{cc}0&\jj\\
\jj&0\end{array}\right],\quad
&\hat{e}_4=\left[\begin{array}{cc}0&\kk\\
\kk&0\end{array}\right]
\end{split}\]
where $\ii$, $\jj$ and $\kk$ are the quaternionic imaginaries:
$\ii^2=\jj^2=\kk^2=-1$, $\ii\jj\kk=-1$. The signs of squared
matrices in accord with \cl{1}{3} algebra signature are
$\hat{e}_1^2=\hat{1}_2$ and
$\hat{e}_2^2=\hat{e}_3^2=\hat{e}_4^2=-\hat{1}_2$. The product of
all matrices gives the pseudoscalar
$I_4=\hat{e}_1\hat{e}_2\hat{e}_3\hat{e}_4=\left[\begin{smallmatrix}0&-1\\1&0\end{smallmatrix}\right]$,
$I_4^2=-\hat{1}_2$, where  $\hat{1}_2$ is the $2\times 2$ unit
matrix.

The representation of \cl{1}{3} by complex $4\times 4$ Dirac gamma
matrices is
\[
\hat{e}_1\equiv\hat{\gamma}_0=\hat{\sigma}_3\otimes\hat{1}=\left[
\begin{array}{cc}
\hat{1} & 0 \\
0 & -\hat{1}
\end{array}
\right], \ \ \
\hat{e}_k\equiv\hat{\gamma}_k=-\ii\hat{\sigma}_2\otimes\hat{\sigma}_k=\left[
\begin{array}{cc}
0 & -\hat{\sigma}_k  \\
\hat{\sigma}_k & 0
\end{array}
\right]
\]
where $\hat{\gamma}_k$ $(k=1,2,3)$, are the Pauli matrices. The
squares  are $\hat{\gamma}_0^2=1$ and $\hat{\gamma}_k^2=-1$. This
$\bbC(4)$  representation is isomorphic  to $\bbH(2)$
representation.
\end{example}

%
\begin{unit}\textit{Complexification by space doubling}.
The other way  to complexify the algebra is to double the space
dimension to $2n$. Complex $n$ dimensional  space can be generated
by $2n$ dimensional real spaces. The complexification is adding to
the basis set $\{\e{i}\}$ vectors an equivalent  basis vectors
$\{\mathbf{f}_i\}$ with properties
\[\mathbf{f}_i\d\mathbf{f}_j=\delta_{ij},\quad\e{i}\d\e{j}=\delta_{ij},\quad\e{i}\d\mathbf{f}_{j}=0.
\]
A complex structure is introduced through a sum $n$ commuting
blades
\[\mathcal{J}=-\mathbf{f}_i\w\e{i}=\e{i}\w\mathbf{f}_i =\e{1}\mathbf{f}_1+\e{2}\mathbf{f}_2+\cdots +\e{n}\mathbf{f}_n.
\]
each playing the role of an imaginary in its own plane. The
imaginary bivector $\mathcal{J}$ satisfies
\[\begin{split}
&\mathcal{J}\d\mathbf{f}_i=\e{i},\quad\mathcal{J}\d\e{i}=-\mathbf{f}_i,\\
&\mathcal{J}\d(\mathcal{J}\d\e{i}) = -\e{i},\quad\mathcal{J}\d(\mathcal{J}\d\mathbf{f}_i)=-\mathbf{f}_i.\\
\end{split}\]
If $a=\sum_{i=1}^n a_i$ is a vector in $2n$ dimensional space with
components $a_i=\alpha_i\e{i}+\beta_i\mathbf{f}_i$, where
$\alpha_i$ and $\beta_i$ are scalars. Then the  product with
$\mathcal{J}$ takes over the role of unit imaginary $i$,
\[\mathcal{J}\d(\mathcal{J}\d a) =(a\d\mathcal{J})\d\mathcal{J}= -a\]
Similarly, $\mathcal{J}$  generates the rotation of the vector $a$
by angle~$\phi$ in each of  the $\e{i}\w\mathbf{f}_i$ planes,
\[e^{-\mathcal{J}\phi/2}a\, e^{\mathcal{J}\phi/2}=\cos\phi\, a+\sin\phi\, a\d\mathcal{J}.\]
Thus the map $a\rightarrow a\d\mathcal{J}$ is a  rotation by
$\pi/2$.
\end{unit}
%

\begin{unit}\textit{The complex spin groups}, Trautman~\cite{Trautman06}
\[\begin{split}
&\mathsf{Spin}(2,\bbC)=\Spin{2}{0}(\bbC)\cong\Spin{0}{2}(\bbC)\cong\bbC\backslash\{0\}\\
&\mathsf{Spin}(3,\bbC)=\Spin{3}{0}(\bbC)\cong\Spin{0}{3}(\bbC)\cong\mathsf{SL}(2,\bbC)\\
&\mathsf{Spin}(4,\bbC)=\Spin{4}{0}(\bbC)\cong\Spin{0}{4}(\bbC)\cong{^2\mathsf{SL}(2,\bbC)}\\
&\mathsf{Spin}(5,\bbC)=\Spin{5}{0}(\bbC)\cong\Spin{0}{5}(\bbC)\cong\mathsf{Sp}(4,\bbC)\\
&\mathsf{Spin}(6,\bbC)=\Spin{6}{0}(\bbC)\cong\Spin{0}{6}(\bbC)\cong\mathsf{SL}(4,\bbC)\\
\end{split}\]
\end{unit}
%
\begin{example} If one imagines two separate  axes $\e{1}$ and
$\mathbf{f}_1$ that belong to two copies of  \textit{Cl}$_{1,0}$
algebras then the complexified version of \textit{Cl}$_{1,0}$ can
be obtained after imposing the additional condition
$\mathcal{J}=\e{1}\mathbf{f}_1=-\mathbf{f}_1\e{1}$ on the
otherwise independent vectors $\e{1}$ and $\mathbf{f}_1$. The
basis of
\textit{Cl}$_{1}^{\bbC}\simeq$\textit{Cl}$_{1,0}\otimes\bbC$
then is
$\{1,\e{1},\mathbf{f}_1,\mathcal{J}\}\in$\textit{Cl}$_{2,0}$.
Therefore, the quantity $\mathcal{J}$ can be interpreted as an
imaginary number formed from two independent frames $\{\e{1}\}$
and $\{\mathbf{f}_1\}$.
\end{example}
 \blue{Apie CPT transformacijas
parasyta kitame skyriuje prie spinoriu} Doran book \cite{Doran03}
p.409, Porteous~\cite{Porteous04} p.37, Gallier~\cite{Gallier08},
p36, Lounesto book p.217



\endinput

%20170203
\blue{Sita pavyzdi radikaliai perdirbau, o kita ismeciau. Gal
zinai gera ir pamokanti pavyzdi is grupiu sitoje vietoje}
\begin{example} \textit{Continuous matrix
group $\mathsf{SO}(n)$}. When $n=2$ the group describes rotations
around a circle drawn in a plane. The general group element and
its inverse are
\[R(\theta)\equiv R_{\theta}=
\left[\begin{array}{cc}
 \cos\theta &-\sin\theta\\
\sin\theta &\cos\theta
\end{array}\right],\quad
R^{-1}_{\theta}= \left[\begin{array}{cc}
 \cos\theta &\sin\theta\\
-\sin\theta &\cos\theta
\end{array}\right],\quad\theta\in (0,2\pi)\]
The action of rotation operator on vector
$\bx=[x_1,x_2]^{\textrm{T}}$ expressed via group elements is
\[R_{\theta}\bx=
\left[\begin{array}{l}
x_1\cos\theta+x_2\sin\theta\\
x_2\cos\theta-x_1\sin\theta
\end{array}\right],\quad
R_{\theta}\bx R^{-1}_{\theta}= \left[\begin{array}{l}
x_1\cos 2\theta+x_2\sin 2\theta\\
x_2\cos 2\theta-x_1\sin 2\theta
\end{array}\right].\]
From this we find that the group axioms are satisfied:
\[R_{\theta_1}R_{\theta_2}=R_{\theta_1+\theta_2},\quad
R_{\theta}R_{\theta}^{-1}=1,\quad \textrm{det}R_{\theta}=1.\]
$R_{\theta_1}R_{\theta_2}=R_{\theta_2}R_{\theta_1}$, so the group
is Abelian. Since $\theta$ is group parameter and functions on
plane manifold are differentiable, $\mathsf{SO}(2)$ is the Lie
group. The generator of the corresponding \textit{Lie algebra} can
be found from
\[\hat{\Gamma}=\lim_{\theta\rightarrow 0}\frac{\partial R(\theta)}{\partial \theta} =
\left[\begin{array}{cc}0&1\\-1&0\end{array} \right]. \]

In terms of GA the Euclidean plane is described by \cl{2}{0}
algebra whose elements are $\{1,\e{1},\e{2},\e{12}\}$ and  which
have the following matrix representation
\begin{equation*}
\hat{1}=\left[
\begin{array}{cc}
1&0\\0&1
\end{array}
\right], \quad \hat{e}_1=\left[
\begin{array}{cc}
1&0\\0&-1
\end{array}
\right],\quad \hat{e}_2=\left[
\begin{array}{cc}
0&1\\1&0
\end{array}
\right],\quad  \hat{e}_{12}=\left[
\begin{array}{cc}
0&1\\-1&0
\end{array}
\right].
\end{equation*}
We see that Lie algebra generator obtained from  $\mathsf{SO}(2)$
and matrix representation of \cl{2}{0} bivector   coincide. This
is also confirmed by the fact that rotations of a vector in plane
can be described by GA rotor
\[\m{R}=\ee^{\theta\e{12}}=\cos\theta+\e{12}\sin\theta=
\left[\begin{array}{cc}
 \cos\theta &\sin\theta\\
-\sin\theta &\cos\theta
\end{array}\right]\]
where in the last expression $1$ and $\e{12}$ were replaced by
representative  matrices. In this case the rotation of the vector
$x=x_1\e{1}+x_1\e{2}$ in $\e{12}$ plane  is done with\footnote{The
underbar is used to distinguish a linear operator from elements of
the group.} $\underline{\m{R}}_{\theta}(x)=
\m{R}_{\theta}x\m{R}^{-1}_{\theta}=\m{R}_{\theta}^2 x$, or in
component form
$\underline{\m{R}}_{\theta}(x)=\ee^{2\theta\e{12}}x=(\cos
2\theta+\e{12}\sin 2\theta)(x_1\e{1}+x_2\e{2})=(x_1\cos
2\theta+x_2\sin 2\theta)\e{1}+( x_2\cos 2\theta-x_1\sin
2\theta)\e{2}$, which gives the same components as
$\mathsf{SO}(2,\bbR)$ group. This simple example once more shows that
Lie algebra generators and GA bivectors are related: both describe
the rotations in a vector space. This is not accidental. As table
in~\ref{tableSG} shows, the spin group of \cl{2}{0} is isomorphic
to $\mathsf{SO}(2)$.


Rotations in $n$-dimensional Euclidean  space may be performed by
more general groups $\mathsf{SO}(n)$, $n\ge 3$, which are
non-Abelian. However, already in $3$-D space the  rotation by two
parameter rotor $R(\theta,\varphi)\in\mathsf{SO}(3)$  cannot be
performed at all points of sphere:  the angle $\theta$ at North
and South poles is undefined. This happens because the group
$\mathsf{SO}(3)$ is not simply connected. On the other hand the
spinor group $\Spin{3}{0}$ of \cl{3}0{} algebra, which is a
covering group $\mathsf{SO}(3,0)$ and isomorphic to $\mathsf{SU}(2)$,
may realize all possible local rotation on sphere. This is
because, in general, the spinor groups $\Spin{n}{0}$ are
topologically simpler than $\mathsf{SO}(n)$, since $\Spin{n}{0}$
is simply connected, but $\mathsf{SO}(n)$ is not.
\end{example}
%




%20170201 ismestas kaip besikartojantis gabalas
The \textit{Lie group} $G$ satisfies group axioms  and, in
addition, depends on parameter(s) that are associated with
(elementary) bivectors. The Lie group is differentiable on a
manifold $\bbR^k$ (e.g. on a surface of sphere imbedded in
$\bbR^3$). If $g\in G$ and vector $x\in\bbR^{k}$ then there is a
differentiable mapping $f:\bbR^k\rightarrow G$ so that for each
$g\in G$ one can define a tangent space $\mathcal{T}_g$.

The Lie algebra elements $m$ and respective Lie group elements are
connected through the exponential map:
$m\rightarrow\cB\rightarrow\ee^{\xi \cB}$, where $\xi$ is the
parameter of the group and $\cB$ is the bivector. Thus, with every
elementary GA bivector one can associate Lie group parameter.
%

\begin{example}The \textit{discrete quaternionic group}
$Q=\Pin{0}{2}$ of algebra \cl{0}{2} is of the order $8$ and
consists of elements $Q=\{\pm 1,\pm\ii,\pm\jj,\pm\kk\}$. Six
elements (imaginaries) are of the order $4$, since for example
$\ii^4=1$. One element, namely $-1$ of order $2$, and the identity
elements $1$ of order $1$. The center of the group consists of two
elements, $Z(Q)=\{1,-1\}$. The multiplication table of the group
is
\[\begin{array}{r||rrrr|rrrr}
Q&1&\ii&\jj&\kk&-1&-\ii&-\jj&-\kk\\ \hline\hline
1&1&\ii&\jj&\kk&-1&-\ii&-\jj&-\kk\\
\ii&\ii&-1&\kk&-\jj&-\ii&1&-\kk&\jj\\
\jj&\jj&-\kk&-1&\ii&-\jj&\kk&1&-\ii\\
\kk&\kk&\jj&-\ii&-1&-\kk&-\jj&\ii&1\\ \hline
-1&-1&-\ii&-\jj&-\kk&1&\ii&\jj&\kk\\
-\ii&-\ii&1&-\kk&\jj&\ii&-1&\kk&-\jj\\
-\jj&-\jj&\kk&1&-\ii&\jj&\kk&-1&\ii\\
-\kk&-\kk&-\jj&\ii&1&\kk&\jj&-\ii&-1\\
\end{array}\]
From this table it follows that the elements are distributed over
$5$ equivalence classes: $\{1\}$,  $\{-1\}$,
$\{\ii,-\ii\}$,$\{\jj,-\jj\}$, $\{\kk,-\kk\}$. The character table
of the group coincides with the point group  $\mathsf{D}_4$. The
representation of group elements $\{1,\ii,\jj,\kk\}$ by complex
matrices, for example, may be of the form
\begin{equation*}
\hat{1}=\left[
\begin{array}{cc}
1 & 0 \\
0 & 1
\end{array}
\right],\quad \hat{\ii}=\left[
\begin{array}{cc}
\ii & 0 \\
0 & -\ii
\end{array}
\right],
\end{equation*}
%
\begin{equation*}
\hat{\jj}=\left[
\begin{array}{cc}
0 & 1 \\
-1 & 0
\end{array}
\right],\quad \hat{\kk}=\left[
\begin{array}{cc}
0 &\ii \\
\ii & 0
\end{array}
\right].
\end{equation*}
The elements $\{-1,-\ii,-\jj,-\kk\}$ are represented by same
matrices but with opposite signs.\\
The discrete groups may be used, for example, to describe  the
point symmetry of crystals. Group character tables (invariants)
are constructed to represent the structure of such groups. They
can be used to decompose an arbitrary MV to group elements.
\end{example}


% Ismesta Coquereaux lentele, nevertinga
$$
\begin{array}{cc|cccccccc}
\multicolumn{10}{c}{\textit{Period-$8$ table of irreducible
matrix representations of GA for $n=p+q>7$}}\\
 &q&0&1&7&8&9&10&11&12\\
 p& &&&&&&&&\\ \hline
 0&
  & & &^2\bbR(8)&\bbR(16)&\bbC(16)&\bbH(16)&^2\bbH(16)& \\
 1& & & &\bbR(16)&^2\bbR(16)&\bbR(32)&\bbC(32)& & \\
 7& &\bbC(8)&\bbH(8)& & & & & & \\
 8& &\bbR(16)&\bbC(16)& & & & & & \\
 9& &^2\bbR(16)&\bbR(32)& & & & & & \\
 10& &\bbR(32)&^2\bbR(32)& & & & & & \\
 11& &\bbC(32)& & & & & & & \\
 12& & & & & & & & &
\end{array}
$$
\blue{Butu gerai uzpildyti daugiau langeliu. Bet as nezinau. Cia
paimta is Coquaraux, AACA, vol.19, 673 (2009), Zvilgtelk.}
%
$\bbR(n)$, $\bbC(n)$ and $\bbH(n)$ are $n$th
order real, complex or quaternionic matrices. Symbol
$^2\mathbb{M}$ denotes a direct sum of matrices:
$^2\mathbb{M}=\mathbb{M}\bigoplus\mathbb{M}$.






****************************************************

A representation of GA by a group $\mathsf{G}$ is a transposition
from $\mathsf{G}$ onto GA vector space $\mathcal{V}$ by some
invertible linear map. Most groups have many different
representations. An arbitrary representation will break up into
irreducible representations. For many groups, the irreducible
representations have been classified. The reduction of an
arbitrary representation to irreducible ones, i.e. to the simplest
or fundamental ones, is widely used in physics, since this helps
to classify elementary states of a physical object, or to uncover
allowed and forbidden transitions (changes) within externally
perturbed physical object in terms of the group theory. Different
approaches are required depending on whether the group
$\mathsf{G}$ is a finite group, an infinite discrete group, or a
continuous Lie group.

*****************************************************

\vspace{5mm}
\noindent \textit{A simple example of
complexification}.  The Clifford algebra \text{Cl}$_{2,0}$ has
basis elements $\{1,\e{1},\e{2},\e{1}\e{2}\}$. The imaginary unit
is played by pseudoscalar $\e{1}\e{2}$ (an elementary oriented
area). To get the complex number $Z$ the vector $x=u\e{1}+v\e{2}$
is premultiplied by $\e{1}$,
\begin{equation*}
\e{1}x=u+v\e{1}\e{2}=u+vI=Z\, .
\end{equation*}
The complex conjugate of $Z$ then is premultiplification by
$\e{1}$ from right: $Z^{\dagger}=u-vI=u+v\e{2}\e{1}=x\e{1}$. Thus,
the complex number  $Z$ consists of scalar and bivector. It
generates rotation and dilation,
\begin{equation*}
Z=|Z| e^{i\phi}=|Z|\big(\cos\phi+i\sin\phi\big)\rightarrow
|Z|\big(\cos\phi+\e{1}\e{2}\sin\phi\big)=|Z|e^{\e{1}\e{2}\phi},
\end{equation*}
where $i=\sqrt{-1}$,  $|Z|=\sqrt{u^2+v^2}$ and $\tan\phi=v/u$.


*****************************************************

%
\begin{example}
The rotors $\m{R}_{\theta}=\ee^{\e{1}\e{2}\theta/2}$ in a plane
formed by \textit{Cl}$_{2,0}$ vectors $\e{1}$ and $\e{2}$  make a
continuous group as a function of angle~$\theta$, i.e. the
elements of the group can be represented as functions
$\ee^{\e{1}\e{2}\theta/2}\in\mathsf{G}
\;|\;\forall\theta\in\bbR, \theta=(0,2\pi)$. The rotor
brings the vector $x=x_1\e{1}+x_2\e{2}$ to a new vector
$x^{\prime}=\ee^{\e{1}\e{2}\theta/2}x\ee^{-\e{1}\e{2}\theta/2}=\m{R}_{\theta}x\tilde{\m{R}}_{\theta}$.
Here, the group operation~$\circ$ is a two-sided action of rotor
on a vector. It satisfies the group axioms:
1)~$\m{R}_{\theta_1}\m{R}_{\theta_2}=\m{R}_{\theta_1+\theta_2}$;
2)~$\big(\m{R}_{\theta_1}\m{R}_{\theta_2}\big)\m{R}_{\theta_3}=\m{R}_{\theta_1}\big(\m{R}_{\theta_2}\m{R}_{\theta_3}\big)$;
3)~The unit element is $\m{R}_{0}=1$;
4)~$\m{R}_{\theta}^{-1}=\m{R}_{-\theta}=\tilde{\m{R}}_{\theta}$.
In this case the group of rotors (or  rotations)  transforms the
vectors to other vectors the ends of which lie on a circle,
$x^2=x_1^2+x_2^2=\mathrm{const}$. The group  is Abelian, since
$\m{R}_{\theta_1}\m{R}_{\theta_2}=\m{R}_{\theta_2}\m{R}_{\theta_1}$,
and belongs to $\mathsf{SO}(2,0)$ (or in short $\mathsf{SO}(2)$ in
the classical notation).  It is a connected group.
\end{example}

***********************************************


\begin{example} The table in item~\ref{tableSG} shows that the
algebra \textit{Cl}$_{3,0}$ induces the spin group
$\mathsf{Sp}(2)\simeq \mathsf{SU}(2,\bbC)$ which can be
represented by $2\times 2$ complex matrices of the form
\[\hat{A}=\left[
\begin{array}{cc}
a&-b^*\\
b& a^*\\
\end{array} \right].\]
To satisfy the unitarity $\hat{A}\hat{A}^{\dagger}=\hat{1}$ the
matrix elements should be related by $aa^*+bb^*=1$. In general
this condition can be satisfied, for example,  by three real
angles, $a=\ee^{\ii(\alpha+\gamma)/2}\cos\frac{\beta}{2}$ and
$b=\ee^{\ii(\alpha-\gamma)/2}\sin\frac{\beta}{2}$. Thus, the group
$\mathsf{SU}(2,\bbC)$ may be thought of as a three
dimensional sphere embedded in a four dimensional space. The
points on a such sphere are characterized by three angles
$\alpha$, $\beta$ and~$\gamma$.
\end{example}

********************************************

$$\begin{array}{cc|ccccccc}
\multicolumn{9}{c}{\textit{Porteous\cite{Porteous95}GA spinor
groups $\mathsf{Spin}^{(+)}(p,q)$,
$p+q=n\le 6$}}\\
\multicolumn{9}{c}{\textit{expressed through classical groups}}\\
 &q&0&1&2&3&4&5&6\\
 p& &&&&&&&\\ \hline
 0& &\{\pm 1\}&\mathsf{O}(1)&\mathsf{U}(1)&\mathsf{Sp}(2)&^2\mathsf{Sp}(2)&\mathsf{Sp}(4)&\mathsf{SU}(4)\\
 1& &\mathsf{O}(1)&\mathsf{GL}(1,\bbR)&\mathsf{Sp}(2,\bbR)&\mathsf{Sp}(2,\bbC)&\mathsf{Sp}(2,2)&\mathsf{SU}^*(4)&\\
 2& &\mathsf{U}(1)&\mathsf{Sp}(2,\bbR)&^2\mathsf{Sp}(2,\bbR)&\mathsf{Sp}(4,\bbR)&\mathsf{SU}(2,2)&&\\
 3& &\mathsf{Sp}(2,\bbR)&\mathsf{Sp}(2,\bbC)&\mathsf{Sp}(4,\bbR)&\mathsf{SL}(4,\bbR)&&&\\
 4& &^2\mathsf{Sp}(2)&\mathsf{Sp}(2,2)&\mathsf{SU}(2,2)&&&&\\
 5& &\mathsf{Sp}(4)&\mathsf{SU}^*(4)&&&&&\\
 6& &\mathsf{SU}(4)&&&&&&\\
\end{array}$$

$^2\mathsf{M}=\mathsf{M}\oplus\mathsf{M}$ is a direct sum, for
example of matrices.
$*$~indicates complex conjugate matrix (or non-zero matrices?).\\



\subsection{Groups of \textit{Cl}$_{p,q}$}

\vspace{5mm} Paimta i{\v s} failo Berhards...pdf
\cite{Berhards11}\\
The Lie algebra includes most of the information of a
corresponding Lie group.
\[\begin{array}{ll}
\multicolumn{2}{c}{\textit{A dictionary between algebras and
groups}}\\ \hline
\textrm{Lie algebra\ }\mathfrak{g}  & \textrm{Lie group\ }\mathsf{G}\\
\textrm{Subalgebra\ }\mathfrak{h}\subseteq \mathfrak{g} &\textrm{Connected Lie subgroup\ }\mathsf{H}\subseteq\mathsf{G} \\
\textrm{Homomorphism\ }\mathfrak{h}\rightarrow\mathfrak{g} &\mathsf{H}\rightarrow\mathsf{G}\textrm{\ provided $\mathsf{H}$ simply connected} \\
\textrm{Module/representation\ }\mathfrak{h}\rightarrow\mathfrak{gl}(\mathcal{V}) &\textrm{Representation\ }\mathsf{G}\rightarrow\mathsf{GL}(\mathcal{V}) \\
 &  (\mathsf{G}\textrm{\ is simply connected})\\
page 32\dots &\dots \\
\dots &\dots \\
\dots &\dots \\
\end{array}\]

From Berhards~\cite{Berhards11} p.8
 $$\begin{array}{lll}
\multicolumn{3}{c}{\textit{The classical Lie groups and their
respective algebras}}\\ \hline
 \textrm{Group/Algebra} & \textrm{Group/Algebra} & \textrm{dim}_{\bbR}(\hat{A}) \\
 \textrm{name} &\textrm{description}& \\ \hline\hline
\mathbf{Compact}& &\\
 \mathsf{SO}(N,\bbR) &\{\hat{A}\in\textrm{Mat}(N,\bbR)|\hat{A}^{\mathsf{T}}\hat{A}=\hat{1},\det\hat{A}=1\}&\binom{N}{2}\\
 \mathfrak{so}(N,\bbR) &\{\hat{A}\in\textrm{Mat}(N,\bbR)|\hat{A}^{\mathsf{T}}+\hat{A}=\hat{0}\} &\\ \hline
 \mathsf{SU}(N) &\{\hat{A}\in\textrm{Mat}(N,\bbC)|\hat{A}^{\dagger}\hat{A}=\hat{1},\det\hat{A}=1\}&N^2-1\\
 \mathfrak{su}(N) &\{\hat{A}\in\textrm{Mat}(N,\bbC)|\hat{A}^{\dagger}+\hat{A}=\hat{0},\textrm{Tr}\hat{A}=0\} &\\ \hline
 \mathsf{Sp}(N,\bbR) &\{\hat{A}\in\textrm{Mat}(N,\bbH)|\hat{A}^{\mathsf{T}}\hat{A}=\hat{1}\}&2N^2+N\\
 \mathfrak{sp}(N,\bbR) &\{\hat{A}\in\textrm{Mat}(N,\bbH)|\hat{A}^{\dagger}+\hat{A}=\hat{0}\} &\\
\hline\hline
& &\\
 \mathsf{GL}(N,\bbH) &\{\hat{A}\in\mathsf{GL}(2N,\bbC)|\hat{J}\hat{A}=\hat{A}^{*}\hat{J}\}&4N^2\\
 \mathfrak{gl}(N,\bbH) &\{\hat{A}\in\textrm{Mat}(2N,\bbC)|\hat{J}\hat{A}=\hat{A}^*\hat{J}\} &\\
\hline\hline
\mathbf{Complex}& &\\
 \mathsf{SL}(N,\bbC) &\{\hat{A}\in\textrm{Mat}(N,\bbC)|\det\hat{A}=1\}&2(N^2-1)\\
 \mathfrak{sl}(N,\bbC) &\{\hat{A}\in\textrm{Mat}(N,\bbC)|\mathrm{Tr}\hat{A}=0\} &\\ \hline
 \mathsf{SO}(N,\bbC) &\{\hat{A}\in\textrm{Mat}(N,\bbC)|\hat{A}^{\dagger}\hat{A}=\hat{1},\det\hat{A}=1\}&N(N-1)\\
 \mathfrak{so}(N,\bbC) &\{\hat{A}\in\textrm{Mat}(N,\bbC)|\hat{A}^{\dagger}+\hat{A}=\hat{0}\} &\\ \hline
 \mathsf{Sp}(N,\bbC) &\{\hat{A}\in\textrm{Mat}(N,\bbC)|\hat{A}^{\mathsf{T}}\hat{J}\hat{A}=\hat{J}\}&2\binom{2n+1}{2}\\
 \mathfrak{sp}(N,\bbC) &\{\hat{A}\in\textrm{Mat}(N,\bbC)|\hat{A}^{\mathsf{T}}\hat{J}+\hat{J}\hat{A}=\hat{0}\} &\\
\end{array}$$
where the block matrix is
\[\hat{J}=\left[\begin{array}{cc}
\hat{0}&-\hat{1}\\
\hat{1}&\hat{0}\end{array}\right]\in\textrm{Mat}(2,\textrm{Mat}(N,\bbC))=
\textrm{Mat}(2n,\bbC).
\]
$\mathsf{GL}(n,\bbH)=\{\hat{A}\in\mathsf{GL}(2N,\bbC),\hat{J}\hat{A}=\hat{A}^*\hat{J}\}$
is closed linear group.

 \vspace{5mm}\noindent $\blacktriangleright$ Dynamical or
differential  representation, for example, the dynamical symmetry
of the Schr{\" o}dinger equation. It describes the angular
momentum operators acting on the intrinsic states and is related
to derivatives

\vspace{5mm} ****************************

 Palyginimui lentel{e}s
i{\v s} Vistoli~\cite{Vistoli07}, Meinrenken~\cite{Meinrenken09},
Porteous~\cite{Porteous95}, Lundholm~\cite{Lundholm09} p.71-82,
Garling~cit{Garling11} p.151
$$\begin{array}{l|l|l|l|l|l}
&Vistoli&Meinrenken&Porteous&Lundholm&Garling\\
\mathsf{Spin}(1,\bbC)&\mathbb{Z}/2\mathbb{Z}\simeq\mathbb{Z}_2&&\mathsf{O}(1)&\mathbb{Z}_2&\\
\mathsf{Spin}(2,\bbC)&\bbC^*&\mathsf{SO}(2)&\mathsf{U}(1)&\mathsf{U}(1)&\mathsf{U}(1)\\
\mathsf{Spin}(3,\bbC)&\mathsf{SL}(2,\bbC)&\mathsf{SU}(2)&\mathsf{Sp}(2)&\mathsf{SU}(2)&\bbH_1\simeq\mathsf{SU}(2)\\
\mathsf{Spin}(4,\bbC)&\mathsf{SL}(2,\bbC)\otimes\mathsf{SL}(2,\bbC)&\mathsf{SU}(2)\otimes\mathsf{SU}(2)&^2\mathsf{Sp}(2)&&\bbH_1\times\bbH_1\\
\mathsf{Spin}(5,\bbC)&\mathsf{Sp}(2,\bbC)&\mathsf{Sp}(2)&\mathsf{Sp}(4)&&\mathsf{Sp}(2)\\
\mathsf{Spin}(6,\bbC)&\mathsf{SL}(4,\bbC)&\mathsf{SU}(4)&\mathsf{SU}(4)&&\mathsf{SU}(4)\\
\end{array}$$
This is for \textit{Cl}$_{0,n}$ and \textit{Cl}$_{n,0}$ only


\textit{Isomorphisms}. According to Garling book. Tai sutampa su Trautman\\
$\mathsf{Spin}^{+}(1,2)\cong\mathsf{SL}(2,\bbR)$\\
$\mathsf{Spin}^{+}(2,2)\cong\mathsf{SL}(2,\bbR)\times
\mathsf{SL}(2,\bbR)$\\
For \textit{Cl}$_{0,n}$ algebras, Spin(n) groups in
Baczkowski~\cite{Baczkowski09}

%%%%%%%%%%%%%%%%%%%%%%%%%%%%%%%%%%%%%%%%%%%%%%%%


\vspace{5mm}\noindent$\blacktriangleright$ A different
representation is also possible~\cite{Gull91}
\[\label{DiracSpinor1}
|\psi\rangle = \begin{bmatrix}
 a^0 + \ii b^3 \\
 -b^2 + \ii b^1 \\
 a^3 + \ii b^0 \\
 a^1 + \ii a^2
  \end{bmatrix}\quad\Longleftrightarrow\quad
 \psi = a_0+\boldsymbol{a}+I(b_0+\boldsymbol{b}) ,
\]
where $\boldsymbol{a}=a^k\sig{k}$ and $\boldsymbol{b}=b^k\sig{k}$
are relative vectors $a^k$ and $b^k$ are the scalars, and
$\sig{1}=$

%%%%%%%%%%%%%%%%%%%%%%%%%%%%%%%%%%%%%%%%%%%%%%%%


\vspace{5mm}\noindent\CL{41}  In \textit{Cl}$_{4,1}$ (i{\v s} mano
Phys. Scr.~\cite{Dargys10bCL}) the most general structure of the
spinor-multivector of $\mathbf{Spin}(4,1)$  group is
\begin{equation}
\psi=c_s+b_{12}\e{12}+b_{13}\e{13}+b_{14}\e{14}+b_{23}\e{23}+b_{24}\e{24}+b_{34}\e{34}+T_{2134}I\e{5},
\end{equation}
where all coefficients are real and satisfy
\[\begin{split}
&c_s^2+b_{12}^2+b_{13}^2+b_{14}^2+b_{23}^2+b_{24}^2+b_{34}^2+T_{2134}^2=1,\\
&c_s T_{2134}+b_{12}b_{34}-b_{13}b_{24}+b_{14}b_{23}=0.\\
\end{split}\]


***************************************************

\begin{example} $\blacktriangleright$ \ex~8. The complexificaion may
be used to describe  $2n$ first order equations -- Hamilton's
equations for coordinate $q_i=\partial H/\partial p_i$ and
momentum $p_i=-\partial H/\partial q_i$.

How calculations are done  with complex quaternions is described
by Kravchenko in book~\cite{Kravchenko03}

 \vspace{5mm}\noindent \textit{Biquaternions} $\mathbb{B}$
are just complexified quaternions:
$\mathbb{B}=\bbH\otimes\bbC\simeq\cl{3}{0}$.


It is always possible to represent the even-dimensional complex GA
in higher-dimensional real GA's.

%%%%%%%%%%%%%%%%%%%%%%%%%%%%%%%%%%%%%
\begin{unit}
\textit{Group isomorphisms} Gilbert~\cite{Gilbert91}, Helgason
book p.520.\blue{Ar verta duoti?}\red{siulau ismest: visu pirma
O(5,\bbC) neapibrezta, antra $\mathsf{Sp}(4,\bbR)\simeq
\mathsf{SO}(5,\bbR)$ nedera su pries tai duotu isomorfizmu
$\mathsf{Sp}(4,\bbR)\cong\mathsf{SO}(3,2,\bbR)$, trecia
$\mathsf{Sp}(2,2,\bbR)$ vel neapibrezta, nes mes kalbedami apie
$\mathsf{Sp}$ neivedem kitu bitiesiniu formu, kuriose butu svarbi
signatura.}
\[\begin{split}
\mathsf{Sp}(2,\bbH)\cong\mathsf{O}(5,\bbC)&\qquad\qquad
\mathsf{Sp}(4,\bbR)\simeq \mathsf{SO}(5,\bbR)\\
\mathsf{Sp}(2,2,\bbR)\cong\mathsf{SO}(4,1,\bbR)&\qquad\qquad\\
\end{split}\]
\end{unit}

\begin{example}
*******************************************

We are interested in establishing relations between $\mathsf{SO}(3)$
and $\mathsf{SU}(2)$. Use of the global parametrization can help
or complicate explicit calculation of generators. Let us take
fundamental representation of $\mathsf{SU}(2)$, which is given by
$2\times 2$ so called Wigner matrices. We will use parametrization
of fixed axes. In this parametrization the representation of the
matrix is obtained by first rotating system around $z$ axis by
angle $\gamma$ (with $0\le\gamma< 2\pi$), then rotating by angle
$\beta$ about old (globally fixed) axis $y$ (with
$0\le\beta\le\pi$) and, at last, rotating by angle $\alpha$ (with
$0\le\alpha< 2\pi$) once more about old (fixed) axis $z$.
Mathematically the above sequence of rotations can be written as
\renewcommand{\kbldelim}{[}% Left delimiter
\renewcommand{\kbrdelim}{]}% Right delimiter
\[D^{1/2}_{m,m^\prime}(\alpha,\beta,\gamma)=\ee^{-\ii m\, \alpha} d^{1/2}_{m,m^\prime} \ee^{-\ii m^\prime\, \gamma} =
\begin{bmatrix}
  \ee^{-\frac{\ii (\alpha+\gamma)}{2}} \cos \frac{\beta}{2} & -\ee^{\frac{\ii (-\alpha+\gamma)}{2}}\sin \frac{\beta}{2}\\
  \ee^{\frac{\ii (\alpha-\gamma)}{2}}\sin \frac{\beta}{2} & \ee^{\frac{\ii (\alpha+\gamma)}{2}} \cos \frac{\beta}{2}
\end{bmatrix},
\]
where $d^{1/2}_{m,m^\prime}=
\kbordermatrix{&m^\prime=\frac12& m^\prime=-\frac12\\
m=\frac12&\cos(\beta/2)&-\sin(\beta/2)\\
m=-\frac12&\sin(\beta/2)&\cos(\beta/2)}$. \blue{Nepasakyta apie
$m,m^{\prime}$, kurie stebuklingai dingsta. Be to, kazkas negerai,
$m^{\prime}=\tfrac{1}{2}m^{\prime}$} For those who fail to see
other two rotations we shall remind that complex exponent
represents general rotations in complex plane (multiplication by
imaginary unit means rotation by $90^{\circ}$ in complex plane,
i.~e. real number becomes complex \blue{imaginary} and complex
\blue{imaginary} turns into negative real). The inverse of
$D^{1/2}_{m,m^\prime}(\alpha,\beta,\gamma)$ is easy to obtain by
sign change of angles: $\alpha \rightarrow -\gamma, \beta
\rightarrow -\beta, \gamma \rightarrow -\alpha $, i.~e.,
\[
  \left(D^{1/2}_{m,m^\prime}(\alpha,\beta,\gamma)\right)^{-1}=D^{1/2}_{m,m^\prime}(-\gamma,-\beta,-\alpha)= \begin{bmatrix}
  \ee^{\frac{\ii (\alpha+\gamma)}{2}} \cos \frac{\beta}{2} & \ee^{\frac{\ii (-\alpha+\gamma)}{2}}\sin \frac{\beta}{2}\\
  -\ee^{\frac{\ii (\alpha-\gamma)}{2}}\sin \frac{\beta}{2} & \ee^{\frac{-\ii (\alpha+\gamma)}{2}} \cos \frac{\beta}{2}
\end{bmatrix}
\]
It is easy to see that this replacement is equivalent to
transposition and complex conjugation, because the matrices that
represent  $\mathsf{SU}(2)$ are unitary, i.~e. $U^{-1}=U^\dagger$.

For \cl{3}{0} algebra the basis  vectors are represented by
standard Pauli matrices $\sigma_1=
\begin{bmatrix}0&1\\1&0\end{bmatrix}$, $\sigma_2=
\begin{bmatrix}0&-\ii\\ \ii&0\end{bmatrix}$ and $\sigma_3=
\begin{bmatrix}1&0\\ 0&-1\end{bmatrix}$. Then bivectors are
$\sigma_{23}= \begin{bmatrix}0&\ii\\\ii&0\end{bmatrix}$,
$\sigma_{13}= \begin{bmatrix}0&-1\\1&0\end{bmatrix}$ and
$\sigma_{12}= \begin{bmatrix}-\ii&0\\ 0&\ii\end{bmatrix}$, and we
hope to discover the match $\sigma_{23}=J_x$, $-\sigma_{13}=J_y$
and  $\sigma_{12}=J_z$ between bivectors of  \cl{3}{0} and
generators of $\mathsf{SU}(2)$ algebra.

Lets calculate $\mathsf{SU}(2)$ generators. We intentionally do
this for arbitrary point of manifold, i.~e. first calculate
general derivatives and then substitute parameter values were
these derivatives to be evaluated. Rotation by angle $\alpha$ is
around $z$ axes, therefore by differentiating with respect to
$\alpha$ we expect to get $J_z$.  This is indeed the case for any
parameter values
\[
J_z=\frac{\partial
D^{1/2}_{m,m^\prime}(\alpha,\beta,\gamma)}{\partial \alpha}\d
\left(D^{1/2}_{m,m^\prime}(\alpha,\beta,\gamma)\right)^{-1}=
\frac12\begin{bmatrix}-\ii&0\\0&\ii\end{bmatrix}.
\]
Doing the same for parameter $\beta$ we hope to get generator
$J_y$. We get
\[
\frac{\partial D^{1/2}_{m,m^\prime}(\alpha,\beta,\gamma)}{\partial
\beta} \d
\left(D^{1/2}_{m,m^\prime}(\alpha,\beta,\gamma)\right)^{-1}=
\frac12
\begin{bmatrix}0&-\ee^{-\ii \alpha}\\\ee^{\ii \alpha}&0\end{bmatrix}.
\]
We see that now the result would depend on parameters  where the
generator is calculated. We can imagine the calculated algebra
element as tangent plane to group manifold. Note, however, that
the plane always is centered at zero matrix, i.~e. planes parallel
to the manifold are moved to neighborhood of zero matrix (not to
identity element of manifold!). Taking  $\alpha=0$ we see that
indeed the $\beta$ is the rotation about $y$ axis in our
parametrization.
\[
J_y=\frac12\left.\begin{bmatrix}0&-\ee^{-\ii \alpha}\\\ee^{\ii
\alpha}&0\end{bmatrix}\right|_{\alpha=0}=\frac12\begin{bmatrix}0&-1\\
1&0\end{bmatrix}.
\]
Differentiation with respect to parameter $\gamma$ yields
\[
\frac{\partial D^{1/2}_{m,m^\prime}(\alpha,\beta,\gamma)}{\partial
\gamma}\d
\left(D^{1/2}_{m,m^\prime}(\alpha,\beta,\gamma)\right)^{-1}=
 \frac12\begin{bmatrix}-\ii\cos\beta&-\ee^{-\ii \alpha}\sin\beta\\-\ii \ee^{\ii \alpha}\sin\beta&\ii\cos\beta\end{bmatrix}.
\]
Now if we evaluate it at $\alpha=0, \beta=0$ we obtain $
\frac12\left.\begin{bmatrix}-\ii\cos\beta&-\ee^{-\ii
\alpha}\sin\beta\\-\ii \ee^{\ii
\alpha}\sin\beta&\ii\cos\beta\end{bmatrix}\right|_{\alpha=0,\beta=0}=\frac12\begin{bmatrix}-\ii&0\\0&\ii\end{bmatrix}=J_z$,
i.~e. the same generator as for $\alpha$. This is not surprising,
because in our  parametrization both parameters corresponds to
rotation about the same fixed $z$ axis.

In order to get the remaining generator $J_x$ we have to evaluate
the  expression at different parameter values. For example,
evaluation at $\alpha=0, \beta=-\pi/2$ yields
\[\frac12\left.\begin{bmatrix}-\ii\cos\beta&-\ee^{-\ii \alpha}\sin\beta\\-\ii \ee^{\ii \alpha}\sin\beta&\ii\cos\beta\end{bmatrix}\right|_{\alpha=0,\beta=-\pi/2}=\frac12\begin{bmatrix}0&\ii\\\ii&0\end{bmatrix}=J_x
.\] If we would take different $\mathsf{SU}(2)$ parametrization,
for example,  well known matrices of the form
$\begin{bmatrix}\ee^{\ii \alpha}\cos\beta&\ee^{-\ii \gamma}\sin\beta\\
-\ee^{\ii \gamma}\sin\beta&\ee^{-\ii
\alpha}\cos\beta\end{bmatrix}$ then in order to get generators in
the form of Pauli matrices we would need to take linear
combination of $\frac{\partial A}{\partial \alpha}\d A^{-1}$ and
$\frac{\partial A}{\partial \gamma}\d A^{-1}$  evaluated at
specific points of parameter space.
 \end{example}


%perkelta 20170510
\begin{unit} \textit{Clifford group}~$\Gamma_{p,q}$. It is associated with every
  Clifford algebra \cl{p}{q}. The Clifford  group
$\Gamma_{p,q}$ (in fact invented by R.~Lipschits)  defines how the
vector $x$ is mapped onto other vector $\m{R}x\m{R}^{-1}$ in the
same space,
\[\Gamma_{p,q}=\{\m{R}\in\cl{p}{q}\,|\,\forall x\in\bbR^{p,q},
\m{R}x\m{R}^{-1}\in\bbR^{p,q}\}.\] Gal ra'syti $\quad
\bbR^{p,q}\rightarrow\mathcal{V}^n$ Thus, the Clifford group
element is determined by its two sided mapping, $x\mapsto
\m{R}x\m{R}^{-1}$, of vectors, where $\m{R}$ has inverse
$\m{R}^{-1}$ but may not necessarily be the rotation.
\end{unit}

\begin{unit} \textit{Pin group}~$\Pin{p}{q}$. Normalization of
the Clifford group, $\m{R}\tilde{\m{R}}=\pm 1$, yields the pin
(also called pinor) group
\[\Pin{p}{q}=\{\m{R}\in\Gamma_{p,q}\,|\,\m{R}\tilde{\m{R}}=\pm 1\}.\]
This group  is a two-fold covering of the orthogonal group
$\mathsf{O}(p,q)$.
\end{unit}
%
\begin{unit}
\textit{Spin group} $\Spin{p}{q}$. This group is important
in physical applications. The Clifford algebra itself allows to
construct the half-spin representations of the spin (or spinor)
groups. The spin group is a subgroup of $\Pin{p}{q}$ and
consists of even-grade elements:
\[\Spin{p}{q}=\Pin{p}{q}\cap\cl{p}{q}^{+} \]
For example, in Euclidean space the elements of the spin group
correspond to rotations $\m{R}$ that map
$\bbR^n\rightarrow\bbR^n$, and transform the vectors
\[x\rightarrow x^{\prime}=\m{R}x\tilde{\m{R}},\quad \m{R}\in\mathsf{Spin}(n,0). \]
An important in quantum mechanics subgroup of
$\Spin{p}{q}$, which has  normalization $+1$, is
\[\Spin{p}{q}^{+}=\{\m{R}\in\Spin{p}{q}\,|\,\m{R}\tilde{\m{R}}=+1\}.\]
The spinor group is symmetric in signature indices,
$\mathsf{Spin}_{pq}\cong \mathsf{Spin}_{qp}$
\end{unit}

% negalejau patikrinti i suprasti is Varlamov:
\item[$\bullet$] For $p-q=1,3,5,7 \mod(8)$ the following isomorphisms hold
 \[
\begin{split}
  \Pin{p}{q} &\cong\Pin{p}{q-1} \cup I \Pin{p}{q-1}\\
    &\cong\Pin{q}{p-1} \cup I \Pin{q}{p-1}\\
    &\cong\Spin{q}{p} \cup I \Spin{q}{p}
\end{split}
  \]
Here $I=\e{1}\e{2}\cdots\e{n}$ is the algebra pseudoscalar. When $p-q=1,5 \mod(8)$ the $$




\begin{unit}\textit{Properties of  spin groups}.
\begin{itemize}
%
\item[$\bullet$] When  $n\le 6$  the spinor  groups
$\Spin{p}{q}$ in table~\ref{tableSG} are isomorphic to
listed here classical groups.  \blue{Iki
galo nesuprantu. Pagal Porteous \cite{Porteous95} psl. 178 tam
reikalinga Cayley algebra su oktanionais. Kadangi jie
neasociatyvus, sektu, kad spinorius negalime pavaizduoti
matricomis kai $n>6$. Nuo 'cia kilo pavadinimas exceptional
groups? I's kitos, ivedus triality savoka ('zr. Porteous psl. 256)
galima dirbti su matricomis? Exceptional phenomena, kaip ra'so
Porteous.}
%
\item[$\bullet$]


%
\item[$\bullet$]  Structural constants of Lie algebra can be found
from commutators of bivectors. They are invariants of the algebra.
I'splesti
%
\item[$\bullet$] The isomorphism between classical and spin groups
allows to apply the GA in the construction of classical Lie groups
and their representation by matrices.


%
\item[$\bullet$] The irreducible  matrices that represent the
spinor group $\Spin{p}{q}$ may be constructed from
products of $8$-fold periodicity table matrices  that represent
basis vectors of \cl{p}{q} algebras. From the table in
item~\ref{8foldperiod} follows that spinor groups may consist of
real $\bbR$, complex $\bbC$ and quaternionic
$\bbH$ matrices including their direct sums.


\begin{unit}\textit{More properties on pin and spin groups}
\begin{itemize}
\item[$\bullet$] The following isomorphism between Clifford groups
and their classical counterparts exists
\[\begin{split}
&\Pin{p}{q}\backslash\{\pm 1\}\simeq\mathsf{O}(p,q)\\
&\Spin{p}{q}\backslash\{\pm 1\}\simeq\mathsf{SO}(p,q)\\
&\mathsf{Spin}^{+}_{p,q}\backslash\{\pm 1\}\simeq\mathsf{SO}^+(p,q)\\
\end{split}\]
where $\mathsf{O}(p,q)$ is the orthogonal group,
$\mathsf{SO}(p,q)$ is special orthogonal group, and
$\mathsf{SO}^+(p,q)$ is special orthogonal group with the
determinant $+1$. Thus, the Clifford groups are two-sheeted
coverings of their classical orthogonal counterparts.


\item[$\bullet$] For a mixed signature $p,q\ge 1$ and any pair of
orthogonal vectors $\e{+}$ and $\e{-}$ such that $\e{+}^2=1$,
$\e{-}^2=-1$ one has that~\cite{Lundholm09} p.65
\[\begin{split}
&\Pin{p}{q}=\mathsf{Spin}^{+}_{p,q}\otimes\{1,\e{+},\e{-},\e{+}\e{-}\}\\
&\Spin{p}{q}=\mathsf{Spin}^{+}_{p,q}\times\{1,\e{+}\e{-}\}.\\
\end{split}\]
\item[$\bullet$] For Euclidean $(p=n,q=0)$ and anti-Euclidean
$(p=0,q=n)$ signatures one has that
\[\Pin{p}{q}=\mathsf{Spin}^{+}_{p,q}\times\{1,\e{}\}, \]
where any $\e{}\in\mathcal{V}$ is such that $\e{}^2=\pm 1$. From
these properties follows  that it is sufficient to study only the
properties of the rotor group $\mathsf{Spin}^{+}_{p,q}$ to
understand the $\mathsf{Spin}$ and $\mathsf{Pin}$ groups.

In the matrix representation the group operation $\circ$ is matrix
multiplication. The matrix multiplication is associative.  The
invertible matrix (by definition) has inverse.
\end{itemize}
\end{unit}

%%%%%%%%%%%%%%%%%%%%%%%%
\textit{Relation of GA with $\mathsf{Spin}(p,q)$ group}. The GA
spin group is constructed from basis vectors that obey
$\e{i}\e{j}+\e{j}\e{i}=2\delta_{ij}$. The number of elements of
the $\mathsf{Spin}(n)$ group (order of the group) is
$\tfrac{1}{2}n(n-1)$, $n=p+q$.
\end{unit}

%
\begin{example} \blue{ TAISYTA``Every Lie algebra can be
represented as a bivector algebra, hence every Lie group can be
represented as a spin group''~\cite{Doran93}. Therefore, the
number of Lie algebra elements (generators), Lie group properties
and their respective matrix  representations follow from  the
$8$-periodicity and Table in~\ref{tableSG}.}

The table in~\ref{tableSG} shows that spin group of \cl{3}{0}
algebra is $\Spin{3}{0}\cong\mathsf{SU}(2,\bbC)$, which is  the
Lie group $\mathsf{SU}(2,\bbC)$. The counterpart of the Lie group
is the Lie algebra  denoted as $\mathfrak{su}(2)$. The elements of
Lie algebra satisfy the axioms in~\ref{Liealgebra}. In GA, the Lie
algebra can be generated by commutators of GA bivectors.

As item~\ref{classgroup} shows the number of elements in the group
$\mathsf{SU}(2,\bbC)$ is $n_{\mathrm{G}}=n^2-1=3$, which is equal
to the number of bivectors of \cl{3}{0}, $\tfrac{1}{2}n(n-1)=3$,
namely, $\e{12}$, $\e{23}$, $\e{31}$. The commutators of bivectors
give a closed multiplication (commutation) table,
\[\begin{array}{c|ccc}
\tfrac{1}{2}[\,.\,,\,.\,] &\e{12}&\e{23}&\e{31}\\
\hline
 \e{12}&0 &-\e{31} &\e{23} \\
 \e{23}&\e{31} &0 &-\e{12} \\
 \e{31}&-\e{23}&\e{12} &0 \\
\end{array}\]
The standard matrices that represent $\mathfrak{su}(2)$ algebra
may be expressed via Pauli matrices as $-\ii\hat{\sigma}_1$,
$-\ii\hat{\sigma}_2$, and $\ii\hat{\sigma}_3$, which give the
following commutation table,
\[\begin{array}{c|ccc}
\tfrac{1}{2}[\,.\,,\,.\,] &-\ii\hat{\sigma}_1&-\ii\hat{\sigma}_2&\ii\hat{\sigma}_3\\
\hline
 -\ii\hat{\sigma}_1&0 &-\ii\hat{\sigma}_3 &-\ii\hat{\sigma}_2 \\
 -\ii\hat{\sigma}_2&\ii\hat{\sigma}_3 &0 &\ii\hat{\sigma}_1 \\
 \ii\hat{\sigma}_3&\ii\hat{\sigma}_2&-\ii\hat{\sigma}_1 &0 \\
\end{array}\]
We see that both commutation table coincide. From tables we find
that the structural constants of $\mathfrak{su}(2)$ are $\pm 1$
and $0$.  Thus, both algebras are isomorphic. The following
relations between  GA generators and standard matrix generators of
$\mathfrak{su}(2)$ exist,
$$
\e{12}\leftrightarrow\left[\begin{array}{cc}0&
-\ii\\-\ii&0\end{array}\right]=-\ii\hat{\sigma}_1,\quad
\e{23}\leftrightarrow\left[\begin{array}{cc}0&-1\\1&0\end{array}\right]=-\ii\hat{\sigma}_2\quad
\e{31}\leftrightarrow\left[\begin{array}{cc}\ii&0\\0&-\ii\end{array}\right]=\ii\hat{\sigma}_3
$$
To obtain the Lie group in GA form or matrix form it is enough to
take an exponential map of obtained matrices multiplied by scalar
parameter (see the item~\ref{Liealgebra}),
\[\begin{split}
&\ee^{\alpha\e{12}}=\cos\alpha+\e{12}\sin\alpha\leftrightarrow\ee^{-\alpha\ii\hat{\sigma}_1}=
\left[\begin{array}{cc}\cos\alpha&-\ii\sin\alpha\\-\ii\sin\alpha&\cos\alpha\end{array}\right],\\
&\ee^{\beta\e{23}}=\cos\beta+\e{23}\sin\beta\leftrightarrow\ee^{-\beta\ii\hat{\sigma}_2}=
\left[\begin{array}{cc}\cos\beta&-\sin\beta\\\sin\beta&\cos\beta\end{array}\right],\\
&\ee^{\gamma\e{31}}=\cos\gamma+\e{31}\sin\gamma\leftrightarrow\ee^{\gamma\ii\hat{\sigma}_3}=
\left[\begin{array}{cc}\ee^{\ii\gamma}&0\\0&\ee^{-\ii\gamma}\end{array}\right].
\end{split}\]
The determinant of matrices is equal to $1$ The matrices satisfy
$\m{SU}(n)$ group requirement
$\hat{U}\hat{U}^{-1}=\hat{U}\hat{U}^{\dagger}=\hat{1}$.  They
describe rotations around three orthogonal axes. In conclusion,
the chain which starts from GA,
$\mathcal{V}^n\rightarrow\mathcal{G}_2^n\rightarrow\mathfrak{su}(n)\rightarrow\mathsf{SU}(n)$,
allows to generate groups. Thus, GA provides a powerful tool to
construct and study the Lie algebra and groups via GA.
\end{example}



%
\begin{unit}\blue{NAUJAS \textit{Lie algebra generators (bivectors) in GA and equivalent irreducible
matrix representations}}. $\hat{\sigma}_i$ are Pauli matrices

\[\bullet\ \m{Spin}(3)\cong\mathfrak{su}(2,\mathbb{C}),\quad\e{12}
\leftrightarrow \ii\hat{\sigma}_3,\quad\e{23}\leftrightarrow
-\ii\hat{\sigma}_1,\quad \e{31}\leftrightarrow
-\ii\hat{\sigma}_2\]
%
\[\bullet\ \m{Spin}(2,1)\cong\mathfrak{sl}(2,\mathbb{R}),\quad\e{12}\leftrightarrow
-\ii\hat{\sigma}_1,
\quad\e{23}\leftrightarrow\ii\hat{\sigma}_2,\quad\e{31}\leftrightarrow
\ii\hat{\sigma}_3\]
%
\[\begin{split}
 \bullet\ \m{Spin}(1,3)\cong\mathfrak{sl}(2,\mathbb{C}),\quad&\e{12}\leftrightarrow\hat{\sigma}_1,
\quad\e{13}\leftrightarrow\hat{\sigma}_2,\quad\e{14}\leftrightarrow
-\hat{\sigma}_3,\quad\\
&\e{23}\leftrightarrow
-\ii\hat{\sigma}_3,\quad\e{42}\leftrightarrow
\ii\hat{\sigma}_2,\quad\e{34}\leftrightarrow
\ii\hat{\sigma}_1\end{split}\]
%
\[\begin{split}
 \bullet\ \m{Spin}(6,0)\cong\mathfrak{su}(4,\mathbb{C}),
\end{split}\]

Lie group elements in GA and matrix form can be found after
exponentiation of bivectors and matrices multiplied by angles
(parameters) $\alpha$, $\beta$ etc, respectively.
\end{unit}
%

\begin{unit}\textit{The real compact spin groups},\label{spingroups}
Trautman~\cite{Trautman06}. Notation:
$\mathsf{Spin}(n)\equiv\Spin{n}{0}=\Spin{0}{n}$.
\[\begin{split}
&\mathsf{Spin}(1)=\Spin{1}{0}\cong\Spin{0}{1}\cong\mathsf{O}(1,\bbR)\\
&\mathsf{Spin}(2)=\Spin{2}{0}\cong\Spin{0}{2}\cong\mathsf{U}(1,\bbC)\cong\mathsf{SO}(2,\bbR)\\
&\mathsf{Spin}(3)=\Spin{3}{0}\cong\Spin{0}{3}\cong\mathsf{SU}(2,\bbC)\cong\mathsf{SL}(1,\bbH)\\
&\mathsf{Spin}(4)=\Spin{4}{0}\cong\Spin{0}{4}\cong{^2\mathsf{SU}(2,\bbC)}\\
&\mathsf{Spin}(5)=\Spin{5}{0}\cong\Spin{0}{5}\cong\mathsf{Sp}(2,\bbH)\\
&\mathsf{Spin}(6)=\Spin{6}{0}\cong\Spin{0}{6}\cong\mathsf{SU}(4,\bbC)\\
\end{split}\]
\end{unit}
 \vspace{2mm}


\begin{unit}\label{Trautman3}
 The groups $\Spin{p}{q}^+$,
Trautman~\cite{Trautman06,Todorov11}
\[\begin{split}
&\mathsf{Spin}^{+}(1,1)=\Spin{1}{1}\cong\mathsf{GL}(1,\bbR)\cong\bbR
\end{split}\]
\vspace{0mm}
\[\begin{split}
&\mathsf{Spin}^{+}(2,1)=\Spin{2}{1}=\Spin{1}{2}\cong\mathsf{SU}(1,1,\bbC)\cong\mathsf{SL}(2,\bbR)\\
\end{split}\]
\vspace{2mm}
\[\begin{split}
&\mathsf{Spin}^{+}(3,1)=\Spin{3}{1}=\Spin{1}{3}\cong\mathsf{SL}(2,\bbC)\\
&\mathsf{Spin}^{+}(2,2)=\Spin{2}{2}\cong{^2\mathsf{SL}(2,\bbR)}\\%?(Sp(2,R)Garling)
\end{split}\]
\vspace{2mm}
\[\begin{split}
&\mathsf{Spin}^{+}(4,1)=\Spin{4}{1}=\Spin{1}{4}\cong\mathsf{Sp}(1,1,\bbH)\\
&\mathsf{Spin}^{+}(3,2)=\Spin{3}{2}=\Spin{2}{3}\cong\mathsf{Sp}(4,\bbR)\\
\end{split}\]
\vspace{2mm}
\[\begin{split}
&\mathsf{Spin}^{+}(5,1)=\Spin{5}{1}=\Spin{1}{5}\cong\mathsf{SL}(2,\bbH)\\
&\mathsf{Spin}^{+}(4,2)=\Spin{4}{2}=\Spin{2}{4}\cong\mathsf{SU}(2,2,\bbC)\\
&\mathsf{Spin}^{+}(3,3)=\Spin{3}{3}\cong\mathsf{SL}(4,\bbR)\\
\end{split}\]
where the notation
$\mathsf{S0}(n)=\mathsf{S0}(n,0)=\mathsf{S0}(0,n)$,
$\mathsf{SU}(n)=\mathsf{SU}(n,0)=\mathsf{SU}(0,n)$ etc is used.
\end{unit}
%

%%%%%%%%%% 2017-05-19 ACUS %%%%%%%%%%%%%%%%%%%%%%%%%%%%%%%%%%%%%%%%
% zemiau surasyta medziaga butinai reiks panaudoti veliau
\begin{unit}\textit{Subalgebras of GA bivector algebra}
  \red{new unit, Sita skyriu kol kas, manau pasalinkime (galite istrinti, as jau nukopijavau ji i failo gala). Jis labai idomus, bet kol kas manau neisvystytas. Pavyzdziui, jis rodo, kad GA be bivektoriu algebru (susijusiu su Lie) egzistuoja ir kitos Lie algebros. Bet tai reikalauja dar daug tikrinimo ir apmastymo. Todel kaip buvo gautos matricu representatcijos geriau kol kas neaiskinti.}
Using Clifford algebra isomorphisms (see \ref{evenoddalgebras}) $\cl{p}{q+1}^+\cong\cl{p}{q}$ and $\cl{p+1}{q}^+\cong\cl{q}{p}$ and noting that
\begin{itemize}
\item number of $\cl{p+1}{q}$  bivectors $(p+1+q)(p+q)/2$ is equal
to a total number of $\cl{q}{p}$ algebra vectors $p+q$ and
bivectors $(p+q)(p+q-1)/2$ \item $\cl{q}{p}$ set of
vectors+bivectors is closed with respect  to commutator operation,
i.e. commutator of two vectors gives bivector (or zero), and
commutator of vector and bivector is a vector (or zero),
\end{itemize}
we have the following\newline \textbf{Theorem}. Bivector
commutator  algebra of $\cl{p+1}{q}$ is isomorphic to $\cl{q}{p}$
commutator algebra made of vectors and bivectors. Bivector algebra
of $\cl{p}{q+1}$ is isomorphic to $\cl{p}{q}$ algebra made of
vectors and bivectors. Bivectors of $\cl{q}{p}$ and $\cl{p}{q}$
then make (maximal) commutator subalgebras of $\cl{p+1}{q}$ and
$\cl{p}{q+1}$, respectively.

Vector+bivector algebras are not the only possible subalgebras of
bivector commutator algebra.  For $p+q\ge5$ there exists
subalgebras which include other grade elements as well. For
example, $\cl{1}{4}$ bivectors have subalgeba, which include
elements of all grades of $\cl{1}{3}$ algebra. It also means that
we can make GA algebra from specific sets of generators which not
all are bivectors. For  $\cl{0}{n}$ and $\cl{n}{0}$ algebras,
however, all subalgebras emerge form vector+bivector split.
\end{unit}
\begin{example}
Represent $\cl{1}{3}$ bivector commutator algebra by $\cl{1}{2}$
and $\cl{3}{0}$ vector+bivector commutator algebras.

Bivectors of $\cl{1}{3}$ are $\{\e{12},\e{13},\e{14},\e{23},\e{24},\e{34}\}$. If we put $\cl{1}{2}$ vectors and bivectors in the following order $\{\e{1},\e{13},\e{12},\e{3},\e{2},\e{23}\}$, then it is easy to check that substitution $\{\e{12},\e{13},\e{14},\e{23},\e{24},\e{34}\}\Rightarrow \{-\e{1},-\e{13},-\e{12},-\e{3},-\e{2},-\e{23}\}$
makes both commutator tables identical.

In the case of $\cl{3}{0}$ the properly ordered set of vectors+bivectors is $\{\e{1},\e{2},\e{3},\e{12},\e{13},\e{23}\}$. The substitution
$\{\e{12},\e{13},\e{14},\e{23},\e{24},\e{34}\}\Rightarrow \{-\e{1},-\e{2},-\e{3},-\e{12},-\e{13},-\e{23}\}$ then gives well known result of time-space split of 4D relativistic algebra into 3D~vectors (boosts in 4D) and bivectors (space rotations). This explains why 4D~space rotations (3D bivectors) make subalgebra, whereas boosts (3D vectors) don't. Analogous splits can be easily constructed for any other Clifford algebra.
\end{example}

%%%%%%%%%%%%%%%%%%%%%%%%%%%%%%%%%%%%%%%%%%%%%%%%%%%%%%%%%%%%%%%

For any bivector $\cB\in\cG^2$ we have that
$\pm\ee^{\cB}\in\m{Spin}^{+}$, in other words
$\pm\ee^{\cG^2}\subseteq\m{Spin}^{+}$.

 ************************************************


The properties of  action of the adjoint operation can be
understood from multiplication table for the triangle group, where
the entries  are $xyx^{-1}$ and $x,y\in\m{G}$,
\[\begin{array}{cc|ccccccc}
xyx^{-1}&x=&R&R^2&R^3&\hat{b}&\hat{b}R&\hat{b}R^2\\
y=&&&&&&&\\   \hline
R&&R&R&R&-R^2&-R^2&-R^2\\
R^2&&R^2&R^2&R^2&-R&-R&-R\\
R^3&&R^3&R^3&R^3&R^3&R^3&R^3\\
\hat{b}&&-\hat{b}R&\hat{b}R^2&\hat{b}&\hat{b}&\hat{b}R^2&-\hat{b}R\\
\hat{b}R&&-\hat{b}R^2&-\hat{b}&\hat{b}R&-\hat{b}R^2&\hat{b}R&-\hat{b}\\
\hat{b}R^2&&\hat{b}&-\hat{b}R&\hat{b}R^2&-\hat{b}R&\hat{b}&\hat{b}R^2\\
\end{array}\]
The table shows that under adjoint map $\textrm{Ad}_x:y\mapsto
xyx^{-1}$ the group is divided into two classes that represent
either rotations or reflections. The division of groups into
classes by adjoint transformation is frequently used in solid
state theory.
